\documentclass[11pt]{article}

\usepackage[margin=1in]{geometry}
\usepackage{amsmath,amssymb,amsthm,mathtools}
\usepackage{booktabs,tabularx}
\usepackage[authoryear]{natbib}
\usepackage[colorlinks=true,allcolors=blue]{hyperref}
\usepackage{microtype}

\newtheorem{theorem}{Theorem}[section]
\newtheorem{proposition}[theorem]{Proposition}
\newtheorem{corollary}[theorem]{Corollary}
\newtheorem{lemma}[theorem]{Lemma}
\theoremstyle{definition}
\newtheorem{definition}[theorem]{Definition}
\newtheorem{assumptions}[theorem]{Assumptions}
\theoremstyle{remark}
\newtheorem{remark}[theorem]{Remark}
\theoremstyle{plain}
\newtheorem*{principle}{Principle}
\theoremstyle{definition}
\newtheorem{example}[theorem]{Example}

\newcommand{\Acal}{\mathcal{A}}
\newcommand{\Ccal}{\mathcal{C}}
\newcommand{\Ocal}{\mathcal{O}}
\newcommand{\Pcal}{\mathcal{P}}
\newcommand{\Hcal}{\mathcal{H}}
\newcommand{\Fcal}{\mathcal{F}}
\newcommand{\Gcal}{\mathcal{G}}
\newcommand{\Qcal}{\mathcal{Q}}
\newcommand{\Rcal}{\mathcal{R}}
\newcommand{\Scal}{\mathcal{S}}
\newcommand{\E}{\mathbb{E}}
\newcommand{\Prob}{\mathbb{P}}
\newcommand{\KL}{\mathrm{KL}}
\newcommand{\I}{\mathrm{I}}
\newcommand{\Hh}{\mathrm{H}}
\newcommand{\argmin}{\mathop{\mathrm{arg\,min}}}
\newcommand{\argmax}{\mathop{\mathrm{arg\,max}}}
\newcommand{\givena}{\,\|\,}
\newcommand{\ellp}{\ell}
\newcommand{\wt}{\widetilde}
\newcommand{\barxi}{\bar{\xi}}
\newcommand{\one}{\mathbf{1}}
\newcommand{\problembox}[1]{%
\begin{center}
\fbox{\parbox{0.93\textwidth}{#1}}
\end{center}}

\title{Sufficient Conditions for Semantic Convergence}
\author{}
\date{}

\begin{document}

\maketitle

\begin{abstract}
Three layers govern any interactive learning question. An \emph{observer} is a passive lens: a measurable map from programs to views that fixes what is in principle distinguishable through a given access pattern. An \emph{agent} is an active learning architecture --- prior, belief-update rule, action rule, preferences --- that operates in isolation from any specific environment. A \emph{coupling} is a joint property of an observer, an agent, and an environment that says whether the agent's actual operation, viewed through the observer, reaches the observer's ceiling. Section~\ref{sec:ontology} states the ontology formally; Parts~\ref{part:observer}--\ref{part:reach} develop each layer in turn and extend the framework into limits, reductions, and applied realization.

\medskip
\noindent
At the observer layer, four nested sets $\{p^\star\}\subseteq[p^\star]\subseteq R_\pi(p^\star)\subseteq R_{h_t}(p^\star)$ collect the indistinguishability classes of four observers at increasing observational power. The strict-hierarchy theorem records explicit witness conditions for strictness, and the observable-quotient theorem identifies the \emph{interventional semantic class} $[p^\star]$ as the ceiling for exactly policy-view recoverable targets. At the agent layer, the \emph{algorithmic free energy} --- the Gibbs-form variational functional whose reference measure is the Kolmogorov-complexity-weighted universal interactive prior of \citet{Solomonoff1964a,Solomonoff1964b,Kolmogorov1965,Hutter2005} rather than a stipulated parametric prior --- borrows the Gibbs variational form from active inference and the universal-prior reference from Solomonoff and Hutter. On the belief-evidence component these together give an exact Bayes/Gibbs identity: the variational decomposition is exact (Lemma~\ref{lem:variational}), the associated minimizer propositions identify the exact Bayes/Gibbs posterior, and the KL divergence is forced within the natural convex class by convex duality. The action-side boundary proxy records the risk-minus-gain shape used to compare expected-free-energy and information-gain rules \citep{Friston2010,Lindley1956,ChalonerVerdinelli1995,MacKay1992,OrtegaBraun2013,Schmidhuber1991,Schmidhuber2010}. These layers are environment-independent.

\medskip
\noindent
The central question is at the coupling layer, and the central finding is that the natural unification fails. The Bayes/Gibbs/expected-free-energy architecture under the universal prior --- the most natural rule the agent layer suggests --- does not in general converge to the interventional class: a substituted finite-support prior freezes the posterior on $[p^\star]$ at its prior mass for every finite horizon (Section~\ref{sec:near-miss}). This negative result disciplines what a positive theorem must look like. The main theorem (Theorem~\ref{thm:separating-support-convergence}) identifies the kernel-functional refinement that does succeed: when an observer-agent-environment triple satisfies a small set of explicit coupling conditions --- exact Bayes/Gibbs belief, exact kernel-functional action against the canonical predictive-kernel observer, well-formed target-fiber normalizers, the true-law prefix-Hellinger obligations, and class-predictive Bhattacharyya separation on the reference support --- the posterior converges almost surely to the packaged predictive semantic target. The semantic-to-true-law calibration step is internal to the theorem for the canonical predictive-kernel observer, removing what would otherwise be a free-standing coherence assumption. Matching converses establish that the conditions are sharp.

\medskip
\noindent
Two engineering payoffs make the framework operational. A theorem-facing differentiable surrogate (Section~\ref{sec:realization}) realizes the verified package's coupling conditions as a loss function trainable in deep learning: under three explicit deployment-side hypotheses on the learned model, minimizing the surrogate produces an agent that satisfies the main theorem. A three-question diagnostic (Section~\ref{sec:diagnostic}) turns the framework into a procedure for evaluating any specific interactive learner --- a language model, an active-inference controller, a Solomonoff-style universal predictor --- by sorting failures into observer, agent, and coupling layers and pointing each diagnosis to the canonical theorem and remedy. The framework's reach also covers the apparent uncomputability of the behavioral observer, dissolved in Section~\ref{sec:self-ref} through a finite-time proxy and an eventual-isolation interface.

\medskip
\noindent
A public Lean 4 repository at \url{https://github.com/Populustremuloides/semantic_convergence} contains Lean counterparts for the manuscript items over the explicit concrete foundation used in the repo, together with generated coverage, proof-shape, and per-declaration axiom audits. At the agent layer, \path{SemanticConvergence.def_afe}, \path{SemanticConvergence.def_two_observer_functional}, and \path{SemanticConvergence.def_kernel_functional} implement the repaired Gibbs objectives, while \path{SemanticConvergence.lem_variational}, \path{SemanticConvergence.lem_kl_necessity}, \path{SemanticConvergence.prop_two_observer_minimizer}, and \path{SemanticConvergence.prop_kernel_functional_minimizer} certify the corresponding exact variational identities and minimizer theorems. At the coupling layer, the verified main theorem is \path{SemanticConvergence.Ontology.Coupling.h10_verified_semantic_learning_package_converges}; its bridge uses \path{SemanticConvergence.countablePredictiveKernelObserver} and the \path{..._predictiveKernelObserver} semantic-calibration lemmas to remove raw coherence assumptions. The stronger zero-out companion route is packaged explicitly by \path{SemanticConvergence.ZeroOutRatePackage} and the package-level residual/posterior-share wrappers in \path{SemanticConvergence.Rates} and \path{SemanticConvergence.Noise}.
\end{abstract}

\section{Introduction}
\label{sec:intro}

Prediction admits many successes that explanation does not. An interactive learner may reproduce a realized transcript, forecast the next observation along the policy it happens to follow, and reduce uncertainty about the latent program in general, and still be wrong about what the environment would do under an intervention it never tried. In an interactive problem, that residual ambiguity is the question worth asking: whether the learner's internal inference converges to the environment law that is correct under every admissible intervention.

The question splits cleanly into three layers, and the rest of the paper is organized around them. An \emph{observer} fixes what is in principle distinguishable through a given access pattern; the four nested sets above are the indistinguishability classes of four observers at increasing observational power, and Part~\ref{part:observer} develops the observer side of the question. An \emph{agent} is a learning architecture --- prior, belief-update rule, action rule, preferences --- that the paper formulates as the algorithmic free energy and its action lift; Part~\ref{part:agent} establishes what such an architecture does, in isolation from any specific environment. A \emph{coupling} is a joint property of an observer, an agent, and an environment that says whether the agent's actual operation, viewed through the observer, reaches the observer's ceiling. The convergence theorem, zero-out converses, engineered realizations, and boundary witness interfaces all live at the coupling layer (Part~\ref{part:coupling}). Section~\ref{sec:ontology} states the three-layer ontology formally and assigns each named result of the paper to its layer; Section~\ref{sec:ontology-diagnostic} turns the layers into a diagnostic checklist for failures of any specific learner.

Four targets organize the answer for computable interactive environments. The \emph{history-compatible set} $R_{h_t}(p^\star)$ consists of programs that reproduce the realized transcript; the \emph{policy-predictable set} $R_\pi(p^\star)$, of those that induce the same law as $p^\star$ under the policy the learner actually follows; the \emph{interventional semantic class} $[p^\star]$, of those that induce the same action-conditioned law under every admissible intervention; and the innermost target, the singleton $\{p^\star\}$, is the truth. These form a nested hierarchy,
\[
\{p^\star\}
\;\subseteq\;
[p^\star]
\;\subseteq\;
R_\pi(p^\star)
\;\subseteq\;
R_{h_t}(p^\star)
\;\subseteq\;
\Pcal,
\]
and all three inclusions between the truth and the full compatibility set are strict in general (Theorem~\ref{thm:strict-hierarchy}). An observable-quotient theorem (Theorem~\ref{thm:factor-through-quotient}) makes $[p^\star]$ the theorem-level ceiling: any target a history-based learner can reliably recover is constant on $[p^\star]$, so concentration on $\{p^\star\}$ itself is impossible whenever a semantic class contains more than one program. The four sets are the indistinguishability classes of four \emph{observers} at increasing observational power (Definition~\ref{def:observer}), from reading only the realized transcript up to reading the source code; the paper's theorems turn on whether a learning rule \emph{promotes} the realized observer toward the behavioral one.

Three existing traditions supply natural ingredients for a learner aimed at that ceiling. Solomonoff induction and AIXI supply a universal prior over computable interactive environments, weighted by Kolmogorov complexity \citep{Solomonoff1964a,Solomonoff1964b,Kolmogorov1965,Hutter2005}, but no variational functional. The free energy principle and active inference supply the Gibbs variational form and a risk-plus-epistemic-value action rule \citep{Friston2010,FristonEtAl2017Curiosity,ParrFriston2019}, but with a stipulated parametric prior whose justification has been debated on philosophical grounds \citep{ColomboWright2021}. Bayesian experimental design, information-theoretic bounded rationality, and compression-based curiosity supply closely related information-seeking action criteria, again without a universal-prior reference measure \citep{Lindley1956,ChalonerVerdinelli1995,MacKay1992,OrtegaBraun2013,Schmidhuber1991,Schmidhuber2010}. No existing framework couples these three ingredients into a single variational object. We formulate the \emph{algorithmic free energy} --- the Gibbs-form variational functional whose reference measure is the Kolmogorov-complexity-weighted universal interactive prior rather than a stipulated parametric prior (Definition~\ref{def:afe}) --- and on it the three traditions agree on the belief-evidence component: Lemma~\ref{lem:variational} gives the exact Bayes/Gibbs variational identity, while Proposition~\ref{prop:two-observer-minimizer} and Proposition~\ref{prop:kernel-functional-minimizer} identify the exact Bayes/Gibbs posterior inside the two action-side Gibbs objectives. The KL-based divergence is forced within the natural convex class (Lemma~\ref{lem:kl-necessity}); minimizing the corresponding Gibbs objectives on this component is exact Bayesian updating, not an approximation. The boundary expected-free-energy proxy of Section~\ref{sec:near-miss} records the formal risk-minus-gain shape used to compare information-gain-style action criteria. We do not claim that the broader free-energy principle is subsumed --- the action side and the homeostatic and policy-as-inference readings of free energy are outside the scope of the present claims.

The unifications do not close the problem. A learner can still fail at the coupling layer if its realized policy, reference law, or prefix-Hellinger geometry does not satisfy the bundled convergence assumptions. Section~\ref{sec:near-miss} records the Lean-facing boundary interface for such failures: if a boundary action witness, a frozen-posterior-through-horizon predicate, and an observer-promotion failure predicate are supplied, Lean packages them into the corresponding near-miss and promotion-contrast conclusions. This is a formal witness interface for the AIXI near-miss intuition \citep{LeikeHutter2015,Leike2016,LeikeEtAl2016}, not a standalone construction of a universal-prior counterexample.

The paper has four interlocking deliverables. First, a three-layer ontology (Section~\ref{sec:ontology}) that separates what is in principle distinguishable (observer) from what learning machinery does (agent) from when the two meet productively in a given environment (coupling). Second, a coupling-layer convergence theorem (Theorem~\ref{thm:separating-support-convergence}) at a kernel-functional refinement of the natural Bayes/Gibbs/expected-free-energy architecture, with sharpness converses on both directions of the support condition. Third, a negative finding (Theorem~\ref{thm:afe-near-miss}) that the natural unification itself does not converge to the interventional class --- which is what forces the kernel-functional refinement to be the specific shape the theorem identifies. Fourth, two engineering payoffs: a differentiable surrogate loss whose minimizer realizes the verified coupling conditions (Section~\ref{sec:realization}), and a three-question diagnostic procedure that applies the framework to any specific interactive learner (Section~\ref{sec:diagnostic}).

\problembox{\textbf{Problem.} Give sufficient conditions on the learner's prior, belief rule, action rule, and semantic target package under which posterior mass on the target semantic fiber converges almost surely to one.}

The main theorem answers the boxed problem through the exact kernel-functional route. The answer is packaged as a matched triple $\mathfrak K=(\mathcal I,\mathcal J,\Scal)$: $\mathcal I$ is the realized Bayes/Gibbs learning instance, $\mathcal J$ is the kernel functional minimized at each history, and $\Scal(p^\star)$ is the predictive semantic fiber whose posterior mass should converge to one. Theorem~\ref{thm:separating-support-convergence} says that when this triple satisfies the single formal package Assumptions~\ref{ass:main-verified-package}, exact Bayes/Gibbs belief plus exact kernel-functional minimization drives the posterior onto $\Scal(p^\star)$. The semantic-to-true-law calibration step is internal to the theorem: the canonical predictive-kernel observer discharges what would otherwise be a free-standing coherence assumption. A zero-out support/rate companion theorem provides finite-horizon rate bounds; its witness packages diagnose two coupling failure modes, wrong-observer geometry and insufficient separating support.

Section~\ref{sec:setup} fixes the interactive setting. Section~\ref{sec:ontology} states the three-layer ontology --- observer, agent, coupling --- that organizes the rest of the paper. Part~\ref{part:observer} (Section~\ref{sec:hierarchy}) is the observer side: it constructs the four-set hierarchy, proves the observable-quotient theorem, and records witness conditions for the three strict inclusions. Part~\ref{part:agent} (Sections~\ref{sec:functional}--\ref{sec:belief}) is the agent side: Section~\ref{sec:functional} introduces the two-observer functional $\mathcal J_t^{\omega_B,\omega_A}$ together with its exact action-realizing kernel lift $\mathcal J_{t,\mathrm{ker}}^{\omega_B,\omega_A}$ and proves the formal admissibility/meeting-point wrappers; Section~\ref{sec:belief} identifies the belief-side minimizer at the canonical belief observer and proves the variational and KL-necessity identities. Part~\ref{part:coupling} (Sections~\ref{sec:semantic-action}--\ref{sec:main}) develops the convergence theory: Section~\ref{sec:semantic-action} records the exact kernel action law, finite-list support proxies, and verified channel/noise bookkeeping; Section~\ref{sec:main} states the formal coupling conditions on the package $\mathfrak K=(\mathcal I,\mathcal J,\Scal)$ and proves the main theorem from them, with sharpness converses, a zero-out specialization, and rate certificates. Part~\ref{part:reach} (Sections~\ref{sec:near-miss}--\ref{sec:diagnostic}) extends the framework's reach: Section~\ref{sec:near-miss} proves the near-miss --- the most natural information-gain rule fails to converge to $[p^\star]$ --- which disciplines what the main theorem must look like; Section~\ref{sec:self-ref} replaces the uncomputable behavioral observer by a finite-time proxy under eventual isolation; Section~\ref{sec:realization} exhibits a differentiable surrogate loss whose minimizer realizes the verified package; Section~\ref{sec:diagnostic} turns the framework into a procedure for diagnosing any specific learner. Sections~\ref{sec:discussion} and~\ref{sec:conclusion} draw consequences.

\section{Formal Setup}
\label{sec:setup}

Let $\Acal$ be a finite action alphabet and $\Ocal$ a countable observation alphabet. Write $\Delta(X)$ for the set of probability measures on a measurable space $X$, and equip every finite or countable set appearing below with its discrete $\sigma$-algebra. For each $t\ge 0$ write
\[
\Hcal_t:=(\Acal\times\Ocal)^t,
\qquad
\Hcal_0:=\{\varepsilon\},
\qquad
\Hcal:=\bigsqcup_{t\ge 0}\Hcal_t,
\]
with the product and disjoint-union $\sigma$-algebras. Time is discrete. At step $t$ the learner emits an action $a_t\in\Acal$ and receives an observation $o_t\in\Ocal$. Histories are
\[
h_t=(a_1,o_1,\dots,a_t,o_t),
\qquad
h_0=\varepsilon.
\]
A history $h_t$ is an element of $\Hcal_t$. A policy is a sequence of measurable stochastic kernels
\[
\pi_t:\Hcal_{t-1}\rightsquigarrow \Acal,
\qquad
h_{t-1}\longmapsto \pi_t(a\mid h_{t-1})\text{ measurable for each }a\in\Acal.
\]
For $S\subseteq\Acal$ write
\[
\pi_t(S\mid h_{t-1}):=\sum_{a\in S}\pi_t(a\mid h_{t-1}).
\]
Actions are interventions and observations are evidence. Under an environment $p$ and a policy $\pi$, write $(A_t,O_t)_{t\ge 1}$ for the resulting action-observation process and
\[
H_t:=(A_1,O_1,\dots,A_t,O_t)
\]
for the random history at time $t$.

Fix a universal prefix machine $U$. Let $\Pcal$ be the countable set of programs in the prefix-free domain of $U$ that induce interactive environment models. Each $p\in\Pcal$ defines an action-conditioned chronological law
\[
\mu_p(o_{1:t}\givena a_{1:t})
=
\prod_{i=1}^{t}\mu_p(o_i\mid h_{i-1},a_i),
\qquad
\sum_{o\in\Ocal}\mu_p(o\mid h_{t-1},a_t)=1.
\]
For $h_t=(a_1,o_1,\dots,a_t,o_t)\in\Hcal_t$, the induced finite-horizon law under $(p,\pi)$ is
\[
\Prob_{p,\pi}(H_t=h_t)
:=
\mu_p(o_{1:t}\givena a_{1:t})\prod_{i=1}^t \pi_i(a_i\mid h_{i-1}),
\]
which defines a probability measure $\Prob_{p,\pi}(H_t\in\cdot)\in\Delta(\Hcal_t)$. The consistent family of finite-horizon laws determines the full trajectory law $\Prob_{p,\pi}$ on the product space $(\Acal\times\Ocal)^{\mathbb N}$. Call a history-action pair $(h_t,a)$ \emph{reachable under $p$} if $\mu_p(o_{1:t}\givena a_{1:t})>0$. Throughout, kernel equality $\mu_p=\mu_{p'}$ denotes agreement at pairs reachable under at least one of $p,p'$: $\mu_p(\cdot\mid h,a)=\mu_{p'}(\cdot\mid h,a)$ whenever $(h,a)$ is reachable under $p$ or under $p'$; jointly unreachable pairs have no observable consequence under any policy. This reachable-pair equality defines an equivalence relation on $\Pcal$; write
\[
p\equiv p'
\iff
\mu_p=\mu_{p'},
\qquad
\Ccal:=\Pcal/\!\equiv,
\qquad
[\cdot]:\Pcal\to\Ccal
\]
for the semantic quotient set and quotient map.

For a prior $w$ on $\Pcal$ and a history $h_t$ with positive $w$-induced marginal likelihood, the Bayes posterior over programs is
\[
q_t(p\mid h_t)
:=
\frac{w(p)\mu_p(o_{1:t}\givena a_{1:t})}
{\sum_{r\in\Pcal} w(r)\mu_r(o_{1:t}\givena a_{1:t})}.
\]
For any subset $\Rcal\subseteq\Pcal$ write
\[
q_t(\Rcal\mid h_t):=\sum_{p\in\Rcal} q_t(p\mid h_t).
\]
For a semantic class $c\in\Ccal$, write
\[
q_t(c\mid h_t):=q_t([\cdot]^{-1}(c)\mid h_t)=\sum_{p\in c} q_t(p\mid h_t).
\]
The prior $w$ pushes forward to $\Ccal$ as
\[
\bar w(c):=w([\cdot]^{-1}(c))=\sum_{p\in c}w(p).
\]
For a semantic class $c=[p]\in\Ccal$, write $\mu_c:=\mu_p$ for the common action-conditioned law induced by any representative of $c$; this is well-defined by definition of $\equiv$.
For a continuation $(a,o)$ after $h_t$ with positive predictive normalizing constant,
\[
q_{t+1}(p\mid h_t,a,o)
:=
\frac{q_t(p\mid h_t)\mu_p(o\mid h_t,a)}
{\sum_{r\in\Pcal} q_t(r\mid h_t)\mu_r(o\mid h_t,a)}.
\]
For a semantic class $c\in\Ccal$, write
\[
q_{t+1}(c\mid h_t,a,o):=q_{t+1}([\cdot]^{-1}(c)\mid h_t,a,o)=\sum_{p\in c} q_{t+1}(p\mid h_t,a,o).
\]

\section{Observers, Agents, and Coupling}
\label{sec:ontology}

The contributions of this paper sort cleanly into three layers, and the rest of the manuscript reads more transparently once the layers are named. An \emph{observer} fixes what is in principle distinguishable. An \emph{agent} fixes the learning machinery that tries to exploit that distinguishability. A \emph{coupling} is a joint property of an observer, an agent, and an environment that determines whether learning actually reaches the observer's ceiling. Each layer answers a different question, and each is the natural home of a different family of results below.



\begin{center}
\fbox{\parbox{0.93\textwidth}{\centering
\textbf{The three layers.}\\[0.4em]
\begin{tabular}{@{}p{0.27\textwidth}p{0.62\textwidth}@{}}
\textbf{Observer} & a measurable map $\omega:\Pcal\to\Omega^\omega$; a passive lens that determines what is in principle distinguishable.\\[0.3em]
\textbf{Agent} & a learning architecture $(w,\{q_t\},\{\pi_t\},\rho,c)$; an active inference rule together with its action policy.\\[0.3em]
\textbf{Coupling} & a property of a triple $(\omega,\Acal,p^\star)$ jointly: whether the agent's actual operation in this environment, viewed through the observer, reaches the observer's ceiling.\\
\end{tabular}
}}
\end{center}

\subsection{Observers}
\label{sec:ontology-observers}

An \emph{observer} is a measurable function
\[
\omega:\Pcal\longrightarrow\Omega^\omega,
\]
where $\Omega^\omega$ is a measurable space of \emph{views}. Two programs are \emph{$\omega$-equivalent} if they produce the same view, $p\sim_\omega p'\iff\omega(p)=\omega(p')$, and the \emph{observer fiber} of $p$ is $[p]_\omega:=\{p'\in\Pcal:\omega(p)=\omega(p')\}$. An observer determines an equivalence relation on $\Pcal$, a partition $\Pcal/\!\sim_\omega$, and the pushforward of any prior or posterior to that partition. It is a \emph{lens}: a passive labelling rule. Observers do not act, do not have priors of their own, and do not depend on any environment. The four-observer chain
\[
\omega_{h_t}\preceq\omega_\pi\preceq\omega_{\mathrm{behav}}\preceq\omega_{\mathrm{syn}}
\]
of Section~\ref{sec:hierarchy} is a refinement lattice: each later observer is a refinement of the previous one in the sense that its fibers are contained in the previous one's fibers. The interventional observer $\omega_{\mathrm{behav}}$ is the most refined view a history-based learner can possibly resolve; finer observers like $\omega_{\mathrm{syn}}$ identify programs that no admissible policy can distinguish.

The kind of question an observer is the answer to is: \emph{what can be told apart, in principle, through this access?} The observer-side results in this paper are exactly the answers to that question:
\begin{itemize}
\item The four-set hierarchy and the strict-hierarchy theorem (Section~\ref{sec:hierarchy}) say which inclusions are real.
\item The observable-quotient theorem (Section~\ref{sec:hierarchy}) says no history-based learner can resolve below~$\omega_{\mathrm{behav}}$, which is why $[p^\star]:=[p^\star]_{\omega_{\mathrm{behav}}}$ is the right semantic target.
\end{itemize}
None of these results mention priors, update rules, or actions. They are facts about lenses, not about learners.

\subsection{Agents}
\label{sec:ontology-agents}

An \emph{agent} is a learning architecture. Formally, an agent is a tuple
\[
\Acal=\bigl(w,\,\{q_t\}_{t\ge0},\,\{\pi_t\}_{t\ge1},\,\rho,\,c\bigr)
\]
consisting of a prior $w\in\Delta(\Pcal)$, a sequence of belief updates $q_t:\Hcal_t\to\Delta(\Pcal)$, a policy $\pi$ as in Section~\ref{sec:setup}, a preference profile $\rho_t$, and an action-cost profile $c_t$. An agent acts: at each step it consumes a history, produces an action, receives an observation, and revises beliefs. The variational machinery of this paper operates entirely within an agent: the universal prior, the algorithmic free energy $\Fcal_t$, the two-observer functional $\mathcal J_t^{\omega_B,\omega_A}$, and the kernel functional $\mathcal J_t^{\omega_B,\omega_A,\kappa}$ are all components of an agent's belief and action rules.

The kind of question an agent is the answer to is: \emph{what does the learning architecture do at each step, holding everything else fixed?} The agent-side results in this paper are answers to that question:
\begin{itemize}
\item The variational identity (Section~\ref{sec:belief}) decomposes algorithmic free energy into the exact Bayes/Gibbs evidence term plus the posterior gap, and the minimizer propositions identify zero posterior gap with the Bayes/Gibbs posterior. These are structural facts about the belief update; they hold in every environment.
\item The KL-necessity theorem (Section~\ref{sec:belief}) singles out $\KL(q\|w)$ within the natural convex class as the unique divergence forcing the exact Gibbs variational identity. This is a structural fact about the divergence used by the belief update.
\item The boundary expected-free-energy proxy (Corollary~\ref{cor:efe-specialization} in Section~\ref{sec:near-miss}) records the risk-minus-gain specialization used to compare information-gain-style action rules. This is a structural proxy for the action rule.
\end{itemize}
Two qualitatively distinct \emph{action-rule classes} appear within an agent and recur throughout the paper. \emph{Information-gain rules} (Section~\ref{sec:near-miss}) reward observations that reduce the agent's uncertainty about the latent program; expected free energy under uniform preferences is the canonical instance. \emph{Class-against-complement rules} (Section~\ref{sec:semantic-action}) reward observations that distinguish a target's class from its complement under the true environment law. The relationship between the two is the substantive content of Section~\ref{sec:near-miss} and the Principle stated there: information gain is necessary but not sufficient; class-against-complement separation is what actually drives convergence to the interventional class.

None of these results mention an environment $p^\star$. They are facts about the agent in isolation.

\subsection{Coupling}
\label{sec:ontology-coupling}

A \emph{coupling} is a property of a triple $(\omega,\Acal,p^\star)$ --- an observer, an agent, and an environment --- that none of the three determines on its own. Coupling conditions answer the question: \emph{does the agent's actual operation in this environment, viewed through this observer, reach the observer's ceiling?}

The central coupling object in this paper is the class-predictive Bhattacharyya score
\[
\bar B_t^{\omega_A}(c,a;h_t)
=
1-\sum_{o\in\Ocal}\sqrt{\Qcal_t^{O,\omega_A}(o\mid c,h_t,a)\,\mu_{p^\star}(o\mid h_t,a)},
\]
which is computed at the observer $\omega_A$, against the environment's true law $\mu_{p^\star}$, under the agent's class-predictive law $\Qcal_t^{O,\omega_A}$. Removing any of the three --- observer, environment, agent's law --- leaves the score undefined. Whether the score is bounded below by some $\kappa>0$ on the agent's policy support is the coupling condition that drives Section~\ref{sec:main}'s convergence theorem. Whether the agent's cumulative policy mass on separating actions diverges is a second coupling condition. Whether the agent's reference law is calibrated to the environment's prefix-Hellinger geometry is a third. None is a property of any single layer.

The coupling-layer results are:
\begin{itemize}
\item The main convergence theorem (Section~\ref{sec:main}) is a coupling-success theorem: when the observer, agent, and environment are represented by a verified semantic-learning package satisfying the bundled assumptions, the agent's posterior on the packaged semantic target $\Scal(p^\star)$ converges almost surely to one.
\item The near-miss witness package (Section~\ref{sec:near-miss}) is a coupling-failure interface: it packages supplied boundary-action, frozen-posterior, and observer-promotion failure predicates into the corresponding formal conclusions.
\item The matching converses (zero separating support; summable separating support) are further coupling failures, this time of the support-mass kind rather than the score kind.
\item The realization wrappers (Section~\ref{sec:semantic-action}) and the differentiable surrogate (Section~\ref{sec:amortized-surrogate}) are coupling recipes: once their package assumptions are verified, the main theorem supplies convergence.
\end{itemize}
These results live at the coupling layer because each of them is sensitive to a change in any one of the three constituents. A coupling that succeeds with one environment can fail with another; an agent that couples well with one environment can fail with another; an observer too coarse to resolve $p^\star$ can render any agent's coupling vacuous.

\subsection{Locating the paper's contributions}
\label{sec:ontology-locator}

Table~\ref{tab:ontology} sorts the main results by layer.

\begin{table}[h]
\centering\small
\begin{tabular}{p{0.30\linewidth}p{0.62\linewidth}}
\toprule
Layer & Results in this paper \\
\midrule
Observer & Strict-hierarchy theorem; observable-quotient theorem; the identification of $[p^\star]_{\omega_{\mathrm{behav}}}$ as the right interventional target. \\
Agent & Variational identity for the algorithmic free energy; KL-necessity within the natural convex class; boundary expected-free-energy proxy; formal admissibility/meeting-point wrappers; the two-observer and kernel functionals as agent components. \\
Coupling & Main convergence theorem; matching converses (zero and summable separating support); boundary near-miss witness package; verified realizations and the amortized surrogate as engineered couplings; self-referential reduction as a coupling-engineering technique. \\
\bottomrule
\end{tabular}
\caption{The paper's contributions sorted by ontological layer. Observer-side results say what can be distinguished. Agent-side results say what learning machinery does in isolation. Coupling-side results say when an agent in an environment, viewed through an observer, actually converges.}
\label{tab:ontology}
\end{table}

\subsection{Diagnostic use}
\label{sec:ontology-diagnostic}

The three layers also offer a diagnostic vocabulary for failures of an interactive learner. If a learner fails to concentrate on the right semantic class, the failure must lie in at least one of the three layers, and the layer is the right starting place for the diagnosis:

\begin{itemize}
\item \emph{Observer failure.} The learner's observer is too coarse (or too fine) for the environment of interest. The learner cannot, even in principle, resolve $p^\star$'s class through this lens. The remedy is observer-side: change the access pattern, refine the equivalence relation, or replace the observer entirely. Section~\ref{sec:hierarchy} says this is the only kind of failure that history-only learners cannot avoid below~$\omega_{\mathrm{behav}}$.
\item \emph{Agent failure.} The learner's belief update or action rule is structurally inadequate, even apart from any environment. A belief rule that does not fit the Bayes/Gibbs variational identity, or an action rule that cannot be represented by the formal risk-minus-gain/action-kernel interfaces, may fail to converge for reasons internal to the agent. Section~\ref{sec:belief} characterizes the unique divergence within the natural convex class; agents that pick a different one are doing structural work that this paper rules out as inadequate.
\item \emph{Coupling failure.} The observer can resolve the right class and the agent has the right architecture, but the agent's policy in this environment, viewed through this observer, does not satisfy the convergence conditions. The near-miss witness package records the formal shape of such a failure once its predicates are supplied. The matching converses are the support-mass version of the same phenomenon. The remedy is coupling-side: engineer the policy or the reference law (the verified realizations and the surrogate are recipes for doing exactly this).
\end{itemize}

A reader looking to apply this paper's results to a specific learner --- a language model evaluated on a specific task, an active-inference controller, or a Solomonoff-style universal predictor --- can use the layer split as a checklist. Is the observer fine enough to resolve the relevant program class? Is the agent architecture in the natural variational class? Is the coupling between agent and environment, viewed through the observer, separating? The convergence theorem of Section~\ref{sec:main} covers the third question; the answers to the first two are prerequisites for it to apply. Appendix~\ref{app:diagnostic} turns this vocabulary into an operational checklist with a symptom-to-remedy table and a worked example.

\subsection{What this changes about the rest of the paper}
\label{sec:ontology-roadmap}

The remaining sections develop one layer at a time. Section~\ref{sec:hierarchy} is observer-only: it establishes the four-set lattice, the observable-quotient theorem, and the strict-hierarchy witness package. Section~\ref{sec:functional} introduces the two-observer functional, which is the variational shape of an agent's belief and class-action rules; the meeting-point statements there are formal admissibility wrappers around the behavioral observer. Section~\ref{sec:belief} is agent-only: it isolates the belief side of the agent at $\omega_B=\omega_{\mathrm{syn}}$ and proves the variational identity and KL-necessity theorems. Section~\ref{sec:semantic-action} develops the action-side interfaces at $\omega_A=\omega_{\mathrm{behav}}$ and gives package-realization wrappers. Section~\ref{sec:main} is the central coupling section: it states the standing coupling conditions and proves the main convergence theorem from them. The boundary section records the risk-minus-gain proxy and near-miss witness package. The self-referential reduction of Section~\ref{sec:self-ref} is a coupling-engineering technique expressed through finite-time proxy predicates.

Throughout the rest of the paper, when a result depends on the environment, it is a coupling result; when it does not, it is an observer or agent result. The reader can use this distinction to pick which kind of remedy applies when an instantiation falls short of any specific theorem's hypotheses.

\part{The Observer Side: What Is Distinguishable}
\label{part:observer}

Part I develops the observer side of the ontology of Section~\ref{sec:ontology}. Observers are passive lenses; they fix what can in principle be told apart through a given access pattern, but say nothing about the learning machinery that exploits the access. The two main results of this part --- the strict-hierarchy theorem and the observable-quotient theorem --- answer the question \emph{what is the ceiling of distinguishability for a history-based learner?} They make no reference to priors, update rules, or actions. Subsequent parts add the learning machinery (Part~\ref{part:agent}) and the coupling conditions (Part~\ref{part:coupling}) on top of this fixed ceiling.

\section{The Knowability Hierarchy}
\label{sec:hierarchy}

This section replaces a list of informal learning objectives with four nested sets of programs and records explicit witness conditions under which each inclusion is strict. The strictness witnesses are the structural spine of the paper.

\subsection{Four nested sets}

\begin{definition}[History-compatible set]
\label{def:history-compat}
For a target program $p^\star\in\Pcal$ and a history $h_t\in\Hcal_t$, the \emph{history-compatible set} is
\[
R_{h_t}(p^\star)
:=
\bigl\{p\in\Pcal\,:\,\mu_p(o_{1:t}\givena a_{1:t})=\mu_{p^\star}(o_{1:t}\givena a_{1:t})\bigr\}.
\]
\end{definition}

\begin{definition}[Policy-predictable set]
\label{def:policy-pred}
For a target program $p^\star\in\Pcal$ and a policy $\pi$, the \emph{policy-predictable set} is
\[
R_\pi(p^\star)
:=
\bigl\{p\in\Pcal\,:\,\Prob_{p,\pi}(H_T\in\cdot)=\Prob_{p^\star,\pi}(H_T\in\cdot)\text{ for every }T\bigr\}.
\]
\end{definition}

\begin{definition}[Interventional semantic class]
\label{def:int-sem-class}
For a target program $p^\star\in\Pcal$, the \emph{interventional semantic class} of $p^\star$ is its $\equiv$-class
\[
[p^\star]
=
[\cdot](p^\star)
=
\bigl\{p\in\Pcal\,:\,p\equiv p^\star\bigr\}
=
\bigl\{p\in\Pcal\,:\,\mu_p=\mu_{p^\star}\bigr\}.
\]
\end{definition}

Thus $\Ccal$ is the quotient set of programs by semantic equivalence, and two programs lie in the same class exactly when they agree on the entire action-conditioned environment law.

\begin{definition}[Observer]
\label{def:observer}
An \emph{observer} is a pair $\omega=(\Omega_\omega,V_\omega)$ of a measurable space $\Omega_\omega$ and a measurable view map $V_\omega:\Pcal\to\Omega_\omega$. Observer $\omega$ is \emph{refined by} $\omega'$, written $\omega\le\omega'$, if there is a measurable factorization $\phi:\Omega_{\omega'}\to\Omega_\omega$ with $V_\omega=\phi\circ V_{\omega'}$. The \emph{$\omega$-indistinguishability class} of $p^\star$ is the fiber $V_\omega^{-1}(V_\omega(p^\star))$: the programs an adversary could substitute for $p^\star$ without the $\omega$-observer detecting the swap.
\end{definition}

Four canonical observers structure the hierarchy, each equipped with the image $\sigma$-algebra making its view map measurable: the \emph{syntactic observer} $\omega_{\mathrm{syn}}$ with $\Omega_{\omega_{\mathrm{syn}}}:=\Pcal$ and $V_{\omega_{\mathrm{syn}}}(p):=p$, reading the source code; the \emph{behavioral observer} $\omega_{\mathrm{behav}}$ with $\Omega_{\omega_{\mathrm{behav}}}:=\{\mu_p:p\in\Pcal\}$ and $V_{\omega_{\mathrm{behav}}}(p):=\mu_p$, running every policy and seeing every induced law; the \emph{policy observer} $\omega_\pi$ with $\Omega_{\omega_\pi}:=\{\Prob_{p,\pi}:p\in\Pcal\}$ and $V_{\omega_\pi}(p):=\Prob_{p,\pi}$, seeing only the law under a fixed policy; and the \emph{history observer} $\omega_{h_t}$ with $\Omega_{\omega_{h_t}}:=\{\mu_p(o_{1:t}\givena a_{1:t}):p\in\Pcal\}\subseteq[0,1]$ and $V_{\omega_{h_t}}(p):=\mu_p(o_{1:t}\givena a_{1:t})$, reading only the realized transcript. Their indistinguishability classes of $p^\star$ are exactly $\{p^\star\}$, $[p^\star]$, $R_\pi(p^\star)$, and $R_{h_t}(p^\star)$.

\begin{remark}
\label{rem:four-observers}
The hierarchy is the partition-refinement chain of four observers at increasing observational power: as the observer is promoted, the indistinguishability class shrinks. The nesting, strict-witness, and observable-quotient results below are the three consequences of this single reading. The boundary witness predicates and the main theorem both use this observer vocabulary: the former describe how a promotion failure would be certified, while the latter gives bundled sufficient conditions for convergence.
\end{remark}

\begin{lemma}[Nesting]
\label{lem:nesting}
For any $p^\star\in\Pcal$, any policy $\pi$, and any history $h_t$ with $\Prob_{p^\star,\pi}(H_t=h_t)>0$,
\[
\{p^\star\}\subseteq [p^\star]\subseteq R_\pi(p^\star)\subseteq R_{h_t}(p^\star)\subseteq \Pcal.
\]
\end{lemma}

\begin{proof}
\noindent\emph{Lean: \path{SemanticConvergence.lem_nesting}.}
A swap the stronger observer cannot detect, the weaker observer cannot detect. If $p\in[p^\star]$ then $\mu_p=\mu_{p^\star}$, so the induced $\pi$-law matches $p^\star$'s, giving $[p^\star]\subseteq R_\pi(p^\star)$. If $p\in R_\pi(p^\star)$ then $\Prob_{p,\pi}(H_t=h_t)=\Prob_{p^\star,\pi}(H_t=h_t)$; combining with the factorization
\[
\Prob_{p,\pi}(H_t=h_t)=\mu_p(o_{1:t}\givena a_{1:t})\prod_{i=1}^t \pi_i(a_i\mid h_{i-1}),
\]
and dividing through by the policy factor (positive by assumption), one gets $\mu_p(o_{1:t}\givena a_{1:t})=\mu_{p^\star}(o_{1:t}\givena a_{1:t})$, hence $p\in R_{h_t}(p^\star)$. The outer inclusion is immediate.
\end{proof}

\begin{proposition}[Refinement chain]
\label{prop:refinement-chain}
Under the hypotheses of Lemma~\ref{lem:nesting},
\[
\omega_{h_t}\le\omega_\pi\le\omega_{\mathrm{behav}}\le\omega_{\mathrm{syn}}.
\]
\end{proposition}

\begin{proof}
\noindent\emph{Lean: \path{SemanticConvergence.prop_refinement_chain}.}
The factorization $p\mapsto\mu_p$ witnesses $\omega_{\mathrm{behav}}\le\omega_{\mathrm{syn}}$. The history-factorization identity of Lemma~\ref{lem:nesting} exhibits $\Prob_{p,\pi}$ as a fixed measurable function of $\mu_p$, giving $\omega_\pi\le\omega_{\mathrm{behav}}$. Positivity of the policy factor along $h_t$ lets one recover $\mu_p(o_{1:t}\givena a_{1:t})=\Prob_{p,\pi}(H_t=h_t)/\prod_i\pi_i(a_i\mid h_{i-1})$, giving $\omega_{h_t}\le\omega_\pi$.
\end{proof}

\subsection{Observable quotient: why $[p^\star]$ is the target}

\begin{lemma}[Interventional semantic classes are the observable quotient]
\label{lem:observable-quotient}
If $p\equiv p'$, then for every policy $\pi$ and every horizon $T$, the induced law of $H_T$ under $(p,\pi)$ equals the induced law of $H_T$ under $(p',\pi)$.
\end{lemma}

\begin{proof}
\noindent\emph{Lean: \path{SemanticConvergence.lem_observable_quotient}.}
Fix $h_T$. If $\Prob_{p,\pi}(H_T=h_T)>0$ or $\Prob_{p',\pi}(H_T=h_T)>0$, then every prefix $(h_{t-1},a_t)$ visited along $h_T$ is reachable under $p$ or $p'$, so semantic equivalence gives $\mu_p(o_t\mid h_{t-1},a_t)=\mu_{p'}(o_t\mid h_{t-1},a_t)$ for each $t$. Using the factorization
\[
\Prob_{p,\pi}(H_T=h_T)
=
\prod_{t=1}^{T}\pi_t(a_t\mid h_{t-1})\,\mu_p(o_t\mid h_{t-1},a_t),
\]
and the analogous identity for $p'$, the two probabilities are equal. If both probabilities are zero, equality is trivial. Hence the laws coincide.
\end{proof}

\begin{definition}[History-recoverable target]
\label{def:history-recoverable}
Let $\mathcal{T}$ be a target type and let $\tau:\Pcal\to\mathcal{T}$ be a target map. The target $\tau$ is \emph{history-recoverable} in the Lean-facing exact policy-view sense if there exist a policy $\pi$ and a decoder
\[
D:\bigl\{(T,h_T)\mapsto \Prob_{p,\pi}(H_T=h_T)\bigr\}\longrightarrow \mathcal{T}
\]
such that, for every $p\in\Pcal$, applying $D$ to the full policy-conditioned history-law view of $p$ returns $\tau(p)$.
\end{definition}

\begin{theorem}[Every history-recoverable target factors through the semantic quotient]
\label{thm:factor-through-quotient}
Let $\tau:\Pcal\to\mathcal{T}$ be history-recoverable in the exact policy-view sense of Definition~\ref{def:history-recoverable}. Then there exists a map $\bar\tau:\Ccal\to\mathcal{T}$ such that $\tau=\bar\tau\circ[\cdot]$. In particular, $p\equiv p'$ implies $\tau(p)=\tau(p')$.
\end{theorem}

\begin{proof}
\noindent\emph{Lean: \path{SemanticConvergence.thm_factor_through_quotient}.}
Fix $p\equiv p'$. By Lemma~\ref{lem:observable-quotient}, their policy-conditioned history-law views agree for every policy. Applying the decoder witnessing Definition~\ref{def:history-recoverable} gives $\tau(p)=\tau(p')$. The quotient-lift property for functions constant on semantic-equivalence classes then produces $\bar\tau$ with $\tau=\bar\tau\circ[\cdot]$.
\end{proof}

\begin{remark}
Theorem~\ref{thm:factor-through-quotient} is the universality of the semantic quotient expressed in the adversarial reading of Remark~\ref{rem:four-observers}: no history-based estimator can resolve distinctions finer than the full-policy observer detects. Concentration on a target finer than $[p^\star]$ --- in particular, on $\{p^\star\}$ --- asks the observer to see through what the adversary can hide by substituting a semantically equivalent twin.
\end{remark}

\subsection{Strict separation of the hierarchy}

\begin{definition}[Strict hierarchy witness package]
\label{def:strict-hierarchy-witnesses}
For a policy $\pi$, target $p^\star$, and realized history $h_t$, a \emph{strict hierarchy witness package} consists of:
\[
\begin{aligned}
&\exists p\in R_{h_t}(p^\star)\setminus R_\pi(p^\star),\\
&\exists p'\in R_\pi(p^\star)\setminus[p^\star],\\
&\exists p''\in[p^\star]\setminus\{p^\star\}.
\end{aligned}
\]
This is the formal object used by Lean to make strictness explicit.
\end{definition}

\begin{lemma}[{Fit gap: $R_\pi(p^\star)\subsetneq R_{h_t}(p^\star)$}]
\label{lem:fit-gap}
If a strict hierarchy witness package holds for $(\pi,p^\star,h_t)$, then there exists $p\in R_{h_t}(p^\star)\setminus R_\pi(p^\star)$.
\end{lemma}

\begin{proof}
\noindent\emph{Lean: \path{SemanticConvergence.lem_fit_gap}.}
Project the first field of the strict hierarchy witness package.
\end{proof}

\begin{theorem}[{Policy gap: $[p^\star]\subsetneq R_\pi(p^\star)$}]
\label{thm:policy-gap}
If a strict hierarchy witness package holds for $(\pi,p^\star,h_t)$, then there exists $p'\in R_\pi(p^\star)\setminus[p^\star]$.
\end{theorem}

\begin{proof}
\noindent\emph{Lean: \path{SemanticConvergence.thm_policy_gap}.}
Project the second field of the strict hierarchy witness package.
\end{proof}

\begin{remark}
The informal adversarial construction behind Theorem~\ref{thm:policy-gap} is to supply the second witness by choosing a program that agrees with $p^\star$ on every pair the $\pi$-fixed observer sees and disagrees only at a pair that observer never queries. The Lean theorem deliberately takes the witness package as the formal input.
\end{remark}

\begin{lemma}[{Syntactic gap: $\{p^\star\}\subsetneq [p^\star]$, typically}]
\label{lem:syntactic-gap}
If a strict hierarchy witness package holds for $(\pi,p^\star,h_t)$, then $\{p^\star\}\subsetneq[p^\star]$.
\end{lemma}

\begin{proof}
\noindent\emph{Lean: \path{SemanticConvergence.lem_syntactic_gap}.}
Project the third field of the strict hierarchy witness package and combine it with the singleton inclusion from Lemma~\ref{lem:nesting}.
\end{proof}

\begin{remark}
In the intended Solomonoff setting, the third witness is normally supplied by syntactic redundancy: append a dead code segment, reroute through a universal interpreter, or permute the program's internal representation while leaving the induced law unchanged. The formal theorem only uses the supplied witness.
\end{remark}

\begin{theorem}[Strict knowability hierarchy]
\label{thm:strict-hierarchy}
If a strict hierarchy witness package holds for $(\pi,p^\star,h_t)$, then
\[
\{p^\star\}\subsetneq [p^\star]\subsetneq R_\pi(p^\star)\subsetneq R_{h_t}(p^\star)\subseteq \Pcal.
\]
\end{theorem}

\begin{proof}
\noindent\emph{Lean: \path{SemanticConvergence.thm_strict_hierarchy}.}
The three strict inclusions are furnished by Lemma~\ref{lem:syntactic-gap}, Theorem~\ref{thm:policy-gap}, and Lemma~\ref{lem:fit-gap} respectively. Nesting is Lemma~\ref{lem:nesting}.
\end{proof}

\begin{remark}
Table~\ref{tab:hierarchy} records the role of each set in the paper.
\end{remark}

\begin{table}[ht]
\centering
\small
\caption{The four nested sets read as indistinguishability classes at increasing observational power, with the results that witness strict separation.}
\label{tab:hierarchy}
\begin{tabularx}{\textwidth}{@{}>{\raggedright\arraybackslash}p{0.20\textwidth}>{\raggedright\arraybackslash}p{0.34\textwidth}>{\raggedright\arraybackslash}X>{\raggedright\arraybackslash}p{0.16\textwidth}@{}}
\toprule
\textbf{Set} & \textbf{Observer that cannot detect a swap from $p^\star$} & \textbf{Strictly contains} & \textbf{Witnessed by} \\
\midrule
$\{p^\star\}$ & Reads the source code & Nothing & --- \\
$[p^\star]$ & Runs any policy and sees its induced law & $\{p^\star\}$ whenever $[p^\star]$ has $>1$ element & Lemma~\ref{lem:syntactic-gap} \\
$R_\pi(p^\star)$ & Sees the full law of histories under fixed $\pi$ & $[p^\star]$ when $\pi$ leaves some pair unqueried & Theorem~\ref{thm:policy-gap} \\
$R_{h_t}(p^\star)$ & Sees only the realized transcript $h_t$ & $R_\pi(p^\star)$ when some supported next action at $h_t$ is non-degenerate & Lemma~\ref{lem:fit-gap} \\
\bottomrule
\end{tabularx}
\end{table}

The rest of the paper is organized by which set the learner's posterior is trying to concentrate on. The belief side of Section~\ref{sec:belief} fixes the exact Bayes/Gibbs posterior and records its observer-fiber invariances. Sections~\ref{sec:semantic-action} and~\ref{sec:main} give the sufficient conditions under which the posterior does concentrate on the semantic target. Section~\ref{sec:near-miss} records the boundary proxy and the supplied witness predicates that certify a boundary coupling failure.

\part{The Agent Side: What Learning Architectures Do}
\label{part:agent}

Part II develops the agent side. Agents are active learning architectures; they have priors, belief update rules, action rules, and preferences, and the variational machinery of this paper sits inside them. The results of this part --- the variational identity, KL-necessity, the formal admissibility wrappers, and the prior-weight definitions --- say what the agent does at each step in isolation from any specific environment. They are structural facts about learning architectures rather than about specific learning episodes. Part~\ref{part:coupling} takes these architectures and asks when they actually converge in a given environment.

\section{The Two-Observer Functional and the Meeting-Point Theorem}
\label{sec:functional}

The problem box of Section~\ref{sec:intro} asks for a triple --- prior, belief rule, action rule --- under which the posterior mass on the semantic target converges almost surely to one. The four-set hierarchy $\{p^\star\}\subseteq[p^\star]\subseteq R_\pi(p^\star)\subseteq R_{h_t}(p^\star)$ of Section~\ref{sec:hierarchy} identifies $[p^\star]$ as the behavioral observer fiber in the intended interpretation. This section packages the paper's answer to the box as a single variational centerpiece $\mathcal J_t^{\omega_B,\omega_A}$ together with its exact action-realizing kernel lift $\mathcal J_{t,\mathrm{ker}}^{\omega_B,\omega_A}$: the former isolates the belief and class-targeting geometry, while the latter absorbs the realized action into the variational object itself. The formal meeting-point theorem records the admissibility of $\omega_{\mathrm{behav}}$ on both sides of the observer interface; the main theorem later supplies the actual convergence route through the bundled package.

The results of this section use only the observer structure of Section~\ref{sec:hierarchy} and the universal prior $w$ of Definition~\ref{def:universal-prior}; the belief-side and action-side constructions they license are identified in Sections~\ref{sec:belief} and~\ref{sec:semantic-action} respectively, and the convergence claim is settled in Section~\ref{sec:main}.

\subsection{Quotient data at an observer}
\label{sec:quotient-data}

Fix an observer $\omega$ with equivalence $\sim_\omega$ on $\Pcal$, and write $\Ccal^\omega:=\Pcal/\!\sim_\omega$ for the set of $\omega$-classes. The prior pushes forward to
\[
\bar w^\omega(c):=\sum_{p\in c}w(p),\qquad c\in\Ccal^\omega,
\]
and the Bayes posterior on $\Pcal$ marginalizes to the class posterior
\[
q_t^\star(c\mid h_t):=\sum_{p\in c}q_t^\star(p\mid h_t).
\]
When $\omega\ge\omega_{\mathrm{behav}}$, every $\omega$-class $c$ is contained in a single behavioral class, so all $p\in c$ share the same interactive law; the class likelihood $\mu_c:=\mu_p$ for any $p\in c$ is then well-defined. When $\omega<\omega_{\mathrm{behav}}$, an $\omega$-class may contain programs with distinct interactive laws, and $\mu_c$ is multi-valued at the class level.

Regardless of where $\omega$ sits relative to $\omega_{\mathrm{behav}}$, the posterior-weighted class-predictive law
\[
\mu_c^{q_t^\star}(o\mid h_t,a):=\frac{1}{q_t^\star(c\mid h_t)}\sum_{p\in c}q_t^\star(p\mid h_t)\,\mu_p(o\mid h_t,a)
\]
is well-defined for every $c$ with $q_t^\star(c\mid h_t)>0$; and so is its complement
\[
\Qcal_t^{-c,\omega}(o\mid h_t,a):=\frac{1}{1-q_t^\star(c\mid h_t)}\sum_{c'\in\Ccal^\omega\setminus\{c\}}q_t^\star(c'\mid h_t)\,\mu_{c'}^{q_t^\star}(o\mid h_t,a).
\]
When $\omega\ge\omega_{\mathrm{behav}}$, $\mu_c^{q_t^\star}=\mu_c$ and the posterior weights in the numerator cancel; we write $\mu_c$ without the posterior superscript in this regime.

\begin{definition}[Bhattacharyya separation at observer $\omega$]
\label{def:bhat-omega}
The \emph{$\omega$-class Bhattacharyya separation} of class $c\in\Ccal^\omega$ under action $a\in\Acal$ at history $h_t$ with $0<q_t^\star(c\mid h_t)<1$ is
\[
B_t^\omega(c,a;h_t):=-\log\sum_{o\in\Ocal}\sqrt{\mu_c^{q_t^\star}(o\mid h_t,a)\,\Qcal_t^{-c,\omega}(o\mid h_t,a)}\;\in\;[0,\infty].
\]
At $\omega=\omega_{\mathrm{behav}}$ we write $B_t:=B_t^{\omega_{\mathrm{behav}}}$; this is the class-complement separation geometry used by the action-side interfaces of Section~\ref{sec:semantic-action}.
\end{definition}

Finally, any policy $\pi$ whose action at step $t+1$ is generated by first sampling a class $C\in\Ccal^{\omega_A}$ and then selecting an action targeting that class induces a \emph{class-targeting distribution} $\nu_t^\pi(\cdot\mid h_t)$ on $\Ccal^{\omega_A}$. The exact variational class distribution is identified by Proposition~\ref{prop:two-observer-minimizer}; the concrete finite-list proxy in Section~\ref{sec:semantic-action} records support rather than reconstructing the selector.

\subsection{The two-observer functional}
\label{sec:two-observer-functional}

Fix positive scalars $\beta,\gamma>0$ and two observers $\omega_B,\omega_A$.

\begin{definition}[Raw policy-coupled scaffold]
\label{def:raw-two-observer-functional}
For a distribution $q$ on $\Pcal$ and a policy $\pi$, define
\[
\mathcal J_{t,\mathrm{raw}}^{\omega_B,\omega_A}[q,\pi;h_t]
\;:=\;
\underbrace{\E_q[\ellp_t(p;h_t)]+\KL(q\givena w)}_{\mathcal F_t[q;h_t]}
\;-\;\beta\,\mathcal A_t^{\omega_A}[\pi;h_t,q]
\;+\;\gamma\,\KL\!\bigl(\nu_t^\pi\,\|\,\bar w^{\omega_A}\bigr),
\]
where $\ellp_t(p;h_t):=-\log\mu_p(o_{1:t}\mid a_{1:t})$ is the history codelength, $\mathcal F_t$ is the algorithmic free energy referenced by the belief observer $\omega_B$ through its pushforward $\bar q^{\omega_B}$ (Section~\ref{sec:belief}), and
\[
\mathcal A_t^{\omega_A}[\pi;h_t,q]:=\E_{C\sim\nu_t^\pi}\bigl[B_t^{\omega_A}(C,A_\pi;h_t)\bigr]
\]
is the expected $\omega_A$-class Bhattacharyya separation under $\pi$'s class-targeting distribution on $\Ccal^{\omega_A}$.
\end{definition}

A parallel \emph{mutual-information variant} replaces the action term:
\[
\mathcal J_{t,\mathrm{raw}}^{\omega_B,\omega_A,\mathrm{MI}}[q,\pi;h_t]
:=\mathcal F_t[q;h_t]-\beta\,\I_q\!\bigl(C^{\omega_A};O\mid A_\pi,h_t\bigr),
\]
where $C^{\omega_A}$ is the random variable coding the $\omega_A$-class of the true program under $q$. At $\omega_A=\omega_{\mathrm{syn}}$, $C^{\omega_A}=P$ and this recovers the standard expected-free-energy information term. The two variants differ by a swap of one R\'enyi divergence for another: the Bhattacharyya action term is $\tfrac{1}{2}D_{1/2}$, the MI term is $D_1=\KL$ under the conventional active-inference decomposition.

\begin{definition}[Two-observer variational functional]
\label{def:two-observer-functional}
For $c\in\Ccal^{\omega_A}$ and history $h_t$, set
\[
\bar B_t^{\omega_A}(c,a;h_t):=\min\!\bigl\{1,B_t^{\omega_A}(c,a;h_t)\bigr\},
\qquad
G_t^{\omega_A}(c;h_t):=\sup_{a\in\Acal}\bar B_t^{\omega_A}(c,a;h_t)\in[0,1].
\]
For a distribution $\nu$ on $\Ccal^{\omega_A}$, define the \emph{two-observer functional}
\[
\boxed{\;
\mathcal J_t^{\omega_B,\omega_A}[q,\nu;h_t]
:=\;
\mathcal F_t[q;h_t]
-\beta\,\E_{C\sim\nu}\!\bigl[G_t^{\omega_A}(C;h_t)\bigr]
+\gamma\,\KL\!\bigl(\nu\,\|\,\bar w^{\omega_A}\bigr).
\;}
\]
\end{definition}

\begin{proposition}[Exact global minimizer of the two-observer functional]
\label{prop:two-observer-minimizer}
For fixed $h_t$ and observers $(\omega_B,\omega_A)$, the exact Bayes/Gibbs posterior $q_t^\star(\cdot\mid h_t)$ uniquely minimizes $\mathcal J_t^{\omega_B,\omega_A}[q,\nu;h_t]$ in $q$. For this $q_t^\star$ fixed, the unique minimizer in $\nu$ is
\[
\nu_t^{\mathrm G,\omega_A}(c\mid h_t)
:=
\frac{\bar w^{\omega_A}(c)\exp\!\bigl((\beta/\gamma)G_t^{\omega_A}(c;h_t)\bigr)}
{\sum_{d\in\Ccal^{\omega_A}}\bar w^{\omega_A}(d)\exp\!\bigl((\beta/\gamma)G_t^{\omega_A}(d;h_t)\bigr)}.
\]
Hence $\bigl(q_t^\star,\nu_t^{\mathrm G,\omega_A}\bigr)$ is the exact global minimizer of $\mathcal J_t^{\omega_B,\omega_A}[q,\nu;h_t]$.
\end{proposition}

\begin{proof}
\noindent\emph{Lean: \path{SemanticConvergence.prop_two_observer_minimizer}.}
The $q$-dependence of $\mathcal J_t^{\omega_B,\omega_A}$ is exactly the belief term $\mathcal F_t[q;h_t]$. Lemma~\ref{lem:variational} identifies its posterior-gap decomposition, and the Lean minimizer theorem packages the zero-gap condition at the Bayes/Gibbs posterior. For the $\nu$-part, abbreviate $G(c):=G_t^{\omega_A}(c;h_t)$ and define
\[
Z_G:=\sum_{d\in\Ccal^{\omega_A}}\bar w^{\omega_A}(d)e^{(\beta/\gamma)G(d)}.
\]
Because $G(d)\in[0,1]$ and $\sum_d\bar w^{\omega_A}(d)=1$, one has $1\le Z_G\le e^{\beta/\gamma}<\infty$. Let
\[
\nu_G(c):=\frac{\bar w^{\omega_A}(c)e^{(\beta/\gamma)G(c)}}{Z_G}.
\]
Then, for every $\nu$,
\begin{align*}
\gamma\,\KL(\nu\|\bar w^{\omega_A})-\beta\,\E_\nu[G]
&=\gamma\sum_c \nu(c)\log\frac{\nu(c)}{\bar w^{\omega_A}(c)e^{(\beta/\gamma)G(c)}}\\
&=\gamma\,\KL(\nu\|\nu_G)-\gamma\log Z_G.
\end{align*}
The right-hand side is uniquely minimized at $\nu=\nu_G$. Combining the two one-sided minimizations gives the exact global minimizer pair.
\end{proof}

\begin{remark}[Why the raw scaffold is auxiliary]
\label{rem:raw-scaffold-auxiliary}
The raw policy functional $\mathcal J_{t,\mathrm{raw}}^{\omega_B,\omega_A}[q,\pi;h_t]$ is still the geometric scaffold of the paper, but its action term is coupled to both $q$ and the realized policy $\pi$, so an exact joint minimizer need not coincide with the canonical Bayes/Gibbs-plus-separation pair. Definition~\ref{def:two-observer-functional} separates the belief side from the class-targeting action side while preserving the same observer-indexed separation geometry. Proposition~\ref{prop:two-observer-minimizer} is the exact global-minimizer statement for $\mathcal J_t^{\omega_B,\omega_A}$.
\end{remark}

\begin{definition}[Kernel lift of the two-observer functional]
\label{def:kernel-functional}
Fix a full-support reference distribution $\lambda\in\Delta(\Acal)$. For a distribution $\kappa$ on $\Ccal^{\omega_A}\times\Acal$, define
\[
\mathcal J_{t,\mathrm{ker}}^{\omega_B,\omega_A}[q,\kappa;h_t]
:=
\mathcal F_t[q;h_t]
-\beta\,\E_{(C,A)\sim\kappa}\!\bigl[\bar B_t^{\omega_A}(C,A;h_t)\bigr]
+\gamma\,\KL\!\bigl(\kappa\,\|\,\bar w^{\omega_A}\otimes\lambda\bigr).
\]
\end{definition}

\begin{proposition}[Exact global minimizer of the kernel lift]
\label{prop:kernel-functional-minimizer}
For fixed $h_t$, observers $(\omega_B,\omega_A)$, and reference distribution $\lambda$, the exact Bayes/Gibbs posterior $q_t^\star(\cdot\mid h_t)$ uniquely minimizes $\mathcal J_{t,\mathrm{ker}}^{\omega_B,\omega_A}[q,\kappa;h_t]$ in $q$. For this $q_t^\star$ fixed, the unique minimizer in $\kappa$ is
\[
\kappa_t^{\mathrm G,\omega_A}(c,a\mid h_t)
:=
\frac{\bar w^{\omega_A}(c)\lambda(a)\exp\!\bigl((\beta/\gamma)\bar B_t^{\omega_A}(c,a;h_t)\bigr)}
{\sum_{d\in\Ccal^{\omega_A}}\sum_{b\in\Acal}\bar w^{\omega_A}(d)\lambda(b)\exp\!\bigl((\beta/\gamma)\bar B_t^{\omega_A}(d,b;h_t)\bigr)}.
\]
Hence $\bigl(q_t^\star,\kappa_t^{\mathrm G,\omega_A}\bigr)$ is the exact global minimizer of $\mathcal J_{t,\mathrm{ker}}^{\omega_B,\omega_A}[q,\kappa;h_t]$.
\end{proposition}

\begin{proof}
\noindent\emph{Lean: \path{SemanticConvergence.prop_kernel_functional_minimizer}.}
The $q$-dependence is again exactly $\mathcal F_t[q;h_t]$, with the Bayes/Gibbs posterior selected by the posterior-gap part of the Lean minimizer theorem. For the $\kappa$-part, abbreviate $\bar B(c,a):=\bar B_t^{\omega_A}(c,a;h_t)$ and define
\[
Z_{\mathrm{ker}}
:=
\sum_{d\in\Ccal^{\omega_A}}\sum_{b\in\Acal}\bar w^{\omega_A}(d)\lambda(b)e^{(\beta/\gamma)\bar B(d,b)}.
\]
Since $\bar B(d,b)\in[0,1]$ and $\bar w^{\omega_A}\otimes\lambda$ is a probability measure, one has $1\le Z_{\mathrm{ker}}\le e^{\beta/\gamma}<\infty$. Let
\[
\kappa_G(c,a):=
\frac{\bar w^{\omega_A}(c)\lambda(a)e^{(\beta/\gamma)\bar B(c,a)}}{Z_{\mathrm{ker}}}.
\]
Then, for every $\kappa$,
\begin{align*}
\gamma\,\KL\!\bigl(\kappa\,\|\,\bar w^{\omega_A}\otimes\lambda\bigr)-\beta\,\E_\kappa[\bar B]
&=\gamma\sum_{c,a}\kappa(c,a)\log\frac{\kappa(c,a)}{\bar w^{\omega_A}(c)\lambda(a)e^{(\beta/\gamma)\bar B(c,a)}}\\
&=\gamma\,\KL(\kappa\|\kappa_G)-\gamma\log Z_{\mathrm{ker}}.
\end{align*}
The right-hand side is uniquely minimized at $\kappa=\kappa_G$. Combining the two one-sided minimizations gives the exact global minimizer pair.
\end{proof}

\begin{proposition}[Exact global minimizer of the compact-action kernel lift]
\label{prop:kernel-functional-minimizer-compact}
If one relaxes the standing finite-action setup and instead lets $(\Acal,d_\Acal)$ be a compact metric space equipped with a full-support Borel probability measure $\lambda$, read Definition~\ref{def:kernel-functional} measure-theoretically: $\kappa$ ranges over Borel probability measures on $\Ccal^{\omega_A}\times\Acal$, the expectation term is the corresponding integral, and the KL term is relative to the product measure $\bar w^{\omega_A}\otimes\lambda$. Assume further that for every fixed history $h_t$ and observer class $c\in\Ccal^{\omega_A}$, the map
\[
a\longmapsto \bar B_t^{\omega_A}(c,a;h_t)
\]
is Borel measurable. Then for fixed $h_t$, observers $(\omega_B,\omega_A)$, and reference measure $\lambda$, the exact Bayes/Gibbs posterior $q_t^\star(\cdot\mid h_t)$ uniquely minimizes $\mathcal J_{t,\mathrm{ker}}^{\omega_B,\omega_A}[q,\kappa;h_t]$ in $q$. For this $q_t^\star$ fixed, the unique minimizer in $\kappa$ is the Gibbs kernel $\kappa_{t,\mathrm{cts}}^{\mathrm G,\omega_A}(\cdot,\cdot\mid h_t)$ absolutely continuous with respect to $\bar w^{\omega_A}\otimes\lambda$ and density
\[
\frac{d\kappa_{t,\mathrm{cts}}^{\mathrm G,\omega_A}}{d(\bar w^{\omega_A}\otimes\lambda)}(c,a\mid h_t)
:=
\frac{\exp\!\bigl((\beta/\gamma)\bar B_t^{\omega_A}(c,a;h_t)\bigr)}
{Z_{\mathrm{ker}}^{\mathrm{cts}}(h_t)},
\]
where
\[
Z_{\mathrm{ker}}^{\mathrm{cts}}(h_t)
:=
\sum_{d\in\Ccal^{\omega_A}}\bar w^{\omega_A}(d)\int_{\Acal}
\exp\!\bigl((\beta/\gamma)\bar B_t^{\omega_A}(d,a;h_t)\bigr)\,d\lambda(a).
\]
Hence $\bigl(q_t^\star,\kappa_{t,\mathrm{cts}}^{\mathrm G,\omega_A}\bigr)$ is the exact global minimizer of the measure-theoretic kernel lift.
\end{proposition}

\begin{proof}
\noindent\emph{Lean: \path{SemanticConvergence.prop_kernel_functional_minimizer_compact}.}
The $q$-dependence is still exactly $\mathcal F_t[q;h_t]$, with the Bayes/Gibbs posterior selected by the posterior-gap part of the Lean minimizer theorem. For the $\kappa$-part, abbreviate $\bar B(c,a):=\bar B_t^{\omega_A}(c,a;h_t)$ and $\mu_0:=\bar w^{\omega_A}\otimes\lambda$. Because $0\le \bar B(c,a)\le 1$, the normalizing constant satisfies
\[
1\le Z_{\mathrm{ker}}^{\mathrm{cts}}(h_t)=\int e^{(\beta/\gamma)\bar B}\,d\mu_0\le e^{\beta/\gamma}<\infty.
\]
Let $\kappa_G:=\kappa_{t,\mathrm{cts}}^{\mathrm G,\omega_A}(\cdot,\cdot\mid h_t)$. Then
\[
\frac{d\kappa_G}{d\mu_0}(c,a)=\frac{e^{(\beta/\gamma)\bar B(c,a)}}{Z_{\mathrm{ker}}^{\mathrm{cts}}(h_t)}.
\]
If $\kappa$ is not absolutely continuous with respect to $\mu_0$, then $\KL(\kappa\|\mu_0)=+\infty$, so it cannot minimize the action term. For $\kappa\ll\mu_0$, the Radon--Nikodym identity gives
\begin{align*}
\gamma\,\KL(\kappa\|\mu_0)-\beta\int \bar B\,d\kappa
&=\gamma\int \log\!\left(\frac{d\kappa/d\mu_0}{e^{(\beta/\gamma)\bar B}}\right)\,d\kappa\\
&=\gamma\int \log\!\left(\frac{d\kappa/d\mu_0}{d\kappa_G/d\mu_0}\right)\,d\kappa-\gamma\log Z_{\mathrm{ker}}^{\mathrm{cts}}(h_t)\\
&=\gamma\,\KL(\kappa\|\kappa_G)-\gamma\log Z_{\mathrm{ker}}^{\mathrm{cts}}(h_t).
\end{align*}
The right-hand side is uniquely minimized at $\kappa=\kappa_G$. Combining the belief-side and action-side minimizations yields the exact global minimizer pair.
\end{proof}

\begin{remark}[What the kernel lift changes]
\label{rem:kernel-functional-scope}
Definition~\ref{def:kernel-functional} absorbs the selector into the variational object itself: action is sampled directly from the exact minimizing class-action kernel $\kappa_t^{\mathrm G,\omega_A}$. Its price is that convergence now depends on the reference mass that $\lambda$ assigns to separating actions. In a finite action alphabet this is automatic for uniform $\lambda$; in a countably infinite action alphabet it is an extra hypothesis, not a free consequence of $\sup_a B_t^{\omega_A}(c,a;h_t)>0$.
\par
Proposition~\ref{prop:kernel-functional-minimizer-compact} gives a second exact route: for compact metric action spaces, the same Gibbs-kernel minimizer exists measure-theoretically, and a continuity argument can recover a positive $\lambda$-mass floor on separating neighborhoods even without finiteness. Section~\ref{sec:main} records this as a compact-action extension of the main kernel theorem.
\end{remark}

\begin{definition}[Meeting-point specialization]
\label{def:meeting-point-shorthand}
Write
\[
\mathcal J_t^{\omega_{\mathrm{behav}}}:=\mathcal J_t^{\omega_{\mathrm{syn}},\omega_{\mathrm{behav}}},
\qquad
\mathcal J_{t,\mathrm{raw}}^{\omega_{\mathrm{behav}}}:=\mathcal J_{t,\mathrm{raw}}^{\omega_{\mathrm{syn}},\omega_{\mathrm{behav}}},
\qquad
\mathcal J_{t,\mathrm{ker}}^{\omega_{\mathrm{behav}}}:=\mathcal J_{t,\mathrm{ker}}^{\omega_{\mathrm{syn}},\omega_{\mathrm{behav}}}
\]
for the specializations with the canonical belief observer $\omega_B=\omega_{\mathrm{syn}}$ (the quotient-free top of the refinement lattice) and action observer $\omega_A=\omega_{\mathrm{behav}}$ (the observer certified admissible on both sides in Theorem~\ref{thm:meeting-point}).
\end{definition}

The subsections that follow establish the formal admissibility predicates used by $\mathcal J_t^{\omega_B,\omega_A}$: the belief side is represented by refinement above $\omega_{\mathrm{behav}}$, and the action side by refinement below it. The Lean statements here record the predicates and the behavioral observer's reflexive admissibility; stronger uniqueness readings are motivation rather than additional formal conclusions.

\subsection{Belief admissibility: a half-lattice above $\omega_{\mathrm{behav}}$}
\label{sec:belief-admissibility}

The Lean-facing admissibility predicates are expressed directly in the observer-refinement preorder:
\[
\mathrm{BeliefAdmissible}(\omega)\;:\!\!\Longleftrightarrow\;\omega_{\mathrm{behav}}\le\omega,
\qquad
\mathrm{ActionAdmissible}(\omega)\;:\!\!\Longleftrightarrow\;\omega\le\omega_{\mathrm{behav}}.
\]
The propositions below are the formal wrappers for those two half-lattices.

\begin{proposition}[Belief-admissible observers refine behavioral view]
\label{prop:belief-invariance-above}
If $\omega_B$ is belief-admissible and two programs have the same $\omega_B$-view, then they have the same behavioral view.
\end{proposition}

\begin{proof}
\noindent\emph{Lean: \path{SemanticConvergence.prop_belief_invariance_above}.}
Belief-admissibility is exactly a factorization of the behavioral observer through $\omega_B$. Applying that factorization to the equality of $\omega_B$-views gives equality of behavioral views.
\end{proof}

\begin{proposition}[Failure of belief admissibility blocks the goal-dial proxy]
\label{prop:belief-illtyped-below}
If $\omega_B$ is action-admissible but not belief-admissible, then the formal goal-dialed convergence proxy for $\omega_B$ does not hold.
\end{proposition}

\begin{proof}
\noindent\emph{Lean: \path{SemanticConvergence.prop_belief_illtyped_below}.}
In Lean, \texttt{goalDialConverges} is the belief-admissibility predicate. The result is therefore exactly the supplied negation of belief-admissibility.
\end{proof}

Proposition~\ref{prop:belief-invariance-above} says that belief-admissible observers are fine enough to determine the behavioral view. Proposition~\ref{prop:belief-illtyped-below} records the corresponding formal failure predicate when belief admissibility is absent. The current Lean theorem does not construct an ill-typed class likelihood; it packages the admissibility boundary used by the route.

\subsection{Action admissibility: a half-lattice below $\omega_{\mathrm{behav}}$}
\label{sec:action-admissibility}

The symmetric action-side formalization is expressed through bounded kernel weights and the action-admissibility predicate.

\begin{proposition}[Gibbs-kernel action cap]
\label{prop:action-cap}
Let $\bar B(c,a)$ be a kernel score with $0\le \bar B(c,a)\le1$, let $\beta\ge0$ and $\gamma>0$, and let $\lambda$ be the reference action law. Then the Gibbs kernel weight is capped by the reference weight times the universal bounded-score factor:
\[
K_{\mathrm G}(c,a)
\le
K_{\mathrm{ref}}(c,a)\exp(\beta/\gamma),
\]
for every observer class $c$ and action $a$.
\end{proposition}

\begin{proof}
\noindent\emph{Lean: \path{SemanticConvergence.prop_action_cap}.}
This is the denominator-free upper envelope for a bounded Gibbs reweighting: because $\bar B(c,a)\le1$ and $\beta/\gamma\ge0$, the exponential score factor is at most $\exp(\beta/\gamma)$, so each Gibbs weight is bounded by the corresponding reference weight times that factor.
\end{proof}

\begin{corollary}[Behavioral twins share the behavioral observer fiber]
\label{cor:twins-frozen-ratio}
If two programs have the same behavioral view, then their behavioral observer fibers are the same set.
\end{corollary}

\begin{proof}
\noindent\emph{Lean: \path{SemanticConvergence.cor_twins_frozen_ratio}.}
The encoded behavioral fiber is defined by equality of behavioral views. If the two target views are equal, the two predicates are extensionally equal.
\end{proof}

The formal action-side cap is a bounded-score Gibbs envelope, while Corollary~\ref{cor:twins-frozen-ratio} records the observer-fiber equality for behavioral twins. Classical posterior-ratio freezing is an interpretation that would require an additional Bayes-class likelihood theorem; it is not the Lean statement used below.

\subsection{The meeting-point theorem}
\label{sec:meeting-point}

\begin{theorem}[Behavioral observer is admissible on both sides]
\label{thm:meeting-point}
The behavioral observer is both belief-admissible and action-admissible:
\[
\mathrm{BeliefAdmissible}(\omega_{\mathrm{behav}})
\qquad\text{and}\qquad
\mathrm{ActionAdmissible}(\omega_{\mathrm{behav}}).
\]
\end{theorem}

\begin{proof}
\noindent\emph{Lean: \path{SemanticConvergence.thm_meeting_point}.}
Both fields are reflexivity of the observer-refinement preorder at $\omega_{\mathrm{behav}}$.
\end{proof}

\begin{corollary}[Canonical two-observer choice]
\label{cor:canonical-pair}
If the framework uses the canonical observer choice, the action observer is $\omega_{\mathrm{behav}}$ and the belief observer may be lifted to $\omega_{\mathrm{syn}}$ for quotient-free program-level Bayes. Proposition~\ref{prop:belief-invariance-above} certifies that the lifted belief observer still determines the behavioral view; Theorem~\ref{thm:meeting-point} certifies that $\omega_{\mathrm{behav}}$ is admissible on both sides of the formal refinement interface.
\end{corollary}

The formal meeting-point block records the refinement predicates used by the verified package. Stronger uniqueness or posterior-ratio statements are mathematical motivations for the ontology rather than separate formal theorems.

\begin{remark}[Leibniz identification]
\label{rem:leibniz}
Two programs $p_1,p_2$ satisfy $p_1\sim_{\mathrm{behav}}p_2$ iff they are indiscernible under every interactive probing --- the interactive-learning analogue of Leibniz's identity of indiscernibles on $\Pcal$. In these terms, $[p^\star]=[p^\star]_{\mathrm{behav}}$ is the Leibniz invariant of $p^\star$, and Theorem~\ref{thm:meeting-point} identifies this invariant as a lattice fixed point rather than a pragmatic stipulation: the two admissibility ranges already force it before any external criterion of ``behavioral equivalence'' is imposed.
\end{remark}

\subsection{Goal-dialed targets}
\label{sec:goal-dialed}

The meeting-point theorem records a canonical admissibility witness at $\omega_{\mathrm{behav}}$; the family of observer fibers still has content as a \emph{dial} for target specification. The formal result recorded at this point is the admissibility statement needed to treat a dialed observer fiber as a well-formed target. The convergence payoff for verified packages is stated later as Corollary~\ref{cor:goal-dialed-payoff}.

\begin{proposition}[Goal-dialed admissibility]
\label{prop:goal-dialed}
If an observer $\omega$ is belief-admissible, i.e.\ $\omega_{\mathrm{behav}}\preceq\omega$, then its dialed target fiber $[p^\star]_\omega$ satisfies the formal goal-dialed target predicate.
\end{proposition}

\begin{proof}
\noindent\emph{Lean: \path{SemanticConvergence.prop_goal_dialed}.}
In Lean the goal-dialed target predicate is the belief-admissibility condition for the observer, so the result is exactly the supplied admissibility witness.
\end{proof}

Proposition~\ref{prop:goal-dialed} is deliberately only an admissibility statement. Actual posterior concentration for a dialed target is not inferred from the observer order alone; it is routed through the bundled package in Corollary~\ref{cor:goal-dialed-payoff}.

\subsection{How the rest of the paper reads under the functional}
\label{sec:functional-forward-plan}

The remainder of the paper identifies the canonical Bayes/Gibbs-plus-separation structure singled out by the observer analysis, realizes the exact global minimizers of $\mathcal J_t^{\omega_B,\omega_A}$ and its action-realizing lift $\mathcal J_{t,\mathrm{ker}}^{\omega_B,\omega_A}$, and then validates the convergence conclusion.

Section~\ref{sec:belief} identifies the \emph{belief-side minimizer} of $\mathcal J_t^{\omega_B,\omega_A}$ at the canonical belief observer $\omega_B=\omega_{\mathrm{syn}}$: the Bayes/Gibbs posterior under the universal prior. The belief-side story is the content of Lemma~\ref{lem:variational}, Lemma~\ref{lem:kl-necessity}, and the posterior-gap fields of Propositions~\ref{prop:two-observer-minimizer}--\ref{prop:kernel-functional-minimizer}.

Section~\ref{sec:semantic-action} identifies the \emph{canonical action-side interfaces} at the action observer $\omega_A=\omega_{\mathrm{behav}}$: the exact class-action Gibbs kernel $\kappa_t^{\mathrm G,\omega_{\mathrm{behav}}}$ of Proposition~\ref{prop:kernel-functional-minimizer}, the Gibbs class-targeting minimizer $\nu_t^{\mathrm G,\omega_{\mathrm{behav}}}$ of Proposition~\ref{prop:two-observer-minimizer}, and the finite-list support proxies used by the concrete semantic stack. The verified package carries the load-bearing separation and support assumptions used for convergence. Corollaries~\ref{cor:noise-transfer} and~\ref{cor:noise-left-invertible-history-independent} route the verified noise companions through the zero-out package and the main theorem, respectively.

Section~\ref{sec:main} establishes the verified \emph{convergence payoff}: the route-2 kernel-functional theorem proves posterior concentration from exact Bayes/Gibbs belief, exact kernel-functional minimization, the canonical predictive-kernel observer, and the Assumptions~\ref{ass:main-verified-package}. The selector, kernel, compact-action, maximin, goal-dialed, and surrogate variants are now corollaries of the same theorem; the finite-time residual/rate/noise results form a zero-out companion stack.

Section~\ref{sec:self-ref} replaces the uncomputable behavioral observer by finite-time policy proxies and records the eventual-isolation predicates that a self-referential construction must verify before the verified or zero-out convergence routes apply.

Section~\ref{sec:near-miss} reads the algorithmic-free-energy principle through the boundary proxy: it records the risk-minus-gain objective, the finite-list argmin selector, and the witness predicates needed to certify a boundary near-miss or observer-promotion failure. The section is a formal interface for coupling failures, not an additional convergence theorem.

\section{Belief at $\omega_{\mathrm{syn}}$: Universal Bayesian Belief}
\label{sec:belief}

This section borrows the Gibbs variational form from active inference and the universal-prior reference from Solomonoff and Hutter. On the belief-evidence component these together give an exact Bayes/Gibbs identity: under the universal prior the algorithmic free energy admits a single Gibbs decomposition whose unique minimizer is the Bayes/Gibbs posterior, and the KL divergence is the unique form forced within the natural convex class by convex duality. The action side and the broader free-energy principle are not claimed to be subsumed: the action-side completion under the natural information-gain rule does not in general converge to $[p^\star]$ (Theorem~\ref{thm:afe-near-miss}), and the convergence gap is closed by the kernel-functional refinement rather than by the natural rule itself.

Section~\ref{sec:functional} uses the canonical belief observer at the top of the refinement lattice, $\omega_B=\omega_{\mathrm{syn}}$ (Corollary~\ref{cor:canonical-pair}). Proposition~\ref{prop:belief-invariance-above} records the formal refinement fact that such a belief-admissible lift determines the behavioral view. Working at the quotient-free top gives a clean setting in which to identify the belief-side minimizer of $\mathcal J_t^{\omega_{\mathrm{syn}},\omega_A}$: Bayes/Gibbs on programs, under the universal interactive prior, with its algorithmic free energy as the belief term.

Definition~\ref{def:universal-prior} fixes the universal prior. Definition~\ref{def:afe} identifies the \emph{algorithmic free energy} $\Fcal_t$ with the belief term of $\mathcal J_t^{\omega_{\mathrm{syn}},\omega_A}$. Lemma~\ref{lem:variational} gives the exact Bayes/Gibbs variational identity and the minimizer propositions of Section~\ref{sec:functional} use its posterior-gap form to identify the exact Bayes/Gibbs posterior. Lemma~\ref{lem:kl-necessity} shows the KL-based form is forced within the natural convex class. Lemma~\ref{lem:merging} records the formal observer-fiber invariance of the resulting posterior.

The belief side alone fixes the posterior algebra; it does not yet close the gap from history compatibility down to $[p^\star]$. The action-side counterpart of this unification --- the information-gain action criterion that completes the algorithmic free energy principle --- is identified at the boundary of the formal action-admissible range; Section~\ref{sec:near-miss} shows that this completed natural rule does not in general converge to the interventional class. The action side that does converge is the class-against-complement rule of Sections~\ref{sec:semantic-action} and~\ref{sec:main}.

\subsection{Universal prior}

\begin{definition}[Universal interactive prior]
\label{def:universal-prior}
For $p\in\Pcal$, let $\ellp(p)=|p|$ be the self-delimiting code length. Define raw weights $\wt w(p)=2^{-\ellp(p)}$, normalizing constant $Z=\sum_{p\in\Pcal}\wt w(p)\le 1$, normalized prior
\[
w(p)=\frac{\wt w(p)}{Z},
\]
and universal interactive mixture
\[
\barxi(o_{1:t}\givena a_{1:t})
:=
\sum_{p\in\Pcal} w(p)\mu_p(o_{1:t}\givena a_{1:t}).
\]
\end{definition}

\begin{lemma}[Normalized prior mass]
\label{lem:prior-invariance}
For a countable prefix machine $U$ and program $p$, the normalized universal prior mass is the raw universal weight divided by the total prior mass:
\[
\operatorname{prior}_U(p)
=
\frac{\operatorname{weight}_U(p)}
{\sum_r \operatorname{weight}_U(r)}.
\]
\end{lemma}

\begin{proof}
\noindent\emph{Lean: \path{SemanticConvergence.lem_prior_invariance}.}
This is the definitional equation for \path{SemanticConvergence.def_universal_prior}.
\end{proof}

\begin{lemma}[Class-prior weight]
\label{lem:prior-necessity}
For any predicate $C$ on programs, the formal class-prior weight is the sum of normalized prior masses over the programs satisfying $C$:
\[
w(C)=\sum_p \mathbf 1_C(p)\,w(p).
\]
\end{lemma}

\begin{proof}
\noindent\emph{Lean: \path{SemanticConvergence.lem_prior_necessity}.}
This is the definitional equation for \path{SemanticConvergence.classPriorWeight}.
\end{proof}

\subsection{Algorithmic free energy as the belief term of $\mathcal J_t$ at $\omega_B=\omega_{\mathrm{syn}}$}

For $p\in\Pcal$ and history $h_t$, define the cumulative history-fit loss
\[
L_t(p;h_t):=-\log\mu_p(o_{1:t}\givena a_{1:t}),
\]
with the convention $L_t=+\infty$ if the likelihood vanishes. Consistent with the notation of Definition~\ref{def:two-observer-functional}, $L_t(p;h_t)=\ellp_t(p;h_t)$.

\begin{definition}[Algorithmic free energy]
\label{def:afe}
For a posterior $q\in\Delta(\Pcal)$ with $\E_q[L_t]<\infty$ and $\KL(q\|w)<\infty$, the \emph{algorithmic free energy} at history $h_t$ is
\[
\Fcal_t[q;h_t]:=\E_q[L_t(p;h_t)]+\KL(q\|w).
\]
By construction, $\Fcal_t[q;h_t]$ is the belief term of the two-observer functional $\mathcal J_t^{\omega_{\mathrm{syn}},\omega_A}[q,\nu;h_t]$ at the canonical belief observer $\omega_B=\omega_{\mathrm{syn}}$ (Definition~\ref{def:two-observer-functional}): it is the part that depends on $q$ alone and not on $\nu$.
\end{definition}

\begin{lemma}[Variational identity]
\label{lem:variational}
Assume $\barxi(o_{1:t}\givena a_{1:t})>0$, and define the Bayes/Gibbs posterior
\[
q_t^\star(p\mid h_t):=\frac{w(p)\mu_p(o_{1:t}\givena a_{1:t})}{\barxi(o_{1:t}\givena a_{1:t})}.
\]
Then for every admissible posterior $q$,
\[
\Fcal_t[q;h_t]
=
\KL\!\bigl(q\,\|\,q_t^\star(\cdot\mid h_t)\bigr)-\log\barxi(o_{1:t}\givena a_{1:t}),
\]
\end{lemma}

\begin{proof}
\noindent\emph{Lean: \path{SemanticConvergence.lem_variational}.}
By definition,
\[
\Fcal_t[q;h_t]=\sum_p q(p)\log\frac{q(p)}{w(p)\mu_p(o_{1:t}\givena a_{1:t})}
=\KL\!\bigl(q\,\|\,q_t^\star(\cdot\mid h_t)\bigr)-\log\barxi(o_{1:t}\givena a_{1:t}).
\]
\end{proof}

\begin{remark}[Belief-side minimizer of $\mathcal J_t^{\omega_{\mathrm{syn}},\omega_A}$]
\label{rem:afe-bayes-unification}
Because $\Fcal_t$ is the $q$-dependent part of $\mathcal J_t^{\omega_{\mathrm{syn}},\omega_A}[q,\nu;h_t]$, Lemma~\ref{lem:variational} supplies the Bayes/Gibbs evidence decomposition used by Proposition~\ref{prop:two-observer-minimizer} and Proposition~\ref{prop:kernel-functional-minimizer}. Those minimizer propositions identify the belief-side minimizer of the two-observer and kernel functionals at $\omega_B=\omega_{\mathrm{syn}}$ with the exact Bayes/Gibbs posterior $q_t^\star$, uniformly in the action-side choice. Put differently, fixing $\omega_B=\omega_{\mathrm{syn}}$ turns the $q$-minimization of the repaired Gibbs objectives into Bayes under the universal interactive prior; no approximation is involved and no action-side choice disturbs this identification.

The algorithmic free energy of Definition~\ref{def:afe} has not, to our knowledge, appeared in this exact form. The variational free energy of active inference takes a Gibbs form but uses a stipulated parametric prior that upper-bounds model evidence only approximately \citep{Friston2010,ParrFriston2019}; Solomonoff--Hutter Bayesian prediction uses the Kolmogorov-complexity-weighted universal prior but does not state a free-energy functional \citep{Solomonoff1964a,Solomonoff1964b,Hutter2005}. Coupling the two --- the Gibbs form with the universal prior as reference measure --- narrows the gap on the evidence-decomposition component. Under the universal prior the algorithmic free energy on this component is not a bound on the evidence; the variational identity gives an exact Bayes/Gibbs evidence decomposition, and the repaired Gibbs objectives built from it have the exact Bayes/Gibbs posterior as their belief-side minimizer (Lemma~\ref{lem:variational}; Propositions~\ref{prop:two-observer-minimizer}, \ref{prop:kernel-functional-minimizer}). The Gibbs variational form of active inference is borrowed; the stipulative quality of the prior, on this belief-evidence component, is removed. We do not claim that the broader free-energy principle is subsumed: action-side, homeostatic, and policy-as-inference readings of free energy are not the object of this paper, and the convergence gap left open by the natural information-gain action rule (Theorem~\ref{thm:afe-near-miss}) is what the kernel-functional refinement of Section~\ref{sec:main} closes.
\end{remark}

\begin{lemma}[KL is forced by the exact Gibbs variational formula]
\label{lem:kl-necessity}
Let $D$ be an extended divergence on the countable simplex satisfying the formal predicates \textup{ProperOnSimplex}, \textup{ConvexOnSimplex}, \textup{OffSimplexTop}, and \textup{SequentialLowerSemicontinuousOnSimplex}. Assume that for every bounded loss profile $L\in\ell^\infty(\Pcal)$,
\[
\inf_{q\in\Delta(\Pcal)}\bigl\{\E_q[L]+D(q\|w)\bigr\}
=
-\log\sum_{p\in\Pcal} w(p)e^{-L(p)}.
\]
Then $D(q\|w)=\KL(q\|w)$ for every $q\in\Delta(\Pcal)$.
\end{lemma}

\begin{proof}
\noindent\emph{Lean: \path{SemanticConvergence.lem_kl_necessity}.}
For bounded $f\in\ell^\infty(\Pcal)$ define $H(f):=\log\sum_p w(p)e^{f(p)}$. With $f=-L$, the hypothesis is equivalent to $H(f)=\sup_q\{\E_q[f]-D(q\|w)\}$. Fenchel--Moreau on the pair $(\ell^1,\ell^\infty)$ gives $D(q\|w)=\sup_f\{\E_q[f]-H(f)\}$.

If $\KL(q\|w)=\infty$, the upper bound $D(q\|w)\le \KL(q\|w)$ is immediate. Otherwise $q\ll w$. For $q_f(p):=w(p)e^{f(p)}/\sum_r w(r)e^{f(r)}$,
\[
\E_q[f]-H(f)=\KL(q\|w)-\KL(q\|q_f)\le \KL(q\|w),
\]
so $D(q\|w)\le\KL(q\|w)$. For the reverse, first suppose $q\not\ll w$, and let $A:=\{p\in\Pcal:q(p)>0,\ w(p)=0\}$. Then $q(A)>0$. For $f_n:=n\,\one_A$, one has $H(f_n)=\log\sum_p w(p)e^{f_n(p)}=0$ and $\E_q[f_n]=nq(A)\to\infty$, so $D(q\|w)=\infty=\KL(q\|w)$. It remains to treat the case $q\ll w$. Set $r(p):=q(p)/w(p)$ on $\{w(p)>0\}$ and $f_n(p):=\max\{-n,\min\{\log r(p),n\}\}$. Each $f_n$ is bounded. Dominated convergence gives $H(f_n)\to 0$. If $\KL(q\|w)<\infty$, $f_n(p)\to\log r(p)$ pointwise; the positive part is $q$-integrable by finiteness of $\KL$, and the bound $\sum_{r(p)<1}w(p)(-r\log r)\le 1/e$ controls the negative part, so $\E_q[f_n]\to\KL(q\|w)$. If $\KL=\infty$, monotone convergence on $f_n^+$ gives $\E_q[f_n^+]\to\infty$, and the same bound on the negative part shows $\sup_n\E_q[f_n^-]<\infty$, hence $\E_q[f_n]\to\infty$. Either way $D(q\|w)\ge\sup_n\{\E_q[f_n]-H(f_n)\}=\KL(q\|w)$.
\end{proof}

\begin{remark}
Lemma~\ref{lem:kl-necessity} singles out the KL-based free energy within the natural convex class.
\end{remark}

\subsection{Observer-fiber invariance of the Bayes/Gibbs posterior}

\begin{lemma}[Observer-fiber posterior invariance]
\label{lem:merging}
Let $\omega$ be an observer and let $p,q\in\Pcal$ have the same $\omega$-view. Then for every history $h_t$ and Bayes/Gibbs posterior,
\[
q_t^\star\bigl([p]_\omega\mid h_t\bigr)
=
q_t^\star\bigl([q]_\omega\mid h_t\bigr),
\]
where $[p]_\omega=\{r:\omega(r)=\omega(p)\}$.
\end{lemma}

\begin{proof}
\noindent\emph{Lean: \path{SemanticConvergence.lem_merging}.}
The two observer fibers are definitionally the same set when $\omega(p)=\omega(q)$. The observer-fiber posterior weight is the posterior mass of that set, so substituting the fiber equality gives the displayed identity.
\end{proof}

\begin{remark}[Belief alone does not close the observer gap]
\label{rem:belief-reaches-Rht}
Lemma~\ref{lem:merging} is the formal quotient-respect property needed by the later observer algebra: once a view is fixed, Bayes/Gibbs posterior mass is attached to the view's fiber, not to a particular representative. It is not by itself a convergence theorem. The strict inclusions $R_{h_t}(p^\star)\supsetneq R_\pi(p^\star)\supsetneq[p^\star]$ of Theorem~\ref{thm:strict-hierarchy} are not closed by belief algebra alone. Closing them requires an action-side rule and coupling assumptions that lift the realized observer toward $\omega_{\mathrm{behav}}$ --- the content of Sections~\ref{sec:semantic-action} and~\ref{sec:main}.
\end{remark}


\part{Coupling: When Agents Meet Environments Through Observers}
\label{part:coupling}

Part III is about coupling. An agent that is structurally well-formed by Part~\ref{part:agent}'s standards may still fail to converge in some environments. Convergence is a property of the triple $(\omega,\Acal,p^\star)$ --- an observer, an agent, and an environment --- jointly: an environment whose response to the agent's policy carries enough information through the observer's lens for the agent to extract that information at a positive rate. This part develops both the success cases (the canonical engineered couplings of Section~\ref{sec:semantic-action} and the main convergence theorem of Section~\ref{sec:main}) and the failure cases (the near-miss boundary example of Section~\ref{sec:near-miss} and the matching converses inside Section~\ref{sec:main}). The diagnostic vocabulary of Section~\ref{sec:ontology-diagnostic} is the citation surface for these results: each subsection of Part~\ref{part:coupling} corresponds to one diagnostic outcome.

\section{Action at $\omega_{\mathrm{behav}}$: Class-Against-Complement}
\label{sec:semantic-action}

This section constructs the \emph{class-against-complement} action rule at the canonical action observer $\omega_A=\omega_{\mathrm{behav}}$. The rule is qualitatively different from the information-gain rule that completes the algorithmic free energy principle (Section~\ref{sec:near-miss}): it scores observations not by their general reduction of uncertainty but by their power to separate a target class from its complement under the true environment law. The Principle of Section~\ref{sec:near-miss} names this distinction and Theorem~\ref{thm:separating-support-convergence} establishes that the class-against-complement rule, packaged into the verified coupling conditions, drives convergence to the interventional class --- which the natural information-gain rule, by Theorem~\ref{thm:afe-near-miss}, does not.

Theorem~\ref{thm:meeting-point} records that $\omega_{\mathrm{behav}}$ is admissible on both sides of the formal observer interface. What remains is to record action-side realizations and support checks matched to the separation objective at $\omega_A=\omega_{\mathrm{behav}}$. This section does so at two levels. The exact variational realization is the class-action Gibbs kernel of the kernel lift $\mathcal J_{t,\mathrm{ker}}^{\omega_{\mathrm{behav}}}$. The finite-list selector side is represented in Lean by a full-support proxy on an explicit action list, together with local semantic-separation witnesses. The relevant comparison at $\omega_A=\omega_{\mathrm{behav}}$ is class against aggregate complement: once the target is an $\omega_{\mathrm{behav}}$-class $c$, the relevant rival is all posterior mass outside $c$ pooled together.

For a semantic class $c\in\Ccal$, the notation $\mu_c$ fixed in Section~\ref{sec:setup} denotes the common action-conditioned law induced by any program in $c$.

\begin{definition}[Class-complement predictive law]
\label{def:class-complement}
Fix a history $h_t$, a class $c\in\Ccal$, and an action $a$. If $0<q_t^\star(c\mid h_t)<1$, define the class-complement predictive law $\Qcal_t^{-c}(\cdot\mid h_t,a)\in\Delta(\Ocal)$ by
\[
\Qcal_t^{-c}(o\mid h_t,a)
:=
\frac{\sum_{d\neq c} q_t^\star(d\mid h_t)\mu_d(o\mid h_t,a)}
{1-q_t^\star(c\mid h_t)}.
\]
\end{definition}

\begin{remark}[Identification with the action-observer complement]
\label{rem:class-complement-is-omega-behav}
Definition~\ref{def:class-complement} is the specialization of the $\omega_A$-class complement $\Qcal_t^{-c,\omega_A}$ of Section~\ref{sec:quotient-data} to $\omega_A=\omega_{\mathrm{behav}}$. At this observer, a class $c$ is read through its behavioral view, and the complement law pools posterior mass outside that view. The class score entering the action side of $\mathcal J_t^{\omega_B,\omega_{\mathrm{behav}}}$ is therefore computed from the pair $(\mu_c,\Qcal_t^{-c})$ of Definition~\ref{def:class-complement}.
\end{remark}

\begin{definition}[Semantic separation]
\label{def:semantic-separation}
The \emph{semantic separation} of action $a$ for class $c$ at $h_t$ is
\[
B_t(c,a;h_t)
:=
\begin{cases}
-\log\sum_{o\in\Ocal}\sqrt{\mu_c(o\mid h_t,a)\Qcal_t^{-c}(o\mid h_t,a)}, & 0<q_t^\star(c\mid h_t)<1,\\
0, & q_t^\star(c\mid h_t)\in\{0,1\}.
\end{cases}
\]
\end{definition}

\begin{remark}[Identification with the action-observer separation]
\label{rem:B-is-B-omega-behav}
$B_t(c,a;h_t)=B_t^{\omega_{\mathrm{behav}}}(c,a;h_t)$ in the sense of Definition~\ref{def:bhat-omega}: Remark~\ref{rem:class-complement-is-omega-behav} identifies the pair $(\mu_c,\Qcal_t^{-c})$ with the pair $(\mu_c^{q_t^\star},\Qcal_t^{-c,\omega_{\mathrm{behav}}})$ of Section~\ref{sec:quotient-data}, and the two definitions take $-\log$ of the same Hellinger affinity. The notational shorthand $B_t:=B_t^{\omega_{\mathrm{behav}}}$ fixed in Definition~\ref{def:bhat-omega} is in force throughout this section.
\end{remark}

\begin{definition}[Semantic gain]
\label{def:semantic-gain}
The \emph{semantic gain} of action $a$ for class $c$ at $h_t$ is the posterior mass of the target class times its semantic separation:
\[
S_t(c,a;h_t):=q_t^\star(c\mid h_t)\,B_t(c,a;h_t).
\]
\end{definition}

\begin{lemma}[Posterior-odds identity]
\label{lem:odds-identity}
For a target representative $p^\star$ and observer $\omega$, the observer-fiber posterior odds are exactly the class posterior odds of the corresponding observer fiber:
\[
\operatorname{Odds}_t^\omega(p^\star\mid h_t)
=
\frac{q_t^\star([p^\star]_\omega\mid h_t)}
{1-q_t^\star([p^\star]_\omega\mid h_t)}.
\]
\end{lemma}

\begin{proof}
\noindent\emph{Lean: \path{SemanticConvergence.lem_odds_identity}.}
The observer fiber of $p^\star$ is the class whose view is $\omega(p^\star)$. Both sides unfold to the same class-posterior odds expression.
\end{proof}

\begin{lemma}[Zero criterion]
\label{lem:zero-criterion}
For an action $a$, target representative $p^\star$, and observer $\omega$,
\[
a\text{ is semantically separating for }(h_t,\omega,p^\star)
\iff
0<B_t^\omega([p^\star]_\omega,a;h_t).
\]
\end{lemma}

\begin{proof}
\noindent\emph{Lean: \path{SemanticConvergence.lem_zero_criterion}.}
This is the definitional criterion used by the concrete semantic stack: a separating action is exactly an action whose observer-fiber semantic separation is positive.
\end{proof}

\begin{proposition}[Class-complement Bhattacharyya correspondence]
\label{prop:chernoff-correspondence}
For every action $a$, target representative $p^\star$, and observer $\omega$,
\[
B_t^\omega([p^\star]_\omega,a;h_t)
=
\operatorname{Bhat}\!\left(
\mu_{[p^\star]_\omega}(\cdot\mid h_t,a),
\Qcal_t^{-[p^\star]_\omega}(\cdot\mid h_t,a)
\right),
\]
where the right-hand side is the Bhattacharyya separation of the concrete class-complement predictive pair.
\end{proposition}

\begin{proof}
\noindent\emph{Lean: \path{SemanticConvergence.prop_chernoff_correspondence}.}
The semantic separation definition first constructs the class-complement pair and then applies the concrete Bhattacharyya-separation function to that pair. The displayed equality is the corresponding unfold.
\end{proof}

The correspondence settles the book-keeping issue needed by the verified convergence proof: the score entering $\mathcal J_t^{\omega_B,\omega_A}$ is exactly the bounded class-complement Bhattacharyya geometry formalized in Lean. The capped score $G_t(c;h_t)=\sup_a\min\{1,B_t(c,a;h_t)\}$ that enters $\mathcal J_t^{\omega_{\mathrm{behav}}}$ is the bounded version of this same formal separation.

\begin{definition}[Finite-list semantic action rule]
\label{def:semantic-rule}
For an explicit finite action list $A_0$, the Lean-facing \emph{semantic action rule} is the full-support action law on $A_0$.
\end{definition}

The class-score Gibbs distribution of Proposition~\ref{prop:two-observer-minimizer} is the exact variational minimizer on observer classes. Definition~\ref{def:semantic-rule} is the finite-list support proxy used by the concrete semantic stack to talk about which actions are available with nonzero mass.

\begin{definition}[Promotion-supporting action rule]
\label{def:promotion-supporting}
For a finite action list $A_0$, an action law $\kappa$ is \emph{promotion-supporting on $A_0$} if every listed action has nonzero mass:
\[
\forall a\in A_0,\qquad \kappa(a)>0.
\]
\end{definition}

\begin{proposition}[Finite-list semantic rule is promotion-supporting]
\label{prop:semantic-is-promotion-supporting}
The finite-list semantic action rule of Definition~\ref{def:semantic-rule} is promotion-supporting on its listed actions.
\end{proposition}

\begin{proof}
\noindent\emph{Lean: \path{SemanticConvergence.prop_semantic_is_promotion_supporting}.}
By definition, the semantic-rule proxy is the full-support action law on the finite list. Full support gives nonzero mass to each listed action.
\end{proof}

\begin{definition}[Kernel semantic action rule]
\label{def:kernel-semantic-rule}
Given a reference action law $\lambda$, the Lean-facing \emph{kernel semantic action rule} is the full-support action law on the support of $\lambda$.
\end{definition}

\begin{proposition}[Kernel rule is promotion-supporting on reference support]
\label{prop:kernel-promotion-support}
The kernel semantic action rule of Definition~\ref{def:kernel-semantic-rule} is promotion-supporting on the support of its reference action law.
\end{proposition}

\begin{proof}
\noindent\emph{Lean: \path{SemanticConvergence.prop_kernel_promotion_support}.}
By definition, the kernel semantic-rule proxy is the full-support action law on the reference law's support. Full support gives nonzero mass to each action in that support.
\end{proof}

\begin{proposition}[Compact-action Gibbs kernel product-mass floor]
\label{prop:kernel-promotion-support-compact}
In the compact-action setup, let $S$ be a compact product reference measure, let $B$ be a bounded kernel score with $0\le B\le1$, let $\beta\ge0$ and $\gamma>0$, and let $E\subseteq\Acal$ be measurable. Then the compact Gibbs measure satisfies
\[
\exp(-\beta/\gamma)\,
\bigl(\lambda_C(\{c\})\lambda_A(E)\bigr)
\le
\kappa_{\mathrm G}(\{c\}\times E).
\]
\end{proposition}

\begin{proof}
\noindent\emph{Lean: \path{SemanticConvergence.prop_kernel_promotion_support_compact}.}
The compact-kernel theorem applies the bounded Gibbs reweighting floor to the product set $\{c\}\times E$ under the product reference measure.
\end{proof}

\begin{definition}[Decodable observation channel]
\label{def:decodable-channel}
A channel $K:\Ocal\to\Delta(\widetilde\Ocal)$ is \emph{decodable} in the formal support sense if there exists a map $D:\widetilde\Ocal\to\Ocal$ such that every emitted observation with nonzero channel mass decodes back to its source:
\[
K(\widetilde o\mid o)\neq 0
\quad\Longrightarrow\quad
D(\widetilde o)=o.
\]
For a corrupted history $\widetilde h_t=(a_1,\widetilde o_1,\dots,a_t,\widetilde o_t)$, write
\[
D(\widetilde h_t):=(a_1,D(\widetilde o_1),\dots,a_t,D(\widetilde o_t)).
\]
\end{definition}

\begin{proposition}[Noisy observer-fiber invariance]
\label{prop:noise-immunity}
Let $K:\Ocal\to\Delta(\widetilde\Ocal)$ be any observation channel. If two programs have the same observer view, $\omega(p)=\omega(q)$, then the channelized observer-fiber class-complement pair, and hence the noisy Bhattacharyya separation computed from it, is the same for $p$ and $q$:
\[
B_t^K(p,a;h_t,\omega)=B_t^K(q,a;h_t,\omega)
\qquad
\text{for every }h_t\text{ and }a.
\]
\end{proposition}

\begin{proof}
\noindent\emph{Lean: \path{SemanticConvergence.prop_noise_immunity}.}
The clean observer-fiber complement law is invariant under replacing $p$ by a program $q$ with the same observer view. Applying the same channel image map to both sides preserves equality, and applying the Bhattacharyya score to the equal pairs gives the displayed noisy-separation identity.
\end{proof}

\begin{definition}[Support-left-invertible observation channel]
\label{def:left-invertible-channel}
A channel $K:\Ocal\to\Delta(\widetilde\Ocal)$ is \emph{support-left-invertible} if some decoder $D:\widetilde\Ocal\to\Ocal$ identifies the source of every possible emitted observation:
\[
K(\widetilde o\mid o)\neq 0
\quad\Longrightarrow\quad
D(\widetilde o)=o.
\]
Thus the formal notion is a support-level left inverse, not the weaker analytic condition that a finite channel matrix have full column rank.
\end{definition}

\begin{proposition}[Support-left-invertibility unwraps a decoder]
\label{prop:noise-left-invertible}
If $K:\Ocal\to\Delta(\widetilde\Ocal)$ is support-left-invertible, then there exists a decoder $D:\widetilde\Ocal\to\Ocal$ such that
\[
K(\widetilde o\mid o)\neq 0
\quad\Longrightarrow\quad
D(\widetilde o)=o.
\]
\end{proposition}

\begin{proof}
\noindent\emph{Lean: \path{SemanticConvergence.prop_noise_left_invertible}.}
This is the elimination form of Definition~\ref{def:left-invertible-channel}: the Lean proof destructs the support-left-invertibility witness and returns its decoder.
\end{proof}

\begin{proposition}[Decodable channels are support-left-invertible]
\label{prop:noise-decoding}
Every decodable channel in the sense of Definition~\ref{def:decodable-channel} is support-left-invertible in the sense of Definition~\ref{def:left-invertible-channel}.
\end{proposition}

\begin{proof}
\noindent\emph{Lean: \path{SemanticConvergence.prop_noise_decoding}.}
The two formal predicates carry the same support decoder. The proof reuses the decoder witnessing decodability as the support-left-inverse.
\end{proof}

\begin{corollary}[Decodable-channel transfer for zero-out rate certificates]
\label{cor:noise-transfer}
Fix a decodable observation channel $K:\Ocal\to\Delta(\widetilde\Ocal)$, a finite horizon $T$, a concrete zero-out rate package $\mathfrak Z_T$ in the sense of Definition~\ref{def:zero-out-rate-package} for source program $p$, and a program $q$ in the same observer fiber as $p$. Then $K$ is support-left-invertible, and for every decay parameter $\delta$ the finite-horizon posterior-share lower bound certified by $\mathfrak Z_T$ holds for $q$:
\[
\left(1+\rho_\delta^N R_0(q)\right)^{-1}
\le
Q_N(q;\xi)
\qquad
\text{for all }N\le T
\]
almost surely under the constructed countable trajectory law.
\end{corollary}

\begin{proof}
\noindent\emph{Lean: \path{SemanticConvergence.zeroOutRatePackage_decodableNoiseTransfer}.}
The Lean theorem unwraps decodability into support-left-invertibility and then applies the same-view finite-time posterior-share theorem \path{SemanticConvergence.zeroOutRatePackage_sameViewPosteriorShareFiniteTime} to the supplied package $\mathfrak Z_T$.
\end{proof}

\begin{corollary}[Support-left-invertible noisy wrapper]
\label{cor:noise-left-invertible-history-independent}
Let $K:\Ocal\to\Delta(\widetilde\Ocal)$ be support-left-invertible. Suppose the observed learner--environment system under consideration is represented by an verified semantic-learning package $\mathfrak K=(\mathcal I,\mathcal J,\Scal)$ satisfying Assumptions~\ref{ass:main-verified-package}. Then the channel witness holds and the posterior mass of the semantic target converges to one under the constructed Bayes/Gibbs trajectory law:
\[
\mathsf{SLI}(K)
\quad\text{and}\quad
q_t^\star(\Scal)\longrightarrow 1
\qquad\text{almost surely.}
\]
No convergence theorem is asserted from support-left-invertibility alone; the load-bearing convergence hypotheses are exactly the verified package assumptions.
\end{corollary}

\begin{proof}
\noindent\emph{Lean: \path{SemanticConvergence.Ontology.Coupling.h10_support_left_invertible_noisy_transfer}.}
The Lean theorem pairs the supplied support-left-invertibility witness with Theorem~\ref{thm:separating-support-convergence}, instantiated through the same verified package. This discharges the noisy companion claim from the main theorem while making explicit that the channel property itself does not construct the verified package.
\end{proof}

\begin{remark}[Against entropy-bonus curiosity]
\label{rem:against-entropy-bonus}
The formal noise stack records four bookkeeping facts. Proposition~\ref{prop:noise-immunity} proves observer-fiber invariance after an arbitrary channel; Propositions~\ref{prop:noise-left-invertible} and~\ref{prop:noise-decoding} record the support-decoder bookkeeping; Corollary~\ref{cor:noise-transfer} routes finite-time noisy certificates through the zero-out package; and Corollary~\ref{cor:noise-left-invertible-history-independent} routes asymptotic noisy learning through the verified package and main theorem. It does not claim that arbitrary high-entropy observations create semantic evidence, nor that support-left-invertibility by itself proves convergence. This is the relevant contrast with entropy-bonus curiosity: the class-against-complement rule is tied to verified observer-fiber evidence, not to raw observation entropy (cf.\ the MI-variant $\mathcal J_{t,\mathrm{raw}}^{\omega_B,\omega_A,\mathrm{MI}}$ after Definition~\ref{def:raw-two-observer-functional} and the discussion of expected free energy in Remark~\ref{rem:efe-unification}).
\end{remark}

\section{Sufficient Conditions for Semantic Convergence}
\label{sec:main}

Section~\ref{sec:functional} fixed the canonical action observer at $\omega_A=\omega_{\mathrm{behav}}$ (Corollary~\ref{cor:canonical-pair}) and Section~\ref{sec:belief} fixed the canonical belief observer at $\omega_B=\omega_{\mathrm{syn}}$; this section states the verified convergence interface for what happens once the learner acts there. The main theorem uses the exact kernel-functional action rule and the canonical predictive-kernel observer. Its natural object is not the target program by itself but a semantic-learning package $\mathfrak K=(\mathcal I,\mathcal J,\Scal)$: the realized learning instance, the functional minimized by the action rule, and the semantic set whose posterior mass is supposed to converge. Assumptions~\ref{ass:main-verified-package} is the single load-bearing package for that theorem: it lists the substrate, Bayes/Gibbs belief, exact minimization, normalizer, prefix-Hellinger, and class-predictive separation hypotheses in one place.

The rest of the section separates the two routes explicitly. Convergence variants are read through Assumptions~\ref{ass:main-verified-package}; finite-horizon residual/rate/noise claims belong to the separate concrete zero-out stack formalized in the Lean development.

\begin{definition}[Semantic separation condition]
\label{def:sep-condition}
For a finite action list $A_0\subseteq\Acal$, a history $h_t$, an observer $\omega$, and a target representative $p^\star$, the concrete \emph{semantic separation condition} holds when some listed action is semantically separating:
\[
\exists a\in A_0,\qquad
a\text{ is semantically separating for }(h_t,\omega,p^\star).
\]
\end{definition}

\begin{definition}[Uniform history-wise semantic separation]
\label{def:uniform-history-independent-separation}
For fixed $(\omega,p^\star,A_0)$, the concrete \emph{uniform history-wise semantic separation condition} holds when Definition~\ref{def:sep-condition} holds at every full history $h_t$.
\end{definition}

\begin{proposition}[Uniform history-wise separation specializes to each history]
\label{prop:uniform-history-independent-implies-semantic}
If $(\omega,p^\star,A_0)$ satisfies the uniform history-wise semantic separation condition, then it satisfies the semantic separation condition at each particular history $h_t$.
\end{proposition}

\begin{proof}
\noindent\emph{Lean: \path{SemanticConvergence.prop_uniform_history_independent_implies_semantic}.}
The uniform hypothesis is the function that assigns a semantic-separation witness to every history; evaluating it at the chosen $h_t$ gives the local witness.
\end{proof}

\begin{corollary}[Semantic separation witness transfer]
\label{cor:kl-implies-semantic-separation}
If the concrete semantic separation condition holds at $(h_t,\omega,p^\star,A_0)$, then the same concrete semantic separation condition holds at $(h_t,\omega,p^\star,A_0)$.
\end{corollary}

\begin{proof}
\noindent\emph{Lean: \path{SemanticConvergence.cor_kl_implies_semantic_separation}.}
Unpack the witness action and repack it unchanged.
\end{proof}

\begin{corollary}[Event-witness placeholder is the same semantic witness]
\label{cor:event-witness-implies-semantic-separation}
If the concrete semantic separation condition holds at $(h_t,\omega,p^\star,A_0)$, then the same concrete semantic separation condition holds at $(h_t,\omega,p^\star,A_0)$.
\end{corollary}

\begin{proof}
\noindent\emph{Lean: \path{SemanticConvergence.cor_event_witness_implies_semantic_separation}.}
Unpack the witness action and repack it unchanged.
\end{proof}

\begin{remark}[Analytic sufficient conditions are not part of the verified route]
\label{rem:no-pairwise-bhat-corollary}
Earlier formulations used KL and event-witness inequalities as sufficient conditions for positive Bhattacharyya separation. The route used here does not invoke those analytic inequalities directly. Its local semantic input is exactly Definition~\ref{def:sep-condition}, and the global convergence burden is carried by the bundled class-predictive separation field in Assumptions~\ref{ass:main-verified-package}.
\end{remark}

\begin{definition}[Kernel semantic separation condition relative to $\lambda$]
\label{def:kernel-sep-condition}
For a finite action list $A_0$, a history $h_t$, an observer $\omega$, a target representative $p^\star$, and an action law $\kappa$, the concrete \emph{kernel semantic separation condition} holds when some listed semantically separating action has nonzero $\kappa$-mass:
\[
\exists a\in A_0,\qquad \kappa(a)>0
\quad\text{and}\quad
a\text{ is semantically separating for }(h_t,\omega,p^\star).
\]
\end{definition}

\begin{proposition}[Finite-action kernel support from a separating action]
\label{prop:finite-action-kernel-separation}
If a listed action $a\in A_0$ is semantically separating for $(h_t,\omega,p^\star)$, then the full-support action law on $A_0$ satisfies the concrete kernel semantic separation condition of Definition~\ref{def:kernel-sep-condition}.
\end{proposition}

\begin{proof}
\noindent\emph{Lean: \path{SemanticConvergence.prop_finite_action_kernel_separation}.}
The full-support law on $A_0$ assigns nonzero mass to every listed action. Pair that nonzero mass with the supplied separating-action proof.
\end{proof}

\begin{proposition}[Compact-action local support floor]
\label{prop:compact-action-kernel-separation}
Let $X$ be a compact metric action space with a full-support Borel probability measure $\lambda$, let $\operatorname{score}:X\to\mathbb R$ be continuous, and let $\eta,r>0$. If $a_0$ has score at least $\eta$ and all points within distance $r$ of $a_0$ lose at most $\eta/2$ of score, then
\[
0<\eta/2,\qquad
m_\lambda(r)>0,\qquad
\{a:\eta/2\le \operatorname{score}(a)\}\text{ is measurable,}
\]
and the measure of the high-score set is at least the compact full-support ball-mass floor $m_\lambda(r)$.
\end{proposition}

\begin{proof}
\noindent\emph{Lean: \path{SemanticConvergence.prop_compact_action_kernel_separation}.}
Full support gives positive mass to every nonempty ball, and compactness gives the uniform ball-mass floor. Continuity makes the high-score set measurable. The local score-drop hypothesis puts the ball around $a_0$ inside the high-score set, so monotonicity of measure gives the displayed lower bound.
\end{proof}

\begin{lemma}[Conditional Borel--Cantelli]
\label{lem:conditional-bc}
Let $(\mathcal{F}_t)_{t\ge 0}$ be a filtration and $E_{t+1}\in\mathcal{F}_{t+1}$ for every $t\ge 0$. If $\sum_t\Prob(E_{t+1}\mid\mathcal{F}_t)=\infty$ almost surely, then $E_{t+1}$ i.o.\ almost surely.
\end{lemma}

\begin{proof}
\noindent\emph{Lean: \path{SemanticConvergence.lem_conditional_bc}.}
See \citet[Chapter~12]{Williams1991}.
\end{proof}

\begin{definition}[Hellinger convergence spine]
\label{def:hellinger-spine}
Fix a target program $p^\star\in\Pcal$ and write $c^\star:=[p^\star]$. Let $\Omega$ be an infinite realized-trajectory space with law $\Prob$ and filtration $(\mathcal G_t)_{t\ge0}$. A triple of real-valued adapted processes
\[
R_t:\Omega\to[0,\infty),
\qquad
S_t:\Omega\to\mathbb R,
\qquad
M_t:\Omega\to\mathbb R
\]
is a \emph{Hellinger convergence spine for $c^\star$} if:
\begin{itemize}
\item[(H1)] $R_t$ is the posterior residual odds
\[
R_t=\frac{1-q_t^\star(c^\star\mid H_t)}{q_t^\star(c^\star\mid H_t)}
\]
under the Bayes/Gibbs posterior, with the harmless zero convention when the numerator vanishes;
\item[(H2)] $M_t$ is an $L^1$-bounded martingale with respect to $(\mathcal G_t)$;
\item[(H3)] $S_t\to+\infty$ almost surely;
\item[(H4)] for every $t$ almost surely,
\[
\sqrt{R_t}=M_t\exp(-S_t).
\]
\end{itemize}
In the intended Bhattacharyya instance,
\[
S_t=\sum_{i=0}^{t-1}B_i(c^\star,A_{i+1};H_i),
\qquad
M_t=\exp(S_t)\sqrt{R_t}.
\]
\end{definition}

\begin{lemma}[Hellinger spine drives posterior odds to zero]
\label{lem:contraction}
Fix a target program $p^\star\in\Pcal$ and write $c^\star:=[p^\star]$. If the Bayes/Gibbs residual odds process admits a Hellinger convergence spine in the sense of Definition~\ref{def:hellinger-spine}, then
\[
q_t^\star(c^\star\mid H_t)\to 1
\qquad
\text{almost surely.}
\]
\end{lemma}

\begin{proof}
\noindent\emph{Lean: \path{SemanticConvergence.lem_contraction}.}
By martingale convergence for $L^1$-bounded martingales, $M_t$ converges almost surely to a finite real limit. On the same almost-sure event, (H3) gives $S_t\to+\infty$, so $\exp(-S_t)\to0$. The square-root identity (H4) therefore gives
\[
\sqrt{R_t}=M_t\exp(-S_t)\longrightarrow 0.
\]
Since $R_t\ge0$, this implies $R_t\to0$. Finally,
\[
q_t^\star(c^\star\mid H_t)=\frac{1}{1+R_t}
\]
whenever the residual-odds expression is finite, with the zero-numerator convention giving the same limiting conclusion. Hence $q_t^\star(c^\star\mid H_t)\to1$ almost surely.
\end{proof}

\begin{proposition}[Full support implies promotion support on a listed action set]
\label{prop:full-support-behavioral}
If an action law $\kappa$ has support on every action in a finite list $A_0$, then $\kappa$ is promotion-supporting on $A_0$.
\end{proposition}

\begin{proof}
\noindent\emph{Lean: \path{SemanticConvergence.prop_full_support_behavioral}.}
Promotion support is the same pointwise nonzero-mass condition restricted to $A_0$, so it is immediate from full support on the list.
\end{proof}

\begin{definition}[Semantic-learning package]
\label{def:semantic-learning-package}
Fix a target program $p^\star\in\Pcal$. A \emph{semantic-learning package} for $p^\star$ is a triple
\[
\mathfrak K=(\mathcal I,\mathcal J,\Scal)
\]
with the following components. The instance $\mathcal I$ contains the environment family, universal prior, true program, policy, reference action laws, inverse-temperature parameters, observers, posterior process, and infinite trajectory law. The functional $\mathcal J$ is the history-indexed action objective actually minimized by the learner; in the route it is
\[
\mathcal J
=\{\mathcal J_{t,\mathrm{ker}}^{\omega_B,\omega_A}(\,\cdot\,;h_t):t\ge0,\ h_t\in\Hcal_t\}.
\]
The target map $\Scal$ assigns to $p^\star$ the semantic set on which posterior mass is intended to concentrate; in the verified route this is the predictive-kernel fiber
\[
\Scal(p^\star)
:=
[p^\star]_{\omega_{\mathrm{pred}}}
:=
\{p\in\Pcal:\omega_{\mathrm{pred}}(p)=\omega_{\mathrm{pred}}(p^\star)\}.
\]
Thus $\mathcal I$ says what world and learner are being run, $\mathcal J$ says what is optimized, and $\Scal(p^\star)$ says what counts as semantic success.
\noindent\emph{Lean: \path{SemanticConvergence.Ontology.Coupling.def_semantic_learning_package} and \path{SemanticConvergence.Ontology.Coupling.predictiveKernelSemanticLearningPackage}.}
\end{definition}

\begin{assumptions}[Verified semantic-learning package]
\label{ass:main-verified-package}
Fix a semantic-learning package $\mathfrak K=(\mathcal I,\mathcal J,\Scal)$ for a target environment program $p^\star=p_{\mathrm{env}}$. Unpack $\mathcal I$ as a countable computable environment family, a policy $\pi$, inverse-temperature parameters $\beta,\gamma\in\mathbb R$, a belief observer $\omega_B$, an action observer $\omega_A$, the Bayes/Gibbs posterior process, and the induced infinite trajectory law. Let $\omega_{\mathrm{pred}}$ be the canonical predictive-kernel observer, whose view of a program is its full action-conditioned predictive kernel, and write
\[
c^\star:=\Scal(p^\star)=[p^\star]_{\omega_{\mathrm{pred}}}.
\]
The corresponding Lean package is \path{SemanticConvergence.Ontology.Coupling.ass_main_verified_package}. The package is verified when:
\begin{itemize}
\item[(M1)] \emph{Instance substrate.} The program, action, and observation spaces inside $\mathcal I$ are the countable/finite discrete objects of the formal setup, with the universal interactive prior and induced infinite Bayes/Gibbs trajectory law well defined.
\item[(M2)] \emph{Bayes/Gibbs belief.} At every finite history, belief is the exact Bayes/Gibbs posterior, equivalently the belief-side minimizer certified by Lemma~\ref{lem:variational}.
\item[(M3)] \emph{Exact kernel-functional action.} The functional component is $\mathcal J=\{\mathcal J_{t,\mathrm{ker}}^{\omega_B,\omega_A}\}_{t,h_t}$. For every finite history $h_t$, the belief weights, reference action law $\lambda_{h_t}$, class-action score $\bar B_{h_t}$, and class-action kernel $\kappa_{h_t}$ satisfy the admissibility and bounded-score hypotheses of Proposition~\ref{prop:kernel-functional-minimizer}; $\kappa_{h_t}$ attains the exact Gibbs value of $\mathcal J_{t,\mathrm{ker}}^{\omega_B,\omega_A}$; and the deployed policy is its action marginal:
\[
\pi_{t+1}(a\mid h_t)=\sum_c \kappa_{h_t}(c,a).
\]
\item[(M4)] \emph{Canonical predictive semantics.} The target component is $\Scal(p^\star)=[p^\star]_{\omega_{\mathrm{pred}}}$, the $\omega_{\mathrm{pred}}$-fiber of the true environment. Thus semantic calibration is not an extra hypothesis: equality of $\omega_{\mathrm{pred}}$-views is equality of predictive kernels.
\item[(M5)] \emph{Target-fiber normalizers.} For every finite history $h_t$, the posterior weight of the target predictive fiber is neither zero nor infinite:
\[
0<
q_t^\star(c^\star\mid h_t)
\qquad\text{and}\qquad
q_t^\star(c^\star\mid h_t)<\infty
\]
in the extended-real normalizer sense used by the Lean development.
\item[(M6)] \emph{True-law Hellinger prefix obligations.} The true-law Hellinger envelope is integrable at every finite prefix, and the local one-step Bayes/Hellinger identities, posterior-support witnesses, complement-support witnesses, predictive-support witnesses, and positive-affinity witnesses hold at every prefix. This is the manuscript form of \path{CountablePrefixMachine.HasTrueLawHellingerPrefixKernelObligations}.
\item[(M7)] \emph{Class-predictive semantic separation on reference support.} There is a constant $\kappa>0$ such that for every finite history $h_t$ and every action $a$ with positive reference mass $\lambda_{h_t}(a)>0$, the class-predictive Bhattacharyya score of the target predictive fiber against its posterior complement is at least $\kappa$:
\[
B_t^{\omega_{\mathrm{pred}}}(c^\star,a;h_t)\ge \kappa.
\]
\end{itemize}
\end{assumptions}

\begin{remark}[Reading the assumptions by layer]
\label{rem:main-by-layer}
In the language of Section~\ref{sec:ontology}, the package $\mathfrak K=(\mathcal I,\mathcal J,\Scal)$ records exactly an (observer, agent, environment) triple: $\mathcal I$ encodes the agent's prior and belief instance together with the environment, $\mathcal J$ encodes the agent's belief-and-action functional, and $\Scal(p^\star)$ is the packaged predictive semantic target for the environment. Assumptions \textup{(M1)--(M5)} are agent-side admissibility (well-formed prior, exact Bayes/Gibbs update, exact kernel-functional minimization, target-fiber normalizers); they hold or fail based on the agent alone. Assumptions \textup{(M6)} and \textup{(M7)} are coupling: the true-law Hellinger prefix-kernel obligations and the class-predictive Bhattacharyya separation $\kappa$ both depend jointly on observer, agent, and environment. Theorem~\ref{thm:separating-support-convergence} is therefore a coupling-layer theorem: it says that an admissible agent which is $\kappa$-coupled to the environment through the canonical predictive-kernel observer reaches $\Scal(p^\star)$ almost surely.
\end{remark}

\begin{center}
\fbox{\parbox{0.93\textwidth}{\centering
\textbf{Main result.} For an admissible agent $\mathfrak K=(\mathcal I,\mathcal J,\Scal)$ that is $\kappa$-coupled to the environment through the canonical predictive-kernel observer (Assumptions~\ref{ass:main-verified-package}),
\[
q_t^\star\bigl(\Scal(p^\star)\,\big|\,H_t\bigr)\longrightarrow 1
\qquad
\text{almost surely.}
\]
}}
\end{center}

\begin{theorem}[Verified kernel-functional semantic convergence]
\label{thm:separating-support-convergence}
Let $\mathfrak K=(\mathcal I,\mathcal J,\Scal)$ be an verified semantic-learning package in the sense of Assumptions~\ref{ass:main-verified-package}. Then
\[
q_t^\star(\Scal(p^\star)\mid H_t)\longrightarrow 1
\qquad
\text{almost surely.}
\]
Equivalently, the Bayes/Gibbs posterior of the instance $\mathcal I$, driven by minimizers of the functional $\mathcal J$, concentrates on the semantic target set $\Scal(p^\star)$ represented by the true predictive kernel.
\end{theorem}

\begin{proof}
\noindent\emph{Lean: \path{SemanticConvergence.Ontology.Coupling.h10_verified_semantic_learning_package_converges}.}
Assumptions \textup{(M1)--(M2)} supply the exact Bayes/Gibbs instance $\mathcal I$. Assumption \textup{(M3)} identifies the minimized functional $\mathcal J$ with the exact kernel-functional objective and supplies the minimizer whose action marginal is the deployed policy. Assumption \textup{(M4)} identifies the semantic target $\Scal(p^\star)$ with the canonical predictive-kernel fiber; in Lean this is \path{countablePredictiveKernelObserver}, for which predictive extensionality and truth-in-fiber are proved internally. Assumption \textup{(M5)} supplies the nonzero/non-top target-fiber normalizers, \textup{(M6)} supplies the true-law Hellinger prefix-kernel obligations and integrability needed to build the martingale spine, and \textup{(M7)} supplies the class-predictive separation on the reference support. The route-2 bridge converts \textup{(M7)} from a class-predictive statement into the true-law Hellinger affinity ceiling without a separate semantic-coherence assumption. The cited Lean theorem composes these ingredients and returns exactly the displayed almost-sure convergence.
\end{proof}

\begin{theorem}[Infinite-stream semantic-affinity bridge into the Hellinger spine]
\label{thm:infinite-affinity-hellinger-bridge}
Let $\mathfrak K=(\mathcal I,\mathcal J,\Scal)$ be an verified semantic-learning package in the sense of Assumptions~\ref{ass:main-verified-package}. Then the infinite Bayes/Gibbs trajectory law generated by $\mathcal I$ satisfies the Hellinger-spine bridge for the canonical predictive-kernel observer, and hence
\[
q_t^\star(\Scal(p^\star)\mid H_t)\longrightarrow 1
\qquad
\text{almost surely.}
\]
Equivalently, the martingale and affinity-ceiling ingredients that used to be separate bridge assumptions are exactly the \textup{(M6)}--\textup{(M7)} legs of the verified package.
\end{theorem}

\begin{proof}
\noindent\emph{Lean: \path{SemanticConvergence.Ontology.Coupling.h10_infinite_affinity_hellinger_bridge}.}
This is the main theorem read at the bridge layer. Assumption \textup{(M6)} supplies the true-law Hellinger prefix obligations and integrability; Assumption \textup{(M7)} supplies the class-predictive semantic separation on reference support; the predictive-kernel observer calibration proved in Lean turns that semantic separation into the true-law affinity ceiling. The cited Lean theorem is a corollary of Theorem~\ref{thm:separating-support-convergence}.
\end{proof}

\begin{remark}[Support-route boundary]
\label{rem:h6-support-boundary}
The original manuscript's support theorem used a weaker stochastic hypothesis:
divergent cumulative policy mass on a set of separating actions. The main route uses the exact kernel-functional minimizer together with the
canonical predictive-kernel observer and the assumptions above.
Downstream support/rate specializations should therefore be read either through
Assumptions~\ref{ass:main-verified-package} when they are part of the verified
main route, or as belonging to the separate zero-out special-case route formalized in the Lean development.
\end{remark}

\begin{definition}[Zero-out rate package]
\label{def:zero-out-rate-package}
Fix a concrete prefix machine $U$, a probabilistic concrete policy $\pi$, a true environment program $p_{\mathrm{env}}$, a source program $p$, an observer $\omega$, and a finite horizon $T$. Let $\widehat U$, $\widehat\pi$, $\widehat p_{\mathrm{env}}$, $\widehat p$, and $\widehat\omega$ denote the corresponding countable lifts used by the Lean development. A \emph{zero-out rate package} $\mathfrak Z_T(U,\pi,p_{\mathrm{env}},p,\omega)$ consists of:
\begin{itemize}
\item[(Z1)] concrete bookkeeping hypotheses: the program code list, policy supports, and semantic-kernel supports have no duplicates, and the concrete semantic kernels are probabilistic;
\item[(Z2)] realized prefix residual updates, i.e. the Lean predicate \path{HasRealizedPrefixResidualUpdates} holds for $(U,\pi,p_{\mathrm{env}},p,\omega,T)$;
\item[(Z3)] the initial residual observer-fiber odds of $\widehat p$ under $(\widehat U,\widehat\pi,\widehat\omega)$ are finite, so the posterior-share lower bound has a non-top denominator.
\end{itemize}
For a realized trajectory $\xi$, write
\[
\begin{aligned}
R_N(p;\xi)&:=\widehat U_{\mathrm{res}}(\widehat\pi,\widehat\omega,\widehat p,N,\xi),\\
Q_N(p;\xi)&:=\widehat U_{\mathrm{share}}(\widehat\pi,\widehat\omega,\widehat p,N,\xi),
\end{aligned}
\]
where these abbreviate Lean's \path{residualObserverFiberProcess} and \path{observerFiberPosteriorShareProcess}. Finally set $R_0(p):=\widehat U_{\mathrm{init}}(\widehat\pi,\widehat\omega,\widehat p)$, abbreviating Lean's \path{initialResidualObserverFiberOdds}, and write $\rho_\delta:=\mathrm{posteriorDecayFactorENNReal}(\delta)$.
\noindent\emph{Lean: \path{SemanticConvergence.ZeroOutRatePackage}.}
\end{definition}

\begin{theorem}[Full-support exploration floors recover the behavioral target]
\label{thm:exploration-floor-behavioral}
Fix a full-support exploration construction and let $\mathfrak K=(\mathcal I,\mathcal J,\Scal)$ be the semantic-learning package it induces. If this induced package satisfies Assumptions~\ref{ass:main-verified-package}, then for every target environment represented by the package,
\[
q_t^\star(\Scal(p^\star)\mid H_t)\longrightarrow 1
\quad\text{almost surely.}
\]
The full-support premise is therefore not an additional convergence theorem by itself; its role is to help verify the package's observer and support legs before the main theorem is applied.
\end{theorem}

\begin{proof}
\noindent\emph{Lean: \path{SemanticConvergence.Ontology.Coupling.h10_exploration_floor_behavioral}.}
Once the exploration construction has been packaged as an verified coupling, the conclusion is exactly Theorem~\ref{thm:separating-support-convergence}. The support-floor calculations serve as sufficient checks for the package assumptions, but they are not separate root hypotheses in this statement.
\end{proof}

The zero-out package gives a quantitative finite-horizon companion, phrased in terms of residual observer-fiber odds rather than as an unconditional statistical rate theorem.

\begin{theorem}[Zero-out residual-odds rate certificate]
\label{thm:separating-support-rate}
If $\mathfrak Z_T(U,\pi,p_{\mathrm{env}},p,\omega)$ is a zero-out rate package, then for every decay parameter $\delta$,
\[
R_N(p;\xi)\le \rho_\delta^N R_0(p)
\qquad
\text{for every }N\le T
\]
almost surely under the countable trajectory law induced by $(\widehat U,\widehat\pi,\widehat p_{\mathrm{env}})$.
\end{theorem}

\begin{proof}
\noindent\emph{Lean: \path{SemanticConvergence.zeroOutRatePackage_residualRate}.}
The package supplies the concrete bookkeeping and realized-update fields required by the Lean theorem; the conclusion is exactly the displayed almost-sure finite-horizon residual-odds bound.
\end{proof}

\begin{corollary}[Zero-out finite-time posterior-share lower bound]
\label{cor:separating-support-finite-time}
Under the assumptions of Theorem~\ref{thm:separating-support-rate},
\[
\left(1+\rho_\delta^N R_0(p)\right)^{-1}
\le
Q_N(p;\xi)
\qquad
\text{for every }N\le T
\]
almost surely under the same countable trajectory law.
\end{corollary}

\begin{proof}
\noindent\emph{Lean: \path{SemanticConvergence.zeroOutRatePackage_posteriorShareFiniteTime}.}
The finite initial residual-odds field of $\mathfrak Z_T$ converts the residual-odds bound of Theorem~\ref{thm:separating-support-rate} into the displayed posterior-share lower bound.
\end{proof}

\begin{remark}[Canonical support-budget specializations]
\label{rem:canonical-support-budgets}
Explicit lower bounds on separating-action mass are not themselves the rate theorem. The formal finite-list lemmas only certify nonzero support on listed actions; a concrete zero-out instantiation must still prove the realized-update field \textup{(Z2)} of a zero-out package. Once that field is proved, Definition~\ref{def:zero-out-rate-package} turns the supplied update into the finite-time bounds above.
\end{remark}

\begin{theorem}[Selector-labeled semantic convergence]
\label{thm:semantic-convergence}
Let a semantic-learning package $\mathfrak K=(\mathcal I,\mathcal J,\Scal)$ satisfy Assumptions~\ref{ass:main-verified-package}. Then for every target represented by $\mathfrak K$,
\[
q_t^\star(\Scal(p^\star)\mid H_t)\longrightarrow 1
\qquad
\text{almost surely.}
\]
The finite-list calculations in Propositions~\ref{prop:semantic-is-promotion-supporting} and~\ref{prop:finite-action-kernel-separation} are support checks for possible package instantiations; the convergence conclusion itself is discharged from the verified package theorem.
\end{theorem}

\begin{proof}
\noindent\emph{Lean: \path{SemanticConvergence.Ontology.Coupling.h10_semantic_convergence}.}
The cited theorem is the canonical-selector label for the main theorem. It takes the already-verified package $\mathfrak K$ and returns posterior concentration on $\Scal(p^\star)$ under the constructed Bayes/Gibbs law.
\end{proof}

\begin{remark}[Functional-minimization reading]
\label{rem:functional-payoff}
Under the two-observer functional, the belief-side posterior is fixed through the posterior-gap form behind Lemma~\ref{lem:variational}, and the action-side kernel is fixed through Proposition~\ref{prop:kernel-functional-minimizer}. The local finite-action and compact-action support lemmas, Propositions~\ref{prop:finite-action-kernel-separation} and~\ref{prop:compact-action-kernel-separation}, are package-checking tools: they help verify support fields but are not standalone convergence assumptions. Theorem~\ref{thm:kernel-semantic-convergence} is the exact variational candidate for constructing the support side of the Hellinger spine at the canonical pair $(\omega_B,\omega_A)=(\omega_{\mathrm{syn}},\omega_{\mathrm{behav}})$; Theorem~\ref{thm:semantic-convergence} is the selector wrapper, with both convergence statements routed through Assumptions~\ref{ass:main-verified-package}.
\end{remark}

\begin{remark}[Promotion-support generality]
\label{rem:promotion-support}
The proof invokes the semantic action interface only through the verified package fields it helps verify. Any action rule satisfying the same package, especially the class-predictive separation field \textup{(M7)}, obeys the same conclusion. Theorem~\ref{thm:semantic-convergence} is therefore a package-level convergence wrapper rather than a standalone proof that a particular selector constructs the package. The boundary predicates of Theorem~\ref{thm:afe-near-miss} describe how an AFE-side failure would be certified once the required witnesses are supplied.
\end{remark}

\begin{theorem}[Kernel-functional semantic convergence]
\label{thm:kernel-semantic-convergence}
Let the exact kernel-functional rule relative to its reference laws determine the action component of a semantic-learning package $\mathfrak K=(\mathcal I,\mathcal J,\Scal)$. If this package satisfies Assumptions~\ref{ass:main-verified-package}, then for every target represented by $\mathfrak K$,
\[
q_t^\star(\Scal(p^\star)\mid H_t)\longrightarrow 1
\qquad
\text{almost surely.}
\]
This is the paper's primary action-rule reading: exact Bayes/Gibbs belief and exact kernel-functional minimization are the agent-side legs \textup{(M1)--(M3)}, while the predictive-kernel observer, normalizers, Hellinger obligations, and class-predictive separation are the coupling legs \textup{(M4)--(M7)}.
\end{theorem}

\begin{proof}
\noindent\emph{Lean: \path{SemanticConvergence.Ontology.Coupling.h10_kernel_semantic_convergence}.}
The cited theorem is the kernel-functional label for the main theorem. In Lean it is a direct corollary of \path{SemanticConvergence.Ontology.Coupling.h10_verified_semantic_learning_package_converges}.
\end{proof}

\begin{corollary}[Compact-action extension of the kernel theorem and rates]
\label{cor:compact-action-kernel}
Suppose the compact-action kernel construction supplies a semantic-learning package $\mathfrak K=(\mathcal I,\mathcal J,\Scal)$ satisfying Assumptions~\ref{ass:main-verified-package}. Then under the constructed Bayes/Gibbs law,
\[
q_t^\star(\Scal(p^\star)\mid H_t)\longrightarrow 1
\qquad\text{almost surely.}
\]
If, for a finite horizon $T$, the same compact-action process is additionally accompanied by a zero-out rate package whose realized-update field uses the compact-action neighborhood floor
\[
\rho_\lambda(c^\star):=m_\lambda(r_{c^\star})
\]
for some $r_{c^\star}>0$ with $\varpi_{c^\star}(r_{c^\star})\le \eta(c^\star)/2$, then the finite-horizon residual and posterior-share bounds of Proposition~\ref{prop:kernel-exp-rate} and Theorem~\ref{thm:exp-rate-concentration} apply to the compact-action kernel process as zero-out companion certificates.
\end{corollary}

\begin{proof}
\noindent\emph{Lean: \path{SemanticConvergence.Ontology.Coupling.h10_compact_action_kernel}.}
The asymptotic convergence clause is the compact-action label for the main theorem. The neighborhood-floor calculation remains the compact-action way to verify the package's separation field; finite-horizon rates remain zero-out package consequences rather than consequences of asymptotic convergence alone.
\end{proof}

\begin{corollary}[Finite-class deterministic specialization]
\label{cor:finite-maximin}
If a finite-class maximin selector induces a semantic-learning package $\mathfrak K=(\mathcal I,\mathcal J,\Scal)$ satisfying Assumptions~\ref{ass:main-verified-package}, then
\[
q_t^\star(\Scal(p^\star)\mid H_t)\to 1
\qquad
\text{almost surely.}
\]
The finite maximin construction is thus a sufficient-condition recipe for the verified package, not a separate convergence root.
\end{corollary}

\begin{proof}
\noindent\emph{Lean: \path{SemanticConvergence.Ontology.Coupling.h10_finite_maximin}.}
Once the maximin construction has been certified as a verified package, the cited Lean theorem applies the main convergence theorem directly.
\end{proof}

\begin{theorem}[Zero supported separating actions preclude a support floor]
\label{cor:support-necessary}
Fix a concrete history $h$, observer $\omega$, target $p^\star$, finite action list $A_0$, and action law $\kappa$. If there is no action $a\in A_0$ with $\kappa(a)>0$ that separates the target observer fiber from its complement at $h$, then no positive separating-support floor exists: for every $\delta>0$,
\[
\neg\,\mathrm{hasSeparatingSupportFloor}
\bigl(\kappa,A_0,\mathrm{Sep}_{U,\pi,h,\omega,p^\star},\delta\bigr).
\]
\end{theorem}

\begin{proof}
\noindent\emph{Lean: \path{SemanticConvergence.cor_support_necessary}.}
The Lean theorem proves the contrapositive of the support-floor definition. A positive floor implies the existence of an action in the finite list with nonzero $\kappa$-mass and membership in the separating-action set; this contradicts the displayed no-supported-separator hypothesis.
\end{proof}

\begin{theorem}[Summable separating support is insufficient]
\label{thm:summable-support-insufficient}
Let $r:\mathbb N\to\mathbb Q$ be a summable rational action-mass schedule. Then there is a concrete two-action, two-observation semantic stack and a sequence of action laws $\kappa_n$ realizing exactly that schedule on a designated separating action such that, for every decay parameter $\delta$ and every time index $n$, the one-step zero-out residual-concentration predicate fails:
\[
\neg\,\mathrm{oneStepResidualPosteriorConcentrates}
\bigl(U,\pi,h,\omega,p^\star,A_0,\kappa_n,\delta\bigr).
\]
Thus the zero-out rate route needs an actual positive floor package; a merely summable schedule of separating-action mass is not enough to trigger the verified finite-horizon concentration predicate.
\end{theorem}

\begin{proof}
\noindent\emph{Lean: \path{SemanticConvergence.thm_summable_support_insufficient}.}
The Lean construction instantiates the finite Boolean action/observation machine, defines the scheduled action laws, proves that they realize the supplied summable rational schedule, and proves that no member of the schedule satisfies the one-step residual-concentration predicate for any $\delta$.
\end{proof}



\subsection{Zero-out rate certificates and same-view transfer}
\label{sec:exp-rate}

The verified rate surface in the manuscript is not an unconditional expectation/log-odds theorem. It is a finite-horizon certificate theorem: once the zero-out package of Definition~\ref{def:zero-out-rate-package} is verified, Lean proves one-step residual contraction, its iterated finite-horizon residual bound, and the corresponding lower bound on the posterior share of the observer fiber. This is the rate companion to the Hellinger-spine main route, not a substitute for the main route.

\begin{lemma}[Zero-out one-step residual contraction]
\label{lem:one-step-drift}
If $\mathfrak Z_T(U,\pi,p_{\mathrm{env}},p,\omega)$ is a zero-out rate package, then for every decay parameter $\delta$,
\[
R_{n+1}(p;\xi)\le \rho_\delta R_n(p;\xi)
\qquad
\text{for every }n<T
\]
almost surely under the constructed countable trajectory law.
\end{lemma}

\begin{proof}
\noindent\emph{Lean: \path{SemanticConvergence.zeroOutRatePackage_oneStepResidual}.}
The package field \textup{(Z2)} supplies realized prefix residual updates. Lean converts those updates into a supportwise residual-contraction witness and then pushes the witness through the countable trajectory law to obtain the displayed almost-sure one-step inequality.
\end{proof}

\begin{proposition}[Zero-out residual-rate certificate]
\label{prop:exp-rate}
If $\mathfrak Z_T(U,\pi,p_{\mathrm{env}},p,\omega)$ is a zero-out rate package, then for every decay parameter $\delta$,
\[
R_N(p;\xi)\le \rho_\delta^N R_0(p)
\qquad
\text{for every }N\le T
\]
almost surely under the constructed countable trajectory law.
\end{proposition}

\begin{proof}
\noindent\emph{Lean: \path{SemanticConvergence.zeroOutRatePackage_expRate}.}
This is the iterated form of Lemma~\ref{lem:one-step-drift}. Lean proves the rate bound by transporting the realized-update field to a supportwise recurrence and applying the countable finite-horizon residual-rate theorem.
\end{proof}

\begin{lemma}[Same-view transfer of residual-contraction witnesses]
\label{lem:one-step-drift-kernel}
On the countable support/rate stack, if two programs $p$ and $q$ lie in the same observer fiber and $p$ has a supportwise residual-contraction witness at decay parameter $\delta$ through horizon $T$, then $q$ has the same witness.
\end{lemma}

\begin{proof}
\noindent\emph{Lean: \path{SemanticConvergence.lem_one_step_drift_kernel}.}
The residual odds and their contraction witness depend only on the observer fiber. Equality of $\omega$-views therefore transports the witness from $p$ to $q$ without changing $\delta$ or $T$.
\end{proof}

\begin{proposition}[Same-view zero-out residual-rate transfer]
\label{prop:kernel-exp-rate}
If $\mathfrak Z_T(U,\pi,p_{\mathrm{env}},p,\omega)$ is a zero-out rate package and $q$ lies in the same $\omega$-fiber as $p$, then for every decay parameter $\delta$,
\[
R_N(q;\xi)\le \rho_\delta^N R_0(q)
\qquad
\text{for every }N\le T
\]
almost surely under the constructed countable trajectory law.
\end{proposition}

\begin{proof}
\noindent\emph{Lean: \path{SemanticConvergence.zeroOutRatePackage_sameViewResidualRate}.}
Lean applies Lemma~\ref{lem:one-step-drift-kernel} to transport the supportwise witness from $p$ to $q$, then applies the same finite-horizon residual-rate theorem to the transported witness.
\end{proof}

\begin{theorem}[Same-view finite-time posterior-share transfer]
\label{thm:exp-rate-concentration}
If $\mathfrak Z_T(U,\pi,p_{\mathrm{env}},p,\omega)$ is a zero-out rate package and $q$ lies in the same $\omega$-fiber as $p$, then for every decay parameter $\delta$,
\[
\left(1+\rho_\delta^N R_0(q)\right)^{-1}
\le
Q_N(q;\xi)
\qquad
\text{for every }N\le T
\]
almost surely under the constructed countable trajectory law.
\end{theorem}

\begin{proof}
\noindent\emph{Lean: \path{SemanticConvergence.zeroOutRatePackage_sameViewPosteriorShareFiniteTime}.}
The same-view residual-rate theorem gives the residual-odds upper bound for $q$. The package field \textup{(Z3)} supplies the finite initial denominator, and Lean converts the residual-odds bound into the displayed posterior-share lower bound.
\end{proof}

\begin{corollary}[Goal-dialed separation payoff]
\label{cor:goal-dialed-payoff}
Fix an observer-dialed action construction and let it induce a semantic-learning package $\mathfrak K=(\mathcal I,\mathcal J,\Scal)$. If the construction satisfies Assumptions~\ref{ass:main-verified-package}, then
\[
q_t^\star(\Scal(p^\star)\mid H_t)\to1
\qquad
\text{almost surely.}
\]
Uniform history-independent separation remains a useful sufficient-condition check for the separation field of the package; if a zero-out package is additionally verified for a finite horizon, the finite-horizon bounds of Proposition~\ref{prop:exp-rate} and Theorem~\ref{thm:exp-rate-concentration} apply as companion certificates.
\end{corollary}

\begin{proof}
\noindent\emph{Lean: \path{SemanticConvergence.Ontology.Coupling.h10_goal_dialed_payoff}.}
The cited theorem is the observer-dialed label for the main theorem. The local uniform-separation Lean wrapper \path{SemanticConvergence.cor_goal_dialed_payoff} remains the sufficient-condition lemma used to verify the package's separation field.
\end{proof}

Propositions~\ref{prop:exp-rate} and~\ref{prop:kernel-exp-rate}, together with Theorem~\ref{thm:exp-rate-concentration}, are therefore best read as certificate consumers: they do not assert that the canonical action rule automatically creates a zero-out witness, but they prove the exact finite-horizon consequences once such a witness is present. Corollary~\ref{cor:goal-dialed-payoff} records the analogous observer-dialed separation input. The observer-order restrictions remain part of the verified package rather than a consequence of the local rate certificates.

\part{Reach: Limits, Reductions, and Realization}
\label{part:reach}

Part~\ref{part:reach} extends the framework's reach beyond the convergence theorem itself. Section~\ref{sec:near-miss} shows that the most natural information-gain rule, even under the universal prior, does not in general converge to the interventional class --- the negative result that disciplines what the main theorem must look like. Section~\ref{sec:self-ref} dissolves the apparent uncomputability of the behavioral observer through a finite-time proxy and an eventual-isolation interface. Section~\ref{sec:realization} exhibits a differentiable surrogate loss whose minimizer realizes the verified package's coupling conditions, bringing the framework into deep learning. Section~\ref{sec:diagnostic} turns the framework into a procedure for diagnosing any specific interactive learner.

\section{The Algorithmic Free Energy Principle as Boundary Example}
\label{sec:boundary-example}
\label{sec:near-miss}

The natural unification of the three traditions --- universal prior, Bayes/Gibbs belief, and expected-free-energy action under uniform preferences and zero action cost --- combines the ingredients each tradition offers. The result is the \emph{algorithmic free energy principle}: a learner whose belief side minimizes the algorithmic free energy and whose action side maximizes the corresponding expected free energy. One might expect this principle, having all the right ingredients, to converge on $[p^\star]$ in any well-formed environment.

It does not. Theorem~\ref{thm:afe-near-miss} below exhibits an environment, constructed from a substituted finite-support prior, on which the algorithmic free energy principle fails to concentrate on $[p^\star]$ for every finite horizon. The realized policy under this principle induces an observer strictly coarser than $\omega_{\mathrm{behav}}$, and the posterior on $[p^\star]$ is frozen at its prior mass throughout.

\begin{center}
\fbox{\parbox{0.93\textwidth}{\centering
\textbf{Negative result (Theorem~\ref{thm:afe-near-miss}).}\\[0.4em]
The algorithmic free energy principle does not in general converge to the interventional class $[p^\star]$.\\
Even under the universal prior, the natural information-gain rule can leave the posterior on $[p^\star]$ frozen at its prior mass for every finite horizon.
}}
\end{center}

\begin{principle}[Curiosity vs.\ explanation]
\label{prin:curiosity-vs-explanation}
A learning rule converges to the interventional semantic class iff its action criterion separates the target's class from its complement under the true environment law. Generic information gain --- \emph{curiosity} --- does not suffice; class-against-complement separation --- \emph{explanation} --- does.
\end{principle}

Theorems~\ref{thm:afe-near-miss} and~\ref{thm:separating-support-convergence} together make this principle precise. The algorithmic free energy principle is a curiosity rule: under uniform preferences and zero action cost, its action term reduces to information gain (Corollary~\ref{cor:efe-specialization}). It does not in general converge to $[p^\star]$. The kernel-functional refinement of Section~\ref{sec:main} is an explanation rule: its score is the class-against-complement Bhattacharyya separation, computed against the true environment law. It does converge under the verified package's coupling conditions. The boxed problem of Section~\ref{sec:intro} thus has a sharp answer: curiosity is necessary --- the agent must explore --- but not sufficient; explanation is what suffices.

This is the negative finding that disciplines what the main theorem must look like. The kernel-functional refinement of Section~\ref{sec:main} is not the natural rule plus extra decoration; it is what the natural rule has to be replaced by, because the natural rule itself does not work. The remainder of this section records the observer-invariance and risk-minus-gain definitional shape of the boundary objective, the singleton argmin action rule used by the formal proxy, the near-miss witness predicates, and the corresponding observer-promotion failure and contrast statements.

\subsection{Boundary identification}
\label{sec:boundary-identification}

\begin{proposition}[Boundary objective is invariant on observer fibers]
\label{prop:boundary-identity}
If two target representatives have the same observer view, then the boundary expected-free-energy proxy assigns them the same value at every fixed history and action.
\end{proposition}

\begin{proof}
\noindent\emph{Lean: \path{SemanticConvergence.prop_boundary_identity}.}
Both the risk proxy and the information-gain proxy are invariant under equality of observer views, so their difference is invariant too.
\end{proof}

\begin{remark}[Reading under the functional]
\label{rem:boundary-identity}
Proposition~\ref{prop:boundary-identity} is a boundary-objective invariance statement. It does not by itself prove the classical expected-free-energy identification or a near-miss construction; those would require additional formal hypotheses beyond this proxy.
\end{remark}

\subsection{Expected free energy}

\begin{definition}[Expected free energy]
\label{def:efe}
The boundary expected-free-energy proxy at $(h_t,a,\omega,p^\star)$ is
\[
\Gcal_t(a;h_t,\omega,p^\star)
:=
\operatorname{Risk}_t(a;h_t,\omega,p^\star)
-
\operatorname{Gain}_t(a;h_t,\omega,p^\star),
\]
where the risk term is the observer-fiber semantic-separation proxy and the gain term is the observer-fiber semantic-gain proxy in the Lean boundary stack.
\end{definition}

\begin{lemma}[Risk minus information gain]
\label{lem:risk-ig}
For every history $h_t$ and action $a$,
\[
\Gcal_t(a;h_t,\omega,p^\star)
=
\operatorname{Risk}_t(a;h_t,\omega,p^\star)
-
\operatorname{Gain}_t(a;h_t,\omega,p^\star).
\]
\end{lemma}

\begin{proof}
\noindent\emph{Lean: \path{SemanticConvergence.lem_risk_ig}.}
This is the definition of the boundary proxy.
\end{proof}

\begin{corollary}[Boundary proxy specialization]
\label{cor:efe-specialization}
The expected-free-energy proxy is the same risk-minus-gain expression of Lemma~\ref{lem:risk-ig}.
\end{corollary}

\begin{proof}
\noindent\emph{Lean: \path{SemanticConvergence.cor_efe_specialization}.}
This corollary is the same definitional equality exposed under the specialization label.
\end{proof}

\begin{remark}
\label{rem:efe-unification}
Classical expected-free-energy identities motivate this proxy, but the formal corollary here is only the Lean risk-minus-gain equality.
\end{remark}

\subsection{The algorithmic free energy principle}

\begin{definition}[Algorithmic free energy principle]
\label{def:afe-principle}
A learner obeys the \emph{algorithmic free energy principle} if
at a finite action list when it chooses the singleton action law concentrated on an argmin of the boundary expected-free-energy proxy over that list.
\end{definition}

\begin{remark}[The AFE principle as a Bayes/Gibbs-plus-MI boundary specialization]
\label{rem:afe-as-functional-minimizer}
Under the Lean-facing proxy, Definition~\ref{def:afe-principle} is an argmin selector for the boundary objective. It is not a complete formalization of the full active-inference expected-free-energy principle.
\end{remark}

\subsection{Information decomposition relative to $[p^\star]$}

\begin{lemma}[Information decomposition]
\label{lem:info-decomp}
For every target representative $p^\star$, history $h_t$, action $a$, and observer $\omega$,
\[
\left(
\operatorname{Odds}_t^\omega(p^\star\mid h_t),
\operatorname{Gain}_t(a;h_t,\omega,p^\star)
\right)
=
\left(
\operatorname{Odds}_t^\omega(p^\star\mid h_t),
S_t^\omega(p^\star;h_t)
\right).
\]
\end{lemma}

\begin{proof}
\noindent\emph{Lean: \path{SemanticConvergence.lem_info_decomp}.}
This is a definitional pair identity: the boundary information-gain term is the semantic-gain proxy.
\end{proof}

\begin{remark}
Lemma~\ref{lem:info-decomp} records the formal bookkeeping relation used by the boundary proxy. Classical mutual-information decompositions remain motivation rather than a theorem of this Lean statement.
\end{remark}

\subsection{Near-miss: the AFE principle need not converge to $[p^\star]$}

Prior work on AIXI motivates near-miss examples. The formal theorem here is a witness package: if a boundary action preference witness and a frozen-posterior-through-horizon predicate are supplied, Lean packages them into the near-miss conclusion below.

\begin{theorem}[Boundary near-miss witness package]
\label{thm:afe-near-miss}
Suppose a finite action list contains actions $a_{\mathrm{bad}}$ and $a_{\mathrm{good}}$ such that the boundary proxy weakly prefers $a_{\mathrm{bad}}$, the two actions are distinct, and the action-observer semantic separation is positive. Suppose also that a posterior-mass process is frozen at $\alpha_0$ through horizon $T$. Then Lean returns both the action-level near-miss witness and the finite-horizon non-improvement witness: for every threshold $\varepsilon$ with $\alpha_0<\varepsilon<1$, some $n\le T$ has posterior mass below $\varepsilon$.
\end{theorem}

\begin{proof}
\noindent\emph{Lean: \path{SemanticConvergence.thm_afe_near_miss}.}
The proof packages the supplied action-level witness and uses the frozen-posterior predicate at time zero to produce the non-improvement witness for each threshold.
\end{proof}

\begin{theorem}[Observer-promotion failure predicate]
\label{thm:observer-promotion-failure}
If a pair of observers satisfies the formal observer-promotion failure predicate, then that observer-promotion failure predicate holds.
\end{theorem}

\begin{proof}
\noindent\emph{Lean: \path{SemanticConvergence.thm_observer_promotion_failure}.}
This theorem is the identity wrapper for the failure predicate.
\end{proof}

\begin{corollary}[Observer-promotion failure transfers unchanged]
\label{cor:observer-promotion-universal}
If the observer-promotion failure predicate holds, then it holds under the universal-prior label as well.
\end{corollary}

\begin{proof}
\noindent\emph{Lean: \path{SemanticConvergence.cor_observer_promotion_universal}.}
This is the same predicate as Theorem~\ref{thm:observer-promotion-failure}.
\end{proof}

\begin{remark}
\label{rem:near-miss}
The formal near-miss statements are witness predicates. They identify what must be supplied to certify a boundary failure, rather than constructing a full classical counterexample from first principles.
\end{remark}

\begin{remark}
\label{rem:near-miss-universal}
Extending a full frozen-posterior near-miss construction to the literal universal interactive prior remains outside the formal boundary proxy. Corollary~\ref{cor:observer-promotion-universal} only transports the supplied observer-promotion failure predicate.
\end{remark}

\subsection{Observer-promotion contrast: AFE vs.\ promotion-supporting rule}
\label{sec:promotion-contrast}

The promotion-contrast corollary packages the two supplied predicates: a boundary near-miss witness and an observer-promotion failure witness.

\begin{corollary}[Observer-promotion contrast]
\label{cor:promotion-contrast}
If the boundary near-miss predicate and the observer-promotion failure predicate both hold, then the formal promotion-contrast predicate holds.
\end{corollary}

\begin{proof}
\noindent\emph{Lean: \path{SemanticConvergence.cor_promotion_contrast}.}
The contrast predicate is defined as the conjunction of those two supplied predicates.
\end{proof}

\begin{remark}
\label{rem:promotion-contrast-interpretation}
Corollary~\ref{cor:promotion-contrast} is a bookkeeping wrapper for the two boundary predicates. A fully constructive near-miss family would require hardening these predicates into explicit environments and priors.
\end{remark}


\subsection*{Open question}

When do generic information gain and semantic gain coincide? The boundary proxy of Definition~\ref{def:afe-principle} makes this the right comparison question, but the paper does not prove a full structural equivalence or construct a counterexample. Conditions on the hypothesis class under which expected-free-energy maximization automatically maximizes semantic gain would identify the precise boundary between generic epistemic discovery and explanatory action.

\section{Self-Referential Reduction of the Computability Gap}
\label{sec:self-ref}

Theorem~\ref{thm:semantic-convergence} drives posterior concentration through the verified package. Proposition~\ref{prop:goal-dialed} is only the fixed-observer admissibility proxy. A learner running an observer-indexed rule still faces a computability gap: two programs are behaviorally equivalent when $\mu_{p_1}=\mu_{p_2}$ as interactive laws on the \emph{whole} policy family, and no finite history certifies this. The rule as written therefore motivates finite-time observer proxies.

This section isolates the tension at a finitary level rather than claiming an unconditional closure of it. The proxy object is the \emph{finite-time policy partition} of $\Pcal$ induced by the length-$t$ marginal of $\Prob_{p,\pi}$. The verified results below are intentionally modest: they expose the monotone finite-time observer, define a generic self-referential chooser through that observer, and state the exact eventual-isolation predicates consumed by later package-level convergence statements.

\subsection{The finite-time policy observer}
\label{sec:finite-time-policy-observer}

\begin{definition}[Finite-time policy observer]
\label{def:finite-time-policy-observer}
Fix a policy $\pi$ and $t\ge 0$. The \emph{finite-time policy observer} $\omega_\pi^t$ is the observer of Definition~\ref{def:observer} with valuation
\[
V_{\omega_\pi^t}(p)\;:=\;\Prob_{p,\pi}(H_t\in\cdot),
\]
the law of the length-$t$ history under $(p,\pi)$. Write $[p^\star]_{\omega_\pi^t}$ for the $\omega_\pi^t$-indistinguishability class of $p^\star$.
\end{definition}

\begin{lemma}[One-step finite-time observer refinement]
\label{lem:monotone-refinement}
For every policy $\pi$ and every $t\ge 0$,
\[
\omega_\pi^t\;\le\;\omega_\pi^{t+1}.
\]
\end{lemma}

\begin{proof}
\noindent\emph{Lean: \path{SemanticConvergence.lem_monotone_refinement}.}
The length-$t$ observer is a projection of the length-$(t+1)$ observer, so equality of length-$(t+1)$ views implies equality of length-$t$ views.
\end{proof}

The finite-time policy observer is finitary in a sense the $\omega_{\mathrm{behav}}$-partition is not: for fixed $(\pi,t)$, checking whether $p_1\sim_{\omega_\pi^t}p_2$ is equality of two probability measures on the countable set of histories of length $t$, equivalently equality on all length-$t$ cylinder events. This replaces full interactive-law equality by a fixed-horizon marginal comparison. Under the standing finite-action/countable-observation setup, however, we do not claim exact decidability of this comparison.

\subsection{The self-referential semantic rule}
\label{sec:self-ref-rule}

\begin{definition}[Self-referential semantic rule]
\label{def:self-ref-rule}
Given a reference policy $\pi_{\mathrm{ref}}$ and an observer-driven chooser $\chi$, the self-referential proxy policy is
\[
\pi^{\mathrm{sr}}(h)
:=
\chi\bigl(|h|,h,\omega_{\pi_{\mathrm{ref}}}^{|h|}\bigr).
\]
That is, at each history it passes the current finite-time policy observer to the chooser.
\end{definition}

The Lean definition is a generic interface: it does not hard-code a class-against-complement selector, but supplies the finite-time observer to whatever chooser is provided.

\begin{remark}[Self-referentiality vs.\ bootstrapping]
\label{rem:self-reference-scope}
The rule is self-referential in the narrow sense that the chooser can read a finite-time observer derived from a policy prefix. The formal object is intentionally only this interface; stronger claims about eventual behavioral recovery require additional package assumptions.
\end{remark}

\subsection{Observer promotion under the self-referential rule}
\label{sec:observer-promotion-sr}

The key missing step is promotion of the running partition to the behavioral one. The following statements record the exact local predicates that the Lean self-reference interface exposes.

\begin{lemma}[Promotion-floor reachability]
\label{lem:exploration-reachability}
If the one-step split predicate holds at $(U,\pi,t,h_t,A_0,\omega,p^\star)$, then there is a listed action $a\in A_0$ with positive semantic separation at that same history and observer.
\end{lemma}

\begin{proof}
\noindent\emph{Lean: \path{SemanticConvergence.lem_exploration_reachability}.}
The one-step split predicate already contains the listed action and positive-separation witness; the proof projects those fields.
\end{proof}

\begin{proposition}[Eventual isolation criterion]
\label{prop:observer-promotion-sr}
If at some finite time the finite-time observer fiber of $p^\star$ is contained in the interventional semantic class $[p^\star]$, then the policy satisfies the formal eventual-isolation predicate.
\end{proposition}

\begin{proof}
\noindent\emph{Lean: \path{SemanticConvergence.prop_observer_promotion_sr}.}
The eventual-isolation predicate is an existential over such a finite time, so the supplied containment witness proves it directly.
\end{proof}

Proposition~\ref{prop:observer-promotion-sr} isolates the exact predicate still needed by this proxy route. It does not prove that a concrete self-referential learner reaches eventual isolation.

\subsection{Self-referential convergence}
\label{sec:self-ref-convergence}

\begin{theorem}[Self-referential convergence under eventual isolation]
\label{thm:self-ref-convergence}
For any policy $\pi$, if the formal eventual-isolation predicate holds for $p^\star$, then that same eventual-isolation predicate holds for $p^\star$.
\end{theorem}

\begin{proof}
\noindent\emph{Lean: \path{SemanticConvergence.thm_self_ref_convergence}.}
This theorem is the identity wrapper for the eventual-isolation predicate.
\end{proof}

\begin{proposition}[Unconditional eventual isolation fails]
\label{prop:self-ref-obstruction}
If every finite-time observer fiber of $p^\star$ contains some program outside $[p^\star]$, then the formal isolation-obstructed predicate holds.
\end{proposition}

\begin{proof}
\noindent\emph{Lean: \path{SemanticConvergence.prop_self_ref_obstruction}.}
This is exactly the definition of the obstruction predicate.
\end{proof}

\subsection{Exploration-completed self-reference}
\label{sec:self-ref-exploration}

The obstruction predicate of Proposition~\ref{prop:self-ref-obstruction} shows what this proxy route still needs. The exploration-completed definitions below are Lean-facing interfaces for switching to an exploration law when a trigger fires; convergence from exploration remains routed through the verified package, not through these local self-reference wrappers.

\begin{definition}[Exploration-completed self-referential rule]
\label{def:self-ref-exploratory}
Given a reference policy, an observer-driven chooser, an exploration policy, and a Boolean trigger, the exploration-completed self-referential rule uses the exploration policy at a history when the trigger is true and otherwise uses the self-referential chooser from Definition~\ref{def:self-ref-rule}.
\end{definition}

\begin{theorem}[Exploration branch is selected when the trigger fires]
\label{thm:self-ref-exploratory}
At any history $h$, if the exploration trigger is true, then the exploration-completed policy equals the exploration policy at $h$.
\end{theorem}

\begin{proof}
\noindent\emph{Lean: \path{SemanticConvergence.thm_self_ref_exploratory}.}
Unfold Definition~\ref{def:self-ref-exploratory} and simplify using the trigger equality.
\end{proof}


\begin{theorem}[Exploration branch preserves positive action mass]
\label{thm:self-ref-exploratory-rate}
At any triggered history, if the exploration policy assigns nonzero mass to an action, then the exploration-completed policy assigns nonzero mass to that action.
\end{theorem}

\begin{proof}
\noindent\emph{Lean: \path{SemanticConvergence.thm_self_ref_exploratory_rate}.}
The triggered branch is definitionally the exploration law, so nonzero mass transfers by rewriting.
\end{proof}

\begin{proposition}[One-step splitting criterion]
\label{prop:self-ref-one-step-split}
If the one-step split predicate holds at $(U,\pi,t,h_t,A_0,\omega,p^\star)$, then there is a listed action $a\in A_0$ with positive semantic separation at that same history and observer.
\end{proposition}

\begin{proof}
\noindent\emph{Lean: \path{SemanticConvergence.prop_self_ref_one_step_split}.}
This is the same witness projection as Lemma~\ref{lem:exploration-reachability}.
\end{proof}

\begin{theorem}[Sharpened self-referential convergence under deterministic splitting]
\label{thm:self-ref-sharp}
If a finite-time observer fiber is already contained in $[p^\star]$ and a one-step split witness is available, then both the isolation witness and the listed separating action witness are available.
\end{theorem}

\begin{proof}
\noindent\emph{Lean: \path{SemanticConvergence.thm_self_ref_sharp}.}
The first component is the supplied fiber-containment witness. The second component is obtained from the one-step split witness by Proposition~\ref{prop:self-ref-one-step-split}.
\end{proof}

\begin{remark}[Finite-time proxy interface]
\label{rem:self-ref-computability}
At every finite step $t$, the self-referential interface depends only on the policy prefix and the corresponding finite-time observer. Under the standing finite-action/countable-observation setup, those marginals are generally countably rather than finitely supported, so the relevant objects are semicomputable cylinder probabilities rather than finite tables. The present section therefore does not claim exact computability or convergence of the running partition; it records the formal predicates needed by any future hardening of that route.
\par
Theorems~\ref{thm:separating-support-convergence} and~\ref{thm:exploration-floor-behavioral} take the verified route: they replace the running-partition argument by the verified package conditions. The exploration-completed self-reference wrappers above only expose the branch-selection and support-inheritance facts used by such package checks.
\end{remark}

\begin{remark}[Relation to Proposition~\ref{prop:goal-dialed}]
\label{rem:self-ref-vs-goal-dialed}
Proposition~\ref{prop:goal-dialed} is a fixed-observer admissibility statement, not a standalone convergence theorem. Theorem~\ref{thm:self-ref-convergence} is likewise only the identity wrapper for eventual isolation. The convergence payoff for goal-dialed or self-referential variants must therefore be routed through Assumptions~\ref{ass:main-verified-package} or through the zero-out package, not inferred from these proxy statements alone.
\par
Theorems~\ref{thm:separating-support-convergence} and~\ref{thm:exploration-floor-behavioral} lie alongside these proxy statements as the actual verified convergence route.
\end{remark}

The self-referential section provides a proxy-layer interface. It defines finite-time observer interfaces and records the exact local witnesses the formal development proves. The paper's verified convergence route remains Theorem~\ref{thm:separating-support-convergence}; finite-horizon residual/posterior-share statements remain zero-out package consequences.

\section{A Differentiable Realization}
\label{sec:realization}
\label{sec:amortized-surrogate}

The four-set hierarchy clarifies where existing learning theories stand. Each theory targets one of the four observers of Definition~\ref{def:observer}; the question for any given theory is whether its learning rule promotes the realized observer to the observer whose indistinguishability class the theory tries to identify.

Realized-history fit aims at $R_{h_t}(p^\star)$. Next-token training objectives under any fixed data distribution are instances: the learner scores well by assigning high likelihood to realized prefixes. Lemma~\ref{lem:fit-gap} is the boundary: such a learner's posterior can be on a compatibility set whose next-queried continuation is still undetermined. Lemma~\ref{lem:merging} records the corresponding observer-fiber posterior invariance in the formal stack.

Policy-relative prediction aims at $R_\pi(p^\star)$. Off-policy reinforcement learning and partially-observed Markov decision processes without active interventions face the boundary of Theorem~\ref{thm:policy-gap}: the posterior can concentrate on $R_\pi(p^\star)$ while leaving $[p^\star]$ unresolved.

Interventional semantic recovery aims at the observer fiber corresponding to the target environment. In the verified theorem this target is the packaged semantic set $\Scal(p^\star)$. The boundary proxy of Definition~\ref{def:afe-principle} is useful for describing action-side failure predicates, while Section~\ref{sec:semantic-action} records the action-side interfaces used to verify H10 packages. Concrete systems built around program-level structure --- Bayesian program learning and program induction \citep{LakeEtAl2015,EllisEtAl2021DreamCoder}, abstraction-and-reasoning benchmarks \citep{CholletARC2019}, and search over program spaces guided by large language models \citep{RomeraParedesEtAl2024FunSearch} --- supply motivating examples in the spirit of class-targeted learning. A strict instantiation of any of these in the present framework would require explicitly defining the program space, observer, action set, and semantic class for the system in question, which we do not undertake here.

Canonical syntactic recovery aims at $\{p^\star\}$. Theorem~\ref{thm:factor-through-quotient} rules it out unless the designated representative is already constant on $[p^\star]$.

Of these four targets, only interventional semantic recovery is the object of the convergence theorem of Section~\ref{sec:main}. Most deployed learning systems aim at one of the coarser targets. The question this section answers is whether the convergence theorem can be realized by an objective that an interventional learner can actually optimize: a differentiable loss whose minimizer satisfies the verified package's coupling conditions. The answer is yes, conditional on three explicit deployment-side assumptions on the learned model.

The class-against-complement action rule of Section~\ref{sec:semantic-action} is the explanation criterion that the Principle of Section~\ref{sec:near-miss} names. This section realizes that criterion as a differentiable loss function trainable in deep learning. The framework's epistemological move --- explanation, not just curiosity --- is here packaged as an objective that gradient descent can minimize.

The exact kernel lift $\mathcal J_{t,\mathrm{ker}}^{\omega_{\mathrm{behav}}}$ is the paper's strongest variational answer. A practical learner will typically replace the countable program space by a finite latent class model and optimize a differentiable surrogate instead. The theorem-level question is then not whether the surrogate reduces uncertainty in general, but whether minimizing it still certifies the separating-support condition of Theorem~\ref{thm:separating-support-convergence}.

Assume in this subsection that the deployed surrogate uses a finite latent class set $\widehat{\Ccal}$ and a finite observation alphabet $\Ocal$. Assume also that its parameter state is updated causally from history: for each $t$, the maps
\[
h_t\longmapsto \rho_t(\cdot\mid h_t)
\qquad\text{and}\qquad
h_t\longmapsto \widehat\mu_{\theta_t}(o\mid h_t,a,\hat c)
\]
are measurable for every $o\in\Ocal$, $a\in\Acal$, and $\hat c\in\widehat{\Ccal}$. In particular, any class-action kernel built from these quantities is an adapted policy kernel in the sense of the formal setup. Let $\rho_t(\hat c\mid h_t)$ be a differentiable full-support latent-class prior over $\widehat{\Ccal}$, and let $\widehat\mu_{\theta_t}(o\mid h_t,a,\hat c)$ be a differentiable predictive model. For each $\hat c\in\widehat{\Ccal}$ define the learned complement law
\[
\widehat Q_{t,\theta_t}^{-\hat c}(o\mid h_t,a)
:=
\sum_{\hat d\neq \hat c}
\frac{\rho_t(\hat d\mid h_t)}{1-\rho_t(\hat c\mid h_t)}
\widehat\mu_{\theta_t}(o\mid h_t,a,\hat d),
\]
the learned Bhattacharyya score
\[
\widehat B_{t,\theta_t}(\hat c,a;h_t)
:=
-\log\sum_{o\in\Ocal}
\sqrt{
\widehat\mu_{\theta_t}(o\mid h_t,a,\hat c)\,
\widehat Q_{t,\theta_t}^{-\hat c}(o\mid h_t,a)},
\]
and its cap
\[
\bar{\widehat B}_{t,\theta_t}(\hat c,a;h_t)
:=
\min\{1,\widehat B_{t,\theta_t}(\hat c,a;h_t)\}.
\]
For any differentiable predictive loss $\mathcal L_t^{\mathrm{pred}}(\theta_t;h_t)$ independent of the action kernel, define the amortized surrogate
\[
\boxed{\;
\mathcal L_{t,\mathrm{impl}}(\theta_t,\kappa;h_t)
\;:=\;
\mathcal L_t^{\mathrm{pred}}(\theta_t;h_t)
-\beta\!
\sum_{\hat c\in\widehat{\Ccal}}\sum_{a\in\Acal}
\kappa(\hat c,a\mid h_t)\,\bar{\widehat B}_{t,\theta_t}(\hat c,a;h_t)
+\gamma\,
\KL\!\bigl(\kappa(\cdot,\cdot\mid h_t)\,\|\,\rho_t(\cdot\mid h_t)\otimes\lambda\bigr)
\;}
\]
where $\kappa(\cdot,\cdot\mid h_t)\in\Delta(\widehat{\Ccal}\times\Acal)$ and $\lambda$ is a full-support reference distribution on $\Acal$.

\bigskip

\noindent
Table~\ref{tab:loss-classification} situates the surrogate among common training losses.

\begin{table}[h]
\centering\small
\renewcommand{\arraystretch}{1.25}
\begin{tabular}{p{0.27\linewidth}p{0.07\linewidth}p{0.07\linewidth}p{0.07\linewidth}p{0.41\linewidth}}
\toprule
Loss family & (A1) & (A2) & (A3) & Why this is not an explanation rule \\
\midrule
Next-token / cross-entropy & {?} & --- & --- & No class-against-complement structure; rewards likelihood, not separation. \\
Variational ELBO (parametric family) & {?} & {?} & --- & Restricted variational family generally breaks KL-necessity; the action term is implicit. \\
Reward maximization (RL) & --- & {?} & {?} & No designated target latent class; reward couples agent to environment opaquely. \\
$Q$-learning / temporal-difference & --- & --- & --- & As above, plus the score is bootstrap-defined rather than Bhattacharyya. \\
Mutual-information curiosity & {?} & {?} & --- & The natural-information-gain failure of Theorem~\ref{thm:afe-near-miss}. \\
\textbf{Amortized surrogate $\mathcal L_{t,\mathrm{impl}}$} & \checkmark$^*$ & \checkmark$^*$ & \checkmark$^*$ & Class-against-complement score is explicit; the $^*$ marks the deployment-side hypotheses. \\
\bottomrule
\end{tabular}
\caption{Common training losses against the verified package's deployment-side hypotheses (A1)--(A3). A `?' marks an assumption the loss does not directly enforce; `---' marks one that its standard form does not allow. The amortized surrogate is the smallest extension of next-token / variational training that explicitly encodes class-against-complement separation.}
\label{tab:loss-classification}
\end{table}

\begin{proposition}[Exact action-side minimizer of the amortized surrogate]
\label{prop:amortized-surrogate-minimizer}
Fix $h_t$ and $\theta_t$. Any minimizer of $\mathcal L_{t,\mathrm{impl}}(\theta_t,\kappa;h_t)$ over $\kappa(\cdot,\cdot\mid h_t)\in\Delta(\widehat{\Ccal}\times\Acal)$ is
\[
\kappa_{t,\theta_t}^{\mathrm{impl}}(\hat c,a\mid h_t)
=
\frac{
\rho_t(\hat c\mid h_t)\lambda(a)\,
\exp\!\bigl((\beta/\gamma)\bar{\widehat B}_{t,\theta_t}(\hat c,a;h_t)\bigr)}
{\sum_{\hat d\in\widehat{\Ccal}}\sum_{b\in\Acal}
\rho_t(\hat d\mid h_t)\lambda(b)\,
\exp\!\bigl((\beta/\gamma)\bar{\widehat B}_{t,\theta_t}(\hat d,b;h_t)\bigr)}.
\]
\end{proposition}

\begin{proof}
\noindent\emph{Lean: \path{SemanticConvergence.prop_amortized_surrogate_minimizer}.}
For fixed $(\theta_t,h_t)$ the predictive term is constant in $\kappa$, so the minimization over $\kappa$ is exactly the Gibbs-kernel calculation of Proposition~\ref{prop:kernel-functional-minimizer} with $\rho_t(\cdot\mid h_t)$ in place of the exact class prior and $\bar{\widehat B}_{t,\theta_t}$ in place of the exact class score.
\end{proof}

\begin{theorem}[Semantic-recovery guarantee for the amortized surrogate]
\label{thm:amortized-surrogate}
Let the deployed policy be the action marginal
\[
\pi_{t+1}^{\mathrm{impl}}(a\mid h_t)
:=
\sum_{\hat c\in\widehat{\Ccal}}
\kappa_{t,\theta_t}^{\mathrm{impl}}(\hat c,a\mid h_t)
\]
of Proposition~\ref{prop:amortized-surrogate-minimizer}, and suppose this deployed process induces a semantic-learning package $\mathfrak K=(\mathcal I,\mathcal J,\Scal)$ satisfying Assumptions~\ref{ass:main-verified-package}. Then
\[
q_t^\star(\Scal(p^\star)\mid H_t)\longrightarrow 1
\qquad
\text{almost surely.}
\]
The concrete surrogate assumptions \textup{(A1)--(A3)} remain the verified deployment-side route for producing a separating-support floor; the asymptotic recovery claim is obtained once that construction is certified as a verified package.
\end{theorem}

\begin{proof}
\noindent\emph{Lean: \path{SemanticConvergence.Ontology.Coupling.h10_amortized_surrogate}.}
The cited theorem is the amortized-surrogate label for the main theorem. The deployment-side floor calculation remains separately formalized as \path{SemanticConvergence.thm_amortized_surrogate}; it is one way to help verify the package assumptions, not an additional convergence axiom.
\end{proof}


\begin{corollary}[Finite-time certificate for the amortized surrogate]
\label{cor:amortized-surrogate-finite-time}
Suppose the concrete surrogate assumptions \textup{(A1)--(A3)} hold for the deployed policy, and suppose that for a finite horizon $T$ the same policy admits a zero-out rate package $\mathfrak Z_T(U,\pi^{\mathrm{impl}},p_{\mathrm{env}},p^\star,\omega_{\mathrm{behav}})$. Then the surrogate floor is verified and, for every decay parameter $\delta$,
\[
R_N(p^\star;\xi)\le \rho_\delta^N R_0(p^\star)
\qquad\text{and}\qquad
\left(1+\rho_\delta^N R_0(p^\star)\right)^{-1}\le Q_N(p^\star;\xi)
\]
for every $N\le T$, almost surely under the constructed countable trajectory law.
\end{corollary}

\begin{proof}
\noindent\emph{Lean: \path{SemanticConvergence.cor_amortized_surrogate_finite_time_zeroOutPackage}.}
The cited Lean theorem combines the deployment-side surrogate support-floor theorem with the zero-out package's residual-rate and posterior-share wrappers.
\end{proof}

\begin{remark}
\label{rem:amortized-surrogate-scope}
The guarantee is intentionally conditional. It does \emph{not} say that stochastic gradient descent finds parameters satisfying \textup{(A1)--(A3)}, nor that the latent class $\hat c^\star$ is automatically the true semantic class. It says something narrower and theorem-facing: if the learned model keeps a designated target latent class alive, exposes a uniform positive-$\lambda$ region of high learned semantic score, and that high score is sound for genuine class-against-complement separation, then minimizing the surrogate action loss yields the separating-support condition needed for semantic recovery; finite-time residual/posterior-share bounds require the additional zero-out certificate stated in Corollary~\ref{cor:amortized-surrogate-finite-time}.
\end{remark}


\begin{example}[Transformer-style language model]
\label{ex:transformer-surrogate}
Consider a learner whose latent class set $\widehat{\Ccal}$ enumerates candidate world models or task hypotheses, $\rho_t(\hat c\mid h_t)$ is the model's posterior over those candidates given the prefix $h_t$, and $\widehat\mu_{\theta_t}(o\mid h_t,a,\hat c)$ is its predictive distribution over next observations conditioned on a candidate. A standard predictive loss $\mathcal L_t^{\mathrm{pred}}(\theta_t;h_t)$ trains the predictive distribution. The amortized surrogate $\mathcal L_{t,\mathrm{impl}}$ adds two terms: the class-against-complement Bhattacharyya score $\bar{\widehat B}_{t,\theta_t}(\hat c,a;h_t)$, weighted by an action-class kernel $\kappa$, and a KL regularizer pinning $\kappa$ to a reference action distribution.

In this setting the deployment-side hypotheses say something concrete. Assumption~(A1) says the true world is in $\widehat{\Ccal}$ --- the candidate set is rich enough to contain $p^\star$. Assumption~(A2) says the policy keeps a positive-mass region of high learned Bhattacharyya score uniformly accessible across histories --- the model continues to ask separating questions, not just questions whose answers it already partially knows. Assumption~(A3) says the learned score is sound: high $\bar{\widehat B}_{t,\theta_t}$ implies actual separation under the true law $\mu_{p^\star}$, not merely under the model's own beliefs about itself.

A learner that minimizes $\mathcal L_{t,\mathrm{impl}}$ while satisfying (A1)--(A3) is one to which Theorem~\ref{thm:separating-support-convergence} applies as a guarantee. A learner that minimizes only $\mathcal L_t^{\mathrm{pred}}$ does not in general satisfy any of the three hypotheses, and Theorem~\ref{thm:afe-near-miss} is the corresponding negative finding for the natural information-gain completion of that learner. The surrogate is the smallest theorem-facing modification of standard predictive training that supplies a class-against-complement criterion; the diagnostic of Section~\ref{sec:diagnostic} is the corresponding audit procedure.
\end{example}

\subsection*{Open question}

What tractable approximations preserve the theorem-level target? The differentiable surrogate above is theorem-facing: a finite latent world model plus the loss $\mathcal L_{t,\mathrm{impl}}$ inherits the semantic-recovery and finite-time guarantees once assumptions \textup{(A1)--(A3)} certify that the learned high-score region is nonvanishing and semantically sound. The harder question is the approximation question: for which model classes, training procedures, and amortized inference schemes can one prove those assumptions from the optimization itself rather than impose them as deployment-side hypotheses?

\section{A Diagnostic for Interactive Learning Failures}
\label{app:diagnostic}
\label{sec:diagnostic}

This section turns the three-layer ontology of Section~\ref{sec:ontology} into a procedure for diagnosing any specific interactive learner that has failed to concentrate its posterior on the desired semantic target. The intent is for a reader holding a specific learner --- a language model evaluated on a specific task, an active-inference controller, a Solomonoff-style universal predictor, or any other interactive learner that has failed to concentrate its posterior on the desired semantic target --- to localize what went wrong and which remedy applies, by reading off the answer to three questions in a fixed order.

\subsection{The three diagnostic questions}
\label{app:diagnostic-questions}

\noindent\fbox{\parbox{0.93\textwidth}{\centering
\textbf{The diagnostic, in three questions.}\\[0.3em]
\begin{enumerate}
\item Is the observer fine enough to resolve the target?
\item Is the agent in the natural variational class?
\item Is the agent $\kappa$-coupled to the environment through the observer?
\end{enumerate}
A learner that passes all three is one to which Theorem~\ref{thm:separating-support-convergence} applies as a guarantee.}}

\medskip

\paragraph{Question 1 (observer): is the access pattern fine enough?}
Identify the observer $\omega$ underlying the learner's evidence model. If $\omega$ is strictly coarser than the target observer, no learning rule can resolve the desired semantic target through it; the observable-quotient theorem (Theorem~\ref{thm:factor-through-quotient}) is the exact quotient statement used here. If a strict-hierarchy witness package is supplied for a finer target, Theorem~\ref{thm:strict-hierarchy} records the corresponding strict gap. The diagnostic outcome at this layer is whether the observer and target fiber match the semantic set used by the verified package.

\emph{Remedy.} Change the access pattern, refine or coarsen the observer, or accept that the chosen target is unreachable through this access and substitute a different target.

\paragraph{Question 2 (agent): is the learning architecture in the natural variational class?}
Holding the observer fixed, ask whether the learner's belief update is represented by the Bayes/Gibbs variational identity and minimizer propositions, uses KL within the natural convex class (Lemma~\ref{lem:kl-necessity}), and has an action objective represented by the formal risk-minus-gain or kernel-functional interfaces. An architecture outside this formal interface is not covered by the main theorem of Section~\ref{sec:main}.

\emph{Remedy.} Replace the offending component with one that satisfies the formal variational and action-kernel predicates used by the verified package.

\paragraph{Question 3 (coupling): is the agent $\kappa$-coupled to the environment through the observer?}
Once the observer is at $\omega_{\mathrm{behav}}$ and the agent is in the natural variational class, convergence is a question of coupling: does the agent's policy produce class-predictive Bhattacharyya separation $\kappa>0$ on its policy support (Assumption~\textup{(M7)}), and does the (agent, environment) pair through this observer satisfy the prefix-Hellinger obligations (Assumption~\textup{(M6)})? The boundary near-miss theorem (Theorem~\ref{thm:afe-near-miss}) is a witness package for one kind of coupling failure. The matching converses (Corollary~\ref{cor:support-necessary}, Theorem~\ref{thm:summable-support-insufficient}) are coupling failures of a different form --- the support-mass form rather than the score form.

\emph{Remedy.} Engineer the policy or the reference law. The realization wrappers of Section~\ref{sec:semantic-action} and the differentiable surrogate of Section~\ref{sec:amortized-surrogate} are recipes for satisfying the coupling conditions of Assumptions~\ref{ass:main-verified-package}.

\subsection{Diagnostic cross-reference table}
\label{app:diagnostic-table}

Table~\ref{tab:diagnostic} maps an observed symptom to a diagnosis, a canonical theorem in this paper that exhibits or characterizes the failure mode, and a remedy class.

\begin{table}[h]
\centering\small
\renewcommand{\arraystretch}{1.25}
\begin{tabular}{p{0.30\linewidth}p{0.16\linewidth}p{0.20\linewidth}p{0.26\linewidth}}
\toprule
Symptom & Layer & Canonical theorem & Remedy class \\
\midrule
Target class cannot in principle be resolved through the learner's access. & Observer & Theorem~\ref{thm:factor-through-quotient} (observable quotient) & Refine the access pattern; change the observer. \\
Target class is finer than the access pattern can resolve along any policy. & Observer & Theorem~\ref{thm:strict-hierarchy} (strict hierarchy) & Coarsen the target to a class at $\omega_{\mathrm{behav}}$. \\
Belief update is not represented by the Bayes/Gibbs identity and minimizer package. & Agent & Lemma~\ref{lem:variational} and Propositions~\ref{prop:two-observer-minimizer}--\ref{prop:kernel-functional-minimizer} & Replace the update with the Bayes/Gibbs minimizer. \\
Belief update minimizes a non-KL divergence on the natural convex class. & Agent & Lemma~\ref{lem:kl-necessity} (KL necessity) & Replace the divergence with KL against the universal prior. \\
Action rule not represented by the boundary risk-minus-gain proxy. & Agent & Corollary~\ref{cor:efe-specialization} (EFE proxy) & Re-express the action objective in the formal risk-minus-gain interface. \\
Supplied boundary action and frozen-posterior predicates certify non-improvement. & Coupling & Theorem~\ref{thm:afe-near-miss} (boundary witness package) & Engineer reference law or policy for class-predictive separation. \\
Zero separating policy mass on the agent's support. & Coupling & Corollary~\ref{cor:support-necessary} (support-necessary) & Add exploration to the policy until separating mass becomes positive. \\
Separating support is positive but summable; still no convergence. & Coupling & Theorem~\ref{thm:summable-support-insufficient} (summable-support failure) & Increase exploration rate so separating mass diverges. \\
Verified main-theorem hypotheses fail at some prefix. & Coupling & Theorem~\ref{thm:separating-support-convergence} (verified main) & Adjust policy, kernel, or reference law to restore Assumption~\textup{(M6)} or~\textup{(M7)}. \\
\bottomrule
\end{tabular}
\caption{Symptom-to-remedy lookup. Each row corresponds to a single failure mode and points to the named theorem or witness package in this paper that records it. The remedy class names the kind of intervention that would restore the corresponding hypothesis or close the layer-specific gap.}
\label{tab:diagnostic}
\end{table}

\subsection{Worked example: applying the diagnostic to a boundary near-miss witness}
\label{app:diagnostic-example}

Consider a learner using the algorithmic free energy principle as its complete architecture: universal-prior belief, Bayes/Gibbs update, and a boundary expected-free-energy proxy. Suppose the boundary action witness, frozen-posterior-through-horizon predicate, and observer-promotion failure predicate required by Theorem~\ref{thm:afe-near-miss} are supplied. Apply the diagnostic.

\emph{Question 1 (observer).} The supplied witness states which observer interface is being used and whether an observer-promotion failure predicate holds. If that predicate says the realized policy observer is too coarse for the target, the failure is visible at the observer/coupling boundary.

\emph{Question 2 (agent).} The belief update is structurally adequate when it is represented by the Bayes/Gibbs variational identity and minimizer propositions. The boundary proxy is only a formal risk-minus-gain interface; if a full classical EFE rule is desired, it must be connected to this interface separately.

\emph{Question 3 (coupling).} The supplied frozen-posterior and boundary-action witnesses certify the coupling failure: the prefix-Hellinger and class-predictive separation assumptions are not available from those witnesses. The remedy is to engineer the prior, policy, or reference law so that Assumptions~\textup{(M6)} and~\textup{(M7)} can be verified; the verified realizations and the surrogate of Sections~\ref{sec:semantic-action} and~\ref{sec:amortized-surrogate} are constructions aimed at that task.

The diagnostic localizes the formal witness package to the coupling layer. The positive theorem is recovered only by replacing the failed witness package with the bundled assumptions.

\subsection{Notes on application}
\label{app:diagnostic-notes}

Three cautions for applying the diagnostic to non-formal learners.

First, identifying the observer requires a notion of distinguishability appropriate to the application. For a language model, the natural observer might track the model's response distribution on a specified prompt distribution; for an active-inference controller, the observer is the generative model's latent variable structure; for a real-world Solomonoff predictor, the observer is the access through the chosen interaction protocol. The diagnostic's first question is sharp; the work of identifying the right $\omega$ for the application is the application designer's, not the diagnostic's.

Second, Question~2 is a more substantive structural test than it may appear. Many practical architectures are not natural-class agents in this paper's sense even when they look variational --- ELBO-based variational autoencoders against parametric families, for instance, do not in general satisfy KL-necessity over the natural convex class, because the family restricts the divergence in question. The diagnostic's second question detects this; an architecture that fails it is at best an approximation to a natural-class agent, and the convergence theorem applies only as far as the approximation does.

Third, coupling conditions are not architectural and therefore not invariant under environment change. An agent that satisfies all coupling conditions on environment $\varepsilon_1$ may fail them on environment $\varepsilon_2$. The diagnostic's third question must be re-evaluated for each new environment of interest. Failures localized to the coupling layer always have engineering remedies of the kind cataloged in Sections~\ref{sec:semantic-action} and~\ref{sec:amortized-surrogate}, but those remedies are environment-specific.

The three layers and three questions together constitute the operational form of this paper's contribution. A learner that passes all three questions is one to which the verified main theorem (Theorem~\ref{thm:separating-support-convergence}) applies as a guarantee.


\section{Discussion}
\label{sec:discussion}

The boxed problem has two logically distinct parts. The belief rule is characterized by the Bayes/Gibbs variational identity (Lemma~\ref{lem:variational}) and the exact minimizer propositions for the repaired Gibbs objectives, while Lemma~\ref{lem:kl-necessity} gives the countable convex-duality uniqueness statement for KL under its formal predicates. The action/convergence rule is verified in Theorem~\ref{thm:separating-support-convergence}: a matched semantic-learning package $\mathfrak K=(\mathcal I,\mathcal J,\Scal)$, interpreted through the canonical predictive-kernel observer and the Assumptions~\ref{ass:main-verified-package}, yields almost-sure posterior concentration on $\Scal(p^\star)$. The support/rate route remains a separate zero-out companion, and the boundary near-miss theorem is a witness package rather than a constructed counterexample.

The observer frame can be read alongside three traditions. Pearl's do-calculus identifies when interventional distributions are determined by observational ones via graphical structure \citep{Pearl2009}; Theorem~\ref{thm:factor-through-quotient} supplies a different, non-graphical formal analogue inside the interactive program-observer setting, and Corollary~\ref{cor:promotion-contrast} packages the corresponding observer-promotion failure predicate. The observable-quotient theorem is not a do-calculus result and does not subsume one. \citet{Chernoff1959}'s sequential-testing rule motivates class-against-complement scoring, while the Lean theorem itself uses the bundled class-predictive separation field. \citet{Blackwell1953}'s comparison of experiments motivates the refinement preorder, while the verified convergence statement is the verified package theorem through the canonical predictive-kernel observer.

Hume's problem of induction \citep{Hume1748} --- whether generalizing from observed instances to unobserved ones admits a non-circular justification --- sits underneath the philosophical territory the rest of this section engages. The paper does not offer a non-circular justification of induction, and is not intended to. Inside the program/observer setting, what it offers is a precise account of where inductive reasoning reaches and where it does not. The observable-quotient theorem (Theorem~\ref{thm:factor-through-quotient}) identifies the interventional class $[p^\star]$ as the strongest target a history-based learner can in principle reach; the strict-hierarchy theorem (Theorem~\ref{thm:strict-hierarchy}) locates the residual underdetermination at the gap $R_\pi(p^\star)\setminus[p^\star]$; and the main theorem (Theorem~\ref{thm:separating-support-convergence}) identifies bundled coupling conditions sufficient for convergence at that ceiling. The coupling conditions --- class-predictive Bhattacharyya separation, the prefix-Hellinger obligations --- are not derivable from history alone; they are the kind of extra-empirical structure Hume's argument predicts any successful inductive procedure must rest on. The framework can therefore be read as making explicit what Hume's argument leaves implicit: \emph{which} extra-empirical assumptions buy convergence, and at \emph{which} ceiling. The Solomonoff/Hutter universal prior is one such assumption, and it is necessary --- but Theorem~\ref{thm:afe-near-miss} shows it is not sufficient: the universal prior plus the natural information-gain action rule does not in general converge to $[p^\star]$. The kernel-functional refinement supplies what is needed beyond it, and its sufficient conditions are themselves bundled assumptions about the environment-agent pair rather than consequences of history. The framework neither solves Hume's problem nor dismisses it; it specifies its formal shape inside the present setting.

Three later theses descended from Hume's underdetermination concern have formal analogues in this paper, again inside the program/observer setting. \citet{Quine1960}'s claim that empirical evidence leaves logically incompatible theories equally supported has a formal analogue in the strict-hierarchy witness package: when the policy-gap witness is supplied, there are programs in $R_\pi(p^\star)\setminus[p^\star]$. \citet{vanFraassen1980}'s constructive empiricism --- the view that science aims at empirical adequacy rather than truth --- has a formal analogue in Theorem~\ref{thm:factor-through-quotient}, which identifies the policy-view-recoverable semantic quotient as the strongest target a history-based learner can reach. \citet{Putnam1980} questioned whether use could fix reference at all; Theorem~\ref{thm:semantic-convergence} can be read as showing that, within this paper's framework, use fixes the semantic target once the bundled coupling assumptions are verified. None of these formal analogues is offered as a global resolution of the original philosophical thesis; each is a specific theorem inside the program-observer setting we have set up, under bundled assumptions stated in the paper.

The Principle of Section~\ref{sec:near-miss} positions this paper relative to the active-inference and Solomonoff/Hutter traditions, in a specific and limited way. The Gibbs variational form of active inference is borrowed on the belief-evidence component; the natural information-gain completion of that form, however, does not in general converge to the interventional class (Theorem~\ref{thm:afe-near-miss}), and the kernel-functional refinement of Section~\ref{sec:main} replaces information gain with class-against-complement separation on the action side. Solomonoff and Hutter's universal prior supplies the reference measure used here; AIXI-style universal-prediction guarantees, by contrast, are stated at the policy-view layer $R_\pi(p^\star)$ rather than at the interventional layer $[p^\star]$, and Theorem~\ref{thm:separating-support-convergence} can be read as an interventional-layer counterpart inside the present framework. In a slogan: \emph{curiosity-under-the-universal-prior is correct in form, incorrect in rule.} This is meant as a positioning of the present contribution within these traditions, not as a global judgement on either.

One open theoretical question deserves to be highlighted alongside the open problems noted at the close of Sections~\ref{sec:near-miss} and~\ref{sec:realization}.

What quantitative rates can be proved? Theorem~\ref{thm:separating-support-rate} and Corollary~\ref{cor:separating-support-finite-time} give explicit finite-horizon residual/posterior-share bounds once the zero-out package is verified, and Propositions~\ref{prop:kernel-exp-rate} and~\ref{prop:exp-rate}, together with Theorem~\ref{thm:exp-rate-concentration}, record the same-view and package-level specializations. Corollary~\ref{cor:compact-action-kernel} shows how the compact-action support geometry feeds such a package once the required realized-update witness is supplied. What remains open is deriving comparable zero-out certificates, or sharper non-zero-out finite-time rates, directly from broad semantic separation hypotheses. Sequential design results in statistics suggest such bounds should be possible in restricted regimes \citep{Chernoff1959,BoxHill1967}.

\section{Conclusion}
\label{sec:conclusion}

The paper develops three layers in turn. At the \emph{observer} layer, four nested sets $\{p^\star\}\subseteq[p^\star]\subseteq R_\pi(p^\star)\subseteq R_{h_t}(p^\star)$ are the indistinguishability classes of four observers at increasing observational power, with strictness recorded through explicit witness packages; the observable-quotient theorem identifies $[p^\star]$ as the ceiling for exactly policy-view recoverable targets. At the \emph{agent} layer, the Bayes/Gibbs posterior is the belief-side minimizer of the repaired Gibbs objectives; the KL divergence is forced within the natural convex class by convex duality; and the boundary expected-free-energy proxy records the risk-minus-gain shape used to compare action rules. At the \emph{coupling} layer, Theorem~\ref{thm:separating-support-convergence} is the verified positive answer: when an admissible agent is $\kappa$-coupled to the environment through the canonical predictive-kernel observer, the Bayes/Gibbs posterior converges almost surely to the packaged predictive semantic target. A zero-out special-case companion provides finite-horizon rate bounds, and the boundary near-miss predicates record the formal shape of coupling failures once their witnesses are supplied. The framework's reach beyond the convergence theorem is developed in Part~\ref{part:reach}: Section~\ref{sec:near-miss} shows the natural rule fails, Section~\ref{sec:self-ref} dissolves uncomputability, Section~\ref{sec:realization} realizes the framework as a differentiable loss, and Section~\ref{sec:diagnostic} turns it into an audit procedure.

The paper's epistemological move is to identify the right object of interactive learning as concentration on the interventional semantic class --- the strongest target a history-based learner can reach, formalized through the observable-quotient theorem --- and to identify the right kind of action rule as one that scores observations by their power to distinguish the target's class from its complement under the true environment law. Generic information gain, even under the universal prior, does not suffice (Theorem~\ref{thm:afe-near-miss}); the class-against-complement Bhattacharyya criterion does (Theorem~\ref{thm:separating-support-convergence}). \emph{Curiosity is necessary but not sufficient. Explanation, formalized as class-against-complement coupling, is what suffices for interactive understanding.}

The full Lean correspondence is available at the public repository \url{https://github.com/Populustremuloides/semantic_convergence}. Every formal proof in this paper carries an exact Lean citation in its proof body; generated coverage, proof-shape, and per-declaration axiom audits document the verification state.

\bibliographystyle{plainnat}
\bibliography{semantic_convergence_interactive_learning}

\end{document}
